% !TEX program = xelatex
% 编译方式：xelatex opf_manual.tex （需两遍以生成目录）
% 注意：MiKTeX 在含中文的路径下编译会崩溃，请复制到纯英文路径（如 C:\latex_build\opf_manual\）再编译。
\documentclass[UTF8,fontset=windows,a4paper,11pt]{ctexart}

\usepackage[margin=2.2cm]{geometry}
\usepackage{amsmath,amssymb}
\usepackage{booktabs,longtable,array}
\usepackage{xcolor}
\usepackage{float}
\usepackage{hyperref}
\hypersetup{colorlinks=true,linkcolor=black,urlcolor=blue}

% ── 统一参数表环境 ────────────────────────────────────────────────
% 六列：字段名 | 符号 | 单位 | 典型范围 | 默认值 | 说明
\emergencystretch=3em
\newcolumntype{L}[1]{>{\raggedright\arraybackslash}p{#1}}
\newenvironment{paramtable}{%
  \begin{footnotesize}
  \begin{longtable}{@{}L{3.2cm}L{1.3cm}L{1.6cm}L{2.3cm}L{1.5cm}L{4.3cm}@{}}
  \toprule
  \textbf{字段名} & \textbf{符号} & \textbf{单位} & \textbf{典型范围} & \textbf{默认值} & \textbf{说明}\\
  \midrule
  \endfirsthead
  \multicolumn{6}{@{}l}{\footnotesize（续表）}\\
  \toprule
  \textbf{字段名} & \textbf{符号} & \textbf{单位} & \textbf{典型范围} & \textbf{默认值} & \textbf{说明}\\
  \midrule
  \endhead
  \midrule
  \multicolumn{6}{r@{}}{\footnotesize（接下页）}\\
  \endfoot
  \bottomrule
  \endlastfoot
}{%
  \end{longtable}
  \end{footnotesize}
}

% 字段名排版（下划线需转义后传入；\_ 处允许换行以防超长字段名溢出版心）
\makeatletter
\newcommand{\fld}[1]{\texttt{\fld@scan#1\fld@stop}}
\def\fld@scan#1{%
  \ifx#1\fld@stop
  \else
    \ifx#1\_\_\allowbreak\else#1\fi
    \expandafter\fld@scan
  \fi}
\makeatother
% 模块元信息行：#1=求解器/类名  #2=头文件路径
\newcommand{\compmeta}[2]{\noindent\textbf{代码结构体：}\texttt{#1}\hfill
  \textbf{定义位置：}\texttt{#2}\par\smallskip}

% 可断行的源码路径/函数名（/ 与 \_ 处允许换行）
\makeatletter
\newcommand{\srcpath}[1]{\texttt{\src@scan#1\src@stop}}
\def\src@scan#1{%
  \ifx#1\src@stop
  \else
    \ifx#1\_\_\allowbreak\else
    \ifx#1/#1\allowbreak\else
    #1\fi\fi
    \expandafter\src@scan
  \fi}
\makeatother

\title{\textbf{交直流混合高弹性能源电力系统仿真平台}\\[0.5em]
       {\Large 最优潮流技术手册}\\[0.3em]
       {\large 数学模型、算法流程、参数选项与验证校核}}
\author{HySim-XJTU-HRPES（西安交通大学 HRPES 团队）}
\date{\today}

\begin{document}

\maketitle
\begin{abstract}
\noindent 本手册依据 HySim-XJTU-HRPES 平台 C++ 源代码中的最优潮流（OPF）模块
（\texttt{src/\allowbreak optimal\_power\_flow/}、\texttt{include/\allowbreak hacdcpf/\allowbreak optimal\_power\_flow/}
与 \texttt{src/\allowbreak power\_models/}、\texttt{include/\allowbreak hacdcpf/\allowbreak power\_models/} 目录）
整理，覆盖 AC OPF 顶层求解器（后端分发、经济调度兜底、目标同伦、Native AC 简约空间路径）、
Parity 全空间 NLP 建模（换流器、DC/DC、储能、可切负荷、可再生、柔性负荷、能量路由器全家族）、
原生原对偶内点法（Mehrotra 预估校正、Gondzio 校正、滤波线搜索、惯性控制）、
DC OPF（QP/LP 链与 LMP）、无功优化（OLTC 离散档位与可投切并联的 MINLP 搜索）、
三相混合 OPF（相域直角坐标、GFL/GFM 换流器、图降阶变体），
以及基于 MIPSolvers AML 的求解器无关建模层（AC/混合/DC OPF、LinDistFlow、SCUC builder）、
OPF 在平台中的下游应用（时序、可靠性、EV-交通、市场边界）与 GUI/HTTP 接入，
并给出全部选项字段的默认值、结果与诊断数据结构、以及测试验证体系。
文档风格与本平台《潮流计算技术手册》一致，全部公式逐条注明源码出处（文件:行号），
近似、未接线、覆盖薄弱处均如实标注。
\end{abstract}

\tableofcontents
\newpage

\section{阅读指南与约定}

\subsection{数据来源}
本手册全部内容取自以下源代码与文档，公式与参数均以代码为准：
\begin{itemize}
  \item \texttt{src/optimal\_power\_flow/} —— 6 个实现文件：\texttt{ac\_opf.cpp}（顶层路由、
        Native AC 简约空间 IPM、Ipopt 适配、经济调度兜底、目标同伦）、
        \texttt{parity\_formulation.cpp}（全空间 NLP 装配与解析导数）、
        \texttt{parity\_ipm.cpp}（原生原对偶 IPM）、\texttt{dc\_opf.cpp}（QP/LP 链）、
        \texttt{reactive\_power\_opt.cpp}（RPO）、\texttt{three\_phase\_hybrid\_opf.cpp}（三相混合 OPF）；
  \item \texttt{include/hacdcpf/optimal\_power\_flow/} —— 10 个头文件（求解器接口、选项、
        结果、formulation、IPM 规范头）；
  \item \texttt{src/power\_models/} 与 \texttt{include/hacdcpf/power\_models/} ——
        AML 求解器无关建模层（acopf/acdcopf/dc\_opf/lindistflow/scuc builder）；
  \item \texttt{MIPSolvers/include/mipsolvers/aml/} —— AML 建模原语（Model/变量/约束/表达式/
        求解结果），经 \texttt{include/hacdcpf/aml/aml.hpp} 转发；
  \item \texttt{src/api/hacdcpf.cpp} —— 公共 API 门面（求解分发、结果回投影、独立审计）；
  \item \srcpath{docs/technical\_notebook/sections/04\_opf.tex}、\srcpath{06\_power\_models.tex}、
        \srcpath{11c\_native\_ipm\_vs\_parity\_ipm.tex}、
        \srcpath{docs/reactive\_power\_optimization\_validation.md}、
        \srcpath{docs/large\_hybrid\_ipm\_diagnostics.md} —— 既有技术笔记，本手册与之交叉核对；
  \item \texttt{tests/} 与 \texttt{external\_data/} —— 验证与校核章节的测试矩阵与算例来源。
\end{itemize}
文中所有“\texttt{文件:行号}”引用基于当前代码检出处，代码演进后行号可能漂移，
但函数名与文件路径保持稳定。

\subsection{单位制约定}
平台统一采用 MW·kV·MVA 单位制配合标幺值：系统基准功率
$S_{\mathrm{base}}=100$~MVA（\texttt{Defaults::kBaseMva}），基准频率 50~Hz。
OPF 方程内部一律以标幺值求解；成本系数以 \$/(${MW}^2$h)、\$/(MWh)、\$/h 录入，
在 pu 域换算为 $a_2=c_2S_b^2$、$a_1=c_1S_b$、$a_0=c_0$。对外结果中功率以 MW/Mvar、
电压幅值以 pu、LMP 以 \$/MWh 给出；相角在结果结构体中以\textbf{弧度}存储
（\fld{ACOPFResult::va}/\fld{DCOPFResult::va}），展示层才换算为度（$^\circ$），
详见第~\ref{sec:options}~章总约定。
ZIP 负荷权重三元组顺序为（恒功率、恒电流、恒阻抗）。

\subsection{编号与符号约定}
\begin{itemize}
  \item \textbf{失配约定}：与潮流一致，$\mathrm{mismatch}=F^{\mathrm{spec}}-F^{\mathrm{calc}}$；
        IPM 等式残差 $g(x)=0$、不等式 $h(x)\le 0$。
  \item \textbf{注入约定}：节点注入以\textbf{发电为正、负荷为负}；换流器端口注入为正时
        满足 $P_{ac}+P_{dc}+P_{loss}=0$。
  \item \textbf{三类 ID 纪律}：组件 \texttt{.index}（对外稳定）、向量下标（内部 0-based）、
        图节点索引（临时）不混用；结果回投影一律经映射回到稳定 ID。
  \item \textbf{AC/DC 域隔离}：AC 与 DC 母线编号空间独立，跨域传递一律走
        domain-qualified 映射（\fld{ac\_bus\_id\_to\_node\_idx} /
        \fld{dc\_bus\_id\_to\_node\_idx}）。
\end{itemize}

\subsection{求解器家族总览}
\begin{center}
\footnotesize
\begin{tabular}{@{}llll@{}}
\toprule
求解器 & 入口函数 & 适用网络 & 定位\\
\midrule
AC OPF（顶层） & \srcpath{opf::solve\_ac\_opf} & 混合 AC/DC & \textbf{生产主路径}（后端分发）\\
Parity 全空间 NLP & \srcpath{parity::build\_problem} & 混合 AC/DC & 唯一同时优化 AC/DC/VSC/DCDC 的路径\\
原生原对偶 IPM & \srcpath{parity::solve\_primal\_dual\_ipm} & 混合 AC/DC & Parity 默认内层求解器\\
嵌入式 Ipopt & \srcpath{solve\_parity\_with\_ipopt} & 混合 AC/DC & 兜底/对照后端\\
Native AC 简约空间 & （\srcpath{OPFSolverPath::NativeAC}） & AC+部分 DC & 旧路径，罚函数处理限值\\
经济调度+PF & \srcpath{solve\_economic\_dispatch} & 纯 AC & 快速兜底（merit-order $\lambda$ 二分）\\
DC OPF & \srcpath{opf::solve\_dc\_opf} & 纯 AC 线性化 & QP/LP 链 + LMP\\
RPO 无功优化 & \srcpath{opf::solve\_rpo} & 混合 AC/DC & 离散档位/并联搜索 + 内层 AC OPF\\
三相混合 OPF & \srcpath{solve\_three\_phase\_hybrid\_opf} & 三相 AC+DC & 活跃研发，Ipopt/NativeIPM 双后端\\
AML builders & \srcpath{power\_models::solve\_*} & 见第~\ref{sec:powermodels}~章 & 求解器无关建模层（非生产链路）\\
\bottomrule
\end{tabular}
\end{center}

\subsection{验证通道总览}
各章末的“验证与校核”小节给出具体对照依据，总体通道如下（详见第~\ref{sec:validation}~章）：
\begin{itemize}
  \item \textbf{手算小算例}：\srcpath{external\_data/opf\_hand\_cases/} 4 个带解析解的
        DC OPF 案例（merit-order、双母线阻塞、切负荷、二次成本线性化），供 GUI 手工核对；
  \item \textbf{MATPOWER 算例回归}：case4gs 至 case13659pegase 共 85 个 .m 算例
        （79 标准 + 4 个 \texttt{contab\_*} + 2 个 \texttt{scenarios\_*}，分档见第~\ref{sec:validation}~章），
        覆盖 parity IPM 收敛、后端语义、暖启动、目标缩放、大系统稀疏 KKT
        （\srcpath{tests/test\_opf\_solver\_backends.cpp}）；
  \item \textbf{三方交叉验证}：AC OPF / DC OPF / PF 调度回放一致性
        （\srcpath{tests/test\_acopf\_dcopf\_crossval.cpp}），含 LMP 符号与阻塞再调度；
  \item \textbf{RPO 交叉验证}：独立 OPF 复算 + PF 回放，判据与容差见
        \srcpath{docs/reactive\_power\_optimization\_validation.md}；
  \item \textbf{解析导数校核}：\srcpath{check\_ac\_opf\_jacobian\_fd}（有限差分对照，
        见第~\ref{sec:acopf}~章诚实口径注）、DC slack KCL Jacobian FD（$<10^{-6}$）、
        三相 OPF Jacobian/Hessian（$<10^{-5}$）；
  \item \textbf{独立审计}：\srcpath{verify\_opf\_result} 对 OPF 结果重算电压限、机组限与目标值；
  \item \textbf{GUI 端到端}：\srcpath{tools/gui\_api\_e2e.py}（\texttt{/api/session/opf} 回归）、
        \srcpath{tools/validate\_case\_capabilities.py}（能力表打点）。
\end{itemize}

\subsection{现状与诚实口径声明}
以下事项在各章节均有对应标注，此处集中声明（如实）：
\begin{enumerate}
  \item \textbf{DC OPF 支路对偶未认证}：原生单纯形不能正确处理有界 $P_f$ 变量，
        \fld{branch\_mu\_valid} \textbf{恒为 false}（dc\_opf.cpp:680, 783），
        \fld{branch\_mu\_lower/upper} 在支路限生效时不可作工程解释，直到实现 KKT 基提取；
        该事实现附机器可读原因串 \fld{branch\_mu\_validity\_reason}（dc\_opf.cpp:681--682）。
        注意区分：\fld{lmp\_valid}=true 仅认证\textbf{平衡行对偶}（节点电价能量分量），
        \textbf{不蕴含}支路拥塞对偶有效（opf\_result.hpp:205--206）。
  \item \textbf{RPO 无全局最优证书}：\fld{globally\_certified}=false、\fld{gap}=1.0，
        只保证可行局部最优；\fld{bc\_stats} 仅为兼容信封，非 B\&B 全局 gap 证书；
        时限/评估上限截断以 \fld{terminated\_by\_time\_limit}/
        \fld{terminated\_by\_evaluation\_limit} 标志与 \fld{model\_limitations} 条目对外暴露
        （reactive\_power\_opt.cpp:1296--1315）。
  \item \textbf{AML power\_models 层不是生产 OPF 链路}：五个 builder 在生产 \texttt{src/}
        中无调用方，仅被测试使用；其中 AC OPF / DC OPF / LinDistFlow / SCUC 的 AML 版本
        无直接单元测试；该声明已形式化为各 builder 结果的 \fld{model\_scope}/
        \fld{model\_limitations} 字段（如 \fld{AMLBuildResult} 默认
        "experimental-aml-builder:not-production-runtime"，aml\_build\_result.hpp:17--19）。
        \fld{power\_models::DCOPFResult} 与 \fld{opf::DCOPFResult} 同名不同物。
  \item \textbf{Native AC 简约空间旧路径}用\textbf{外部二次罚}（非障碍）处理支路热限与
        VSC 容量圆，换流器不钉 DC\_V 母线电压
        （\fld{vsc\_vdc\_control\_modelled}=false，ac\_opf.cpp:2723）；
        该路径结果附加 "legacy 外部障碍/罚循环，生产建议 ParityIPM" 的
        \fld{model\_limitations} 条目（ac\_opf.cpp:2831--2848）。
  \item \textbf{储能单时段无 SOC 模型}，parity 路径中钉在 UC 调度值 $\pm10^{-4}$/base
        （\srcpath{parity\_formulation.cpp:674--717}）；\textbf{能量路由器为无损守恒行} $\sum P_{port}=0$，
        \fld{loss\_percent} 未接入（:1563--1564）。
  \item \textbf{AML 混合 OPF 的 DC/DC $I^2R$ 损耗现已建模}：\fld{r\_eq\_pu} 纳入数据转换
        （acdcopf\_builder.cpp:179），约束为 $P_{in}=P_{out}/\eta + r_{eq}(P_{out}/\eta)^2/V_{dc,in}^2$
        （:586--591），平衡审计与约束采用同一效率+$I^2R$ 约定（:745--758），
        修复了审计把正常 $I^2R$ 损耗误报为平衡违规的问题；残余限制为仅正向传输的
        smooth 模型、不表示双向符号切换（hybrid\_opf\_model\_builder.hpp:83--86）。
  \item \textbf{Ipopt 适配器不返回约束乘子}，Ipopt 路径下 $\lambda$ 为空、LMP 提取跳过
        （ac\_opf.cpp:2105--2108）；该事实现已形式化为 \fld{lmp\_valid}=false 及
        \fld{lmp\_validity\_reason} 原因串（opf\_result.hpp:101--102），Parity IPM
        路径在等式乘子齐全时置 \fld{lmp\_valid}=true（ac\_opf.cpp:2619--2628）。
  \item \textbf{DC OPF 的 LP 回退采用 lambda 形式 PWL 凸组合}：仅对 $c_2>10^{-12}$ 且
        $P^{max}>P^{min}$ 的机组，按 \fld{pwl\_segments}（默认 4，clamp 到 $[1,1000]$，
        dc\_opf.cpp:166）取等距断点，$\lambda\in[0,1]$ 凸组合插值真实成本、无需 SOS2
        （:471--491）；无二次项机组保留纯 $c_1$ 线性（:189--203）。
        \fld{objective\_model} 如实区分 "QP"/"LP-PWL"/"LP"，生效段数记入
        \fld{pwl\_segments\_effective}（QP 后端恒 0，:1200--1207）。
  \item \textbf{LinDistFlow builder} 仅建固定负荷：头文件注释与实现一致，明确不建
        DG/可控源/切负荷变量（lindistflow\_builder.hpp:15--16, 30--31）；热限为
        $|P|,|Q|\le S_{max}$ 分离盒式近似；结果为电压平方 pu$^2$。
  \item \textbf{AML SCUC 无网络约束}（无潮流、无 N-1）；MILP 下价格对偶通常不可得，
        已形式化为 \fld{price\_valid}=false 与 \fld{price\_validity\_reason}
        （scuc\_builder.hpp:86--92）。
  \item \texttt{opf\_problem.hpp} 的 \texttt{OPFProblem} 已有消费方
        \srcpath{opf\_solver\_interface.hpp}（\fld{OPFSolveResult} variant 与
        \texttt{DefaultOPFSolver} 按域转发到 \srcpath{solve\_ac\_opf}/\srcpath{solve\_dc\_opf}，
        已纳入公共门面 \srcpath{api/hacdcpf.hpp:12}），被测试消费
        （test\_hacdcdss.cpp:420）但生产 \texttt{src/} 无调用方；
        RPO 的 \fld{branching}/\fld{node\_sel}/\fld{int\_tol} 为遗留字段。
  \item \textbf{雅可比诊断入口未接线}：\srcpath{ac\_opf\_solver.hpp:54} 声明的公开入口
        \fld{compute\_jacobian\_diagnostics} 在全库\textbf{无实现}（调用会链接错误）；
        实际的 FD 诊断 \srcpath{check\_ac\_opf\_jacobian\_fd}（ac\_opf.cpp:3826--4024，
        签名与声明不同）目前无任何调用方，属死代码。
  \item 环境变量开关众多（\texttt{HACDCPF\_OPF\_KKT\_FORM}/\texttt{\_LINEAR\_SOLVER}/
        \texttt{\_MU\_STRATEGY}/\texttt{\_FILTER\_ALL} 等，全集见第~\ref{sec:options}~章），
        inertia gate 默认关，注释明确其效果因算例而异（parity\_ipm.cpp:1244--1259）。
  \item 三相混合 OPF 属\textbf{活跃研发/论文验证性质}，接口与产物格式仍可能调整；
        AC 支路热限与 VSC 调制比尚未进入相域模型；该路径不导出约束乘子，
        GUI 端点固定回显 \fld{lmp\_valid}=false 及原因。
\end{enumerate}

\newpage
% 第 1 章：架构总览与数据流
\section{架构总览与数据流}\label{sec:overview}

\subsection{两条 OPF 对外面}
平台 OPF 代码分两层，定位不同（docs/technical\_notebook/sections/04\_opf.tex:4--22）：
\begin{itemize}
  \item \textbf{生产链路} \srcpath{src/optimal\_power\_flow/}：\srcpath{opf::solve\_ac\_opf}
        （ac\_opf.cpp:2699）是全功能入口——含负荷削减、DER、储能、DC/DC、柔性负荷、
        能量路由器、多后端 fallback 链与诚实诊断；\srcpath{opf::solve\_dc\_opf}
        （dc\_opf.cpp:731）与 \srcpath{opf::solve\_rpo}（reactive\_power\_opt.cpp:621）同层。
  \item \textbf{AML 建模层} \srcpath{src/power\_models/}：基于兄弟仓库 MIPSolvers 的 AML
        的求解器无关 builder（acopf/acdcopf/dc\_opf/lindistflow/scuc），
        \textbf{不含}负荷削减、DER、储能、柔性负荷、能量路由器与 fallback 链，
        生产 \texttt{src/} 中无调用方，定位为独立/备用的公式化与验证层
        （详见第~\ref{sec:powermodels}~章）。
\end{itemize}

\subsection{数据流流水线}
\begin{quote}\footnotesize
\texttt{Rich HybridPowerSystem} $\to$ 校验/预处理（参考母线资格、孤岛预检、external grid 提升）
$\to$ rich$\to$canonical 投影 $\to$ NLP/LP/QP 装配（parity formulation 或 DC OPF LP）
$\to$ 求解（原生 IPM / Ipopt / QP-LP 链）$\to$ 结果回投影（ authored 母线/机组序）
$\to$ 独立审计（\srcpath{verify\_opf\_result}）
\end{quote}
混合算例（含 VSC/DC/DC/储能等 DC 侧组件）在分发处自动强制启用 parity 全空间路径
（ac\_opf.cpp:2930--2934）；纯 AC 算例可走 Native AC 简约空间旧路径或经济调度快路径。

\subsection{后端分发路由}\label{sec:overview-routing}
\fld{ACOPFSolverBackend} 语义（ac\_opf.cpp:2905--2934）：
\begin{center}
\footnotesize
\begin{tabular}{@{}lL{4.2cm}L{8.3cm}@{}}
\toprule
后端枚举 & 内层求解器 & 说明\\
\midrule
\texttt{Auto} & ParityIPM 先行 $\to$ Ipopt 兜底 & \fld{allow\_fallback} 时生效；
  \fld{ac\_pf\_warm\_start} 时 Ipopt 先行 + 原生 IPM 精炼\\
\texttt{ParityIPM} & 原生原对偶 IPM & 强制 \fld{use\_parity\_ipm}\\
\texttt{Ipopt} & 嵌入式 Ipopt & 强制 \fld{use\_parity\_ipm}；受
  \texttt{HACDCPF\_DISABLE\_IPOPT\_OPF} 门控\\
\texttt{EconomicDispatch} & merit-order $\lambda$ 二分 + 单次 AC NR & 仅纯 AC；
  混合算例禁用（会丢掉 DC/VSC/DCDC 方程）\\
\bottomrule
\end{tabular}
\end{center}
DC OPF 的后端链（\fld{solver\_chain} 全程记录，dc\_opf.cpp:945--1097）：
Auto $\to$ NativeQP(LCQP) $\to$ Gurobi QP $\to$ HiGHS（QP 路由，失败回退 HiGHS-LP）
$\to$ NativeDualSimplex。

\subsection{公共 API 门面}
\compmeta{hacdcpf::（公共门面）}{\srcpath{include/hacdcpf/api/hacdcpf.hpp}}
\begin{itemize}
  \item \srcpath{solve\_ac\_opf}（src/api/hacdcpf.cpp:1182--1207）：调用
        \srcpath{opf::solve\_ac\_opf} 后，收敛且有开关/断路器时补算设备端子潮流
        （canonical 电压投影 $\to$ 支路潮流 $\to$ \srcpath{unproject\_branch\_flows}
        $\to$ \srcpath{compute\_device\_terminal\_flows}）。
  \item \srcpath{solve\_dc\_opf} / \srcpath{solve\_rpo}（:1209--1215）：直通。
  \item \srcpath{verify\_opf\_result}（:1348--1451）：独立审计电压限、机组限、目标值重算
        （含 external grid 成本），生成 \fld{infeasibility\_hints}，写入 \fld{result.audit}。
  \item \srcpath{get\_solver\_capabilities}（:1297--1344）：运行时后端能力查询
        （HiGHS / Gurobi / Ipopt / KLU / UMFPACK 按编译宏上报）。
\end{itemize}
求解器命名空间入口（\texttt{hacdcpf::opf::}）：
\srcpath{solve\_ac\_opf}（ac\_opf\_solver.hpp:42）、
\srcpath{compute\_jacobian\_diagnostics}（:54）、
\srcpath{solve\_dc\_opf} 与 \srcpath{check\_dc\_opf\_feasibility}（\srcpath{dc\_opf\_solver.hpp:17,22}）、
\srcpath{solve\_rpo} / \srcpath{inspect\_rpo\_controls} / \srcpath{apply\_rpo\_discrete\_solution}
（\srcpath{reactive\_power\_opt.hpp:236--302}）、
\srcpath{solve\_three\_phase\_hybrid\_opf}（\srcpath{three\_phase\_hybrid\_opf.hpp:186}）。

\subsection{缓存与暖启动}
API 层无结果缓存；暖启动链存在于：
\begin{itemize}
  \item \fld{ACOPFOptions::warm\_start} + 结果中的 \fld{ipm\_*\_state} 字段续传
        （x/λ/μ/z 全套，opf\_result.hpp）；
  \item \fld{ac\_pf\_warm\_start}：先跑 AC 潮流取电压初值（实测使 case13659pegase
        189 次迭代/$\sim$40~s 收敛）；
  \item RPO 内部 \fld{eval\_cache}：离散候选解缓存，就地修改避免深拷贝；
  \item 目标同伦的 Davidenko 预测器以前一步 KKT 切线预热下一步（见第~\ref{sec:acopf}~章）。
\end{itemize}

\subsection{GUI/HTTP 接入}
生产路由见 \srcpath{tests/run\_gui\_server.cpp}：
\texttt{POST} \srcpath{/api/session/opf}（统一端点，solver = ac / parity / dc / robust，约束族开关，
\fld{network\_model} 支持 three\_phase\_hybrid）、\srcpath{/api/session/opf\_ac}、
\srcpath{/api/session/opf\_parity}（固定 parity IPM 无 fallback，400 迭代）、
\srcpath{/api/session/opf\_dc}、\srcpath{/api/session/run\_rpo}，
以及无状态轻量端点 \srcpath{/api/opf/ac}、\srcpath{/api/opf/dc}。
收敛后自动跑 post-PF 与碳流分析；parity/native/dc 路径以 \fld{scope.model\_scope}
返回实际约束范围。

% 第 2 章：AC OPF 顶层求解器
\section{AC OPF 顶层求解器}\label{sec:acopf}

\subsection{功能定位与入口}
\compmeta{opf::solve\_ac\_opf}{\srcpath{include/hacdcpf/optimal\_power\_flow/ac\_opf\_solver.hpp}}
入口签名（ac\_opf\_solver.hpp:42；实现 src/optimal\_power\_flow/ac\_opf.cpp:2712）：
\begin{quote}\footnotesize
\texttt{ACOPFResult solve\_ac\_opf(const HybridPowerSystem\& sys, const ACOPFOptions\& opt = \{\});}
\end{quote}
适用纯 AC 与混合 AC/DC 网络（VSC、DC/DC、储能、能量路由器、可切负荷、可再生、
柔性负荷全家族）。

\textbf{诚实口径注}：头文件 ac\_opf\_solver.hpp:54 声明的公开诊断入口
\texttt{compute\_jacobian\_diagnostics} 在全库\textbf{无实现}；实际存在的 FD 诊断是
\srcpath{check\_ac\_opf\_jacobian\_fd(sys, max\_columns, fd\_eps, ac\_eval\_threads)}
（ac\_opf.cpp:3826--4024，判据 $\max|\mathrm{err}|\le5\times10^{-4}$ 或
相对误差 $\le5\times10^{-3}$），目前无调用方（死代码）。

\subsection{预处理}
\begin{itemize}
  \item \textbf{参考母线资格校验} \srcpath{validation::validate\_reference\_bus\_eligibility}
        （ac\_opf.cpp:2768--2787）：裸理想 DC\_V 边界（无可调度 DC 平衡设备）直接拒绝
        并在 \fld{infeasibility\_hints} 写明原因；
  \item \textbf{拓扑孤岛预检}（:2903--2932）：无含 slack 的岛屿判不可行并给 hints；
  \item \textbf{external grid 提升为合成发电机}（:2792--2822）：界限
        $\pm\max(1, s\_sc\_max\_mva, base\_mva)$，合成 index = \texttt{INT\_MIN}+i；
        求解后拆分回 \fld{external\_grid\_p\_mw/q\_mvar} 与 authored \fld{pg\_mw}（:2829--2885），
        映射错位直接 \texttt{throw}（:2858--2862）。
        该统一后处理 lambda 开头回填 \fld{model\_limitations}（opf\_result.hpp:159）：
        非凸 NLP 只认证\textbf{局部 KKT 点}而非全局最优（:2831--2832）；
        有储能时追加``储能按单时段注入、无 SOC 跨时段动态''（:2833--2836）；
        有能量路由器时追加``端口平衡无损、内部变换损耗未表示''（:2837--2840）；
        \fld{lmp\_valid}=false 且无原因串时补通用原因（:2841--2844）；
        NativeAC 路径追加``legacy 外部障碍/罚循环，生产建议 ParityIPM''（:2845--2848）；
  \item 混合算例禁用 AC-only 经济调度兜底并追加 hint（:2887--2901）。
\end{itemize}

\subsection{后端分发与目标同伦}
分发语义见第~\ref{sec:overview-routing}~节。\textbf{目标联动}：非经济目标
（VoltageDeviation / ActiveLoss / VoltageDeviationAndLoss）下 parity 问题
\textbf{自动禁用切负荷}——\fld{form\_opt.load\_shedding} = (objective == Economic)
（ac\_opf.cpp:2277），即不含 $\Delta P_d/\Delta Q_d$ 变量与 VOLL 项
（RPO 语义：不许靠丢负荷改善电压/网损）。

\textbf{目标同伦}
（\fld{objective\_homotopy}，ac\_opf.cpp:2940--3066）求解路径
$P(t):\min\ t\cdot f$，将 \fld{cost\_c0/c1/c2} 与
\fld{voltage\_deviation\_weight}/\fld{active\_loss\_weight} 乘以 $t$（:2955--2969），
$t:0\to1$ 最多 120 步：
\begin{itemize}
  \item 步长自适应：快步（$\le20$ 次迭代）$dt$ 翻倍至 0.5 上限；失败减半，
        $dt<0.005$ 停止（:2979--3027）；
  \item Davidenko 预测：$w_{pred}=w+h\,w'$，$h$ 受参数度量信赖半径 $\kappa=0.05$ 限制，
        $h=\min(dt,\ 0.05/\lVert dx/dt\rVert_\infty)$（:2993--2994），
        松弛/不等式对偶地板 $10^{-8}$（:3014--3018）；
  \item 诚实失败报告：$t_{ok}<1$ 时返回 ``minimal-violation region; full objective not
        solved''，\fld{converged}=false（:3057--3062）。
\end{itemize}

\subsection{经济调度 + AC PF 快路径}
\srcpath{try\_dispatch\_pf\_fallback}（ac\_opf.cpp:1759）：merit-order $\lambda$ 二分经济调度
\srcpath{solve\_economic\_dispatch}（:1496--1610）+ 单次 AC NR 潮流（:1763--1840）。
二次成本 $C_g(P)=c_2P^2+c_1P+c_0$ 下增量成本 $\lambda=2c_2P+c_1$，
$\lambda$ 二分区间内按 $P_g=(\lambda-c_1)/(2c_2)$ 裁剪到 $[P^{min},P^{max}]$。
仅纯 AC；目标值由调度成本重算（:1827--1832）。

\subsection{嵌入式 Ipopt 适配器}
\srcpath{solve\_parity\_with\_ipopt}（ac\_opf.cpp:2105--2247）把 parity 问题经
engine NLPModel 接口交给 Ipopt：
\begin{itemize}
  \item NLPModel 映射：box 界直通变量界，无限界改 $\pm10^{20}$；
        \fld{solver\_options.tolerance}$=\min(\text{三个容差})$、
        \fld{acceptable\_tolerance}$=\max(\text{三个容差})$（:2164--2170）；
  \item \textbf{接受判据}（Ipopt 报 MaxIter 等非 success 时）：
        $\fld{primal\_feas}\le\fld{tol\_primal}$ 且
        $\fld{dual\_feas}\le1000\cdot\fld{tol\_dual}$ 且
        互补 $\le\max(10^{-2},\ 1000\cdot\fld{tol\_comp})$；
        接受时 status 记 ``Ipopt: acceptable residuals after ...''（:2214--2238）；
  \item 暖启动裁剪到界内 $\pm10^{-8}$，置 \fld{warm\_start\_used}（:2127--2138）；
        \fld{initial\_primal\_inf} 在 $x_0$ 处评估（:2140--2149）；
  \item \textbf{诚实口径}：Ipopt 适配器不返回约束乘子，$\lambda$ 为空（:2107--2108）。
        该事实现已\textbf{形式化}为 \fld{lmp\_valid}/\fld{lmp\_validity\_reason} 字段
        （opf\_result.hpp:101--102）：Parity IPM 路径在等式乘子维度齐全时置
        \fld{lmp\_valid}=true（:2619--2621）；否则 false，
        \fld{lmp\_validity\_reason} 区分``嵌入式 Ipopt 适配器不返回乘子''与
        ``未返回完整等式对偶''两种情形（:2622--2628）。
        受 \texttt{HACDCPF\_DISABLE\_IPOPT\_OPF} 门控（:2119--2124）；
        MUMPS 5.6.2 ABI 兼容注释见 :2114--2118。
\end{itemize}

\subsection{暖启动与状态续传}
\paragraph{\fld{ac\_pf\_warm\_start}（混合潮流暖启动）}
实际跑的是 \srcpath{powerflow::solve\_hybrid}（耦合混合潮流，非纯 AC PF），选项
\fld{max\_iter}$=\max(100,\text{inner})$、\fld{tol}$=\min(\fld{feas\_tol},10^{-8})$，
开启 converter mode switching / PV$\to$PQ / auto swing（ac\_opf.cpp:2305--2326）；
除 $(v_m,v_a,v_{dc})$ 外还从 \fld{pf.vsc\_transfers}/\fld{pf.dcdc\_transfers}
按组件 index 提取 $P_{ac}/Q_{ac}/P_{dcdc}$ 种子（:2052--2093）。

\paragraph{parity 原变量暖启动 \srcpath{build\_parity\_primal\_warm\_start}（:1874--2003）}
\begin{itemize}
  \item 优先消费 \fld{ipm\_primal\_state}（维数$==n_{total}$ 且 allFinite），
        裁剪到 $[x_{min}+10^{-8},\ x_{max}-10^{-8}]$（:1886--1901）；
  \item 退化路径按公共物理量逐族映射（vm/va/vdc/pac/qac/dpd/dqd/pren/qren/pdcdc/
        pflex/erp），任意族缺失则回退物理初值（:1903--1947）；
  \item $P_{dc}$ 由 $P_{ac}/Q_{ac}$ 经损耗式重建：
        $p_{dc}=-(p+a+bI+cI^2)$，$I=\sqrt{p^2+q^2+\max(\fld{eps\_iac},10^{-12})}/v_m$
        （:1919--1939）；
  \item 合成 external-grid 机组按 $\fld{data\_index}-\fld{pg\_mw.size()}$ 偏移识别并
        映射到 \fld{external\_grid\_p/q\_mw}（:1948--1972）；
        ER 无功按 \fld{port.qvar} 映射（:1992--2000）；
  \item 对偶三族（equality/inequality/slack）原样透传（:2327--2350）；
        IPM 侧兼容性检查：尺寸 + allFinite + minCoeff$>$0（parity\_ipm.cpp:910--919）。
\end{itemize}
生产级用例：时序流水线逐时步 \fld{warm\_start=\&previous}（见第~\ref{sec:applications}~章）。

\subsection{Native AC 简约空间旧路径}
\srcpath{OPFSolverPath::NativeAC}（ac\_opf.cpp:3112--3793）是不含 DCDC 调度的
简约空间原对偶 IPM。

\paragraph{变量与等式}
变量布局（\srcpath{build\_index}，:224--343）：
$x=[\theta_{\text{非slack}};\ v_{m,\text{PQ}};\ p_g;\ q_g;\ v_{dc};\ p_{ac};\ q_{ac};\ p_{dc}]$；
等式 $=[P\text{平衡}(n_p),\ Q\text{平衡}(n_q),\ DC\text{平衡}(n_{dc}),\ \text{换流器平衡}(n_c)]$。

\paragraph{目标}
二次发电成本 $f=\sum_g(c_2P_g^2+c_1P_g+c_0)$（:539--549），梯度/对角 Hessian 按 pu 换算（:551--571）。

\paragraph{约束处理（诚实口径）}
\begin{itemize}
  \item 变量界 $\to$ 对数障碍 $-\mu(\log s^-+\log s^+)$（:573--592）；
  \item 非线性限值 $\to$ \textbf{外部二次罚}（非障碍）：$\tfrac12\rho\max(h,0)^2$，
        支路两端 $h=P_f^2+Q_f^2-S_{max}^2$（from 端 :899--924、to 端 :926--951，
        $\pi$ 模型含 tap \textbf{与} shift），换流器容量圆 $P_{ac}^2+Q_{ac}^2-S_{max}^2$（:954--958）；
        \textbf{相移修复}：\fld{BranchLimitEntry} 新增 \fld{shift\_rad} 字段（:646），
        装配时 \fld{e.shift\_rad} = \fld{shift\_deg}$\times$\fld{kDegToRad}（:685），
        非线性支路限值评估的角度差为 $\theta_i-\theta_j-\mathrm{shift}$（:891--893）——
        修复前带相移支路（复变比 $|t|\cdot e^{j\cdot\mathrm{shift}}$）的 $S$ 限值评估
        漏掉相移角、与 Ybus/支路潮流装配不一致，修复后一致；
        罚系数 $\rho=100\cdot\fld{merit\_penalty}\cdot\fld{ineq\_gate}$，
        $\fld{ineq\_gate}=\mathrm{clamp}(10^{-2}/\max\_eq\_mismatch,\ 10^{-3},\ 1)$（:3262--3263）；
  \item DC 平衡 $P_{dc,k}=p_{d,k}+V_k(G_{dc}V)_k-\sum p_{dc}$（:738--752）；
        VSC 耦合 $p_{ac}+p_{dc}+a+bI_{ac}+cI_{ac}^2=0$，
        $I_{ac}=\sqrt{P_{ac}^2+Q_{ac}^2+\varepsilon}/V_m$，$a=\fld{loss\_mw}/S_b$、
        $b=\fld{loss\_percent}/100$、$c=1-\eta$（:754--769）；
  \item VSC 容量圆半径为\textbf{推断值}：$s_{max}=\max(|p_{max}|,|p_{min}|,|q_{max}|,|q_{min}|)$
        （无独立 s\_rated 字段，:689--701；parity 路径不同，优先 \fld{p\_rated\_mw}，
        见第~\ref{sec:parity}~章）。DC\_V 母线 $v_{dc}$ 界 $v_{set}\pm0.05$，
        其余 DC 母线 0.8--1.2（:414--427）。
\end{itemize}

\paragraph{KKT 与全局化}
KKT 系统 $\begin{bmatrix}H_{diag}+\delta I & J^{\mathsf T}\\ J & -\delta I\end{bmatrix}$
（:1318--1429），线性求解器 \srcpath{core::make\_default\_sparse\_solver()}。
全局化：6 次信赖尝试、正则化 $\times10$ 递增（:3270--3425）；
$\Delta x$ 限幅 $0.25\times$ 区间宽且 $\infty$-范数 $\le0.25$，$\Delta\lambda$ 限 $\pm10^3$
（:3306--3318）；fraction-to-boundary + merit 滤波接受
（不可行度下降 $(1-5\times10^{-3}\alpha)$ 或近可行时 merit 下降，:3387--3397）；
失败后依次：恢复阶段（纯可行性步，$H_{diag}=0$，:3427--3542）$\to$
梯度回退步 $\Delta x=-J^{\mathsf T}g$（:3544--3619）。
$\mu$ 更新：内层收敛或近可行时 $\mu=\max(\mu\cdot\fld{barrier\_mu\_reduction},\ \mu_{min})$，
否则加大罚与正则化（:3673--3681）。
收敛判据（:3251--3257）：
\[
\max\_mismatch\le\fld{feasibility\_tol}\ \land\
\max\_stationarity\le\fld{stationarity\_tol}\ \land\ \mu\le1.01\,\mu_{min}.
\]
不可行度 $\le10^{-3}$ 时用对偶最小二乘估计 stationarity（:598--626, 3243--3245）。
动态权重：\fld{objective\_weight} 初值 0.05，停滞/恢复时自适应（:3185, 3650--3664）。

\paragraph{全局化实测常量补充}
merit 函数显式形式 $\mathrm{merit}=w_{obj}\,f+\mathrm{barrier}+\tfrac12\rho(\lVert g\rVert^2+\lVert h^+\rVert^2)$
（:3260--3262）；near-neutral 步接受：近可行且 infeas/merit 各 $\le1.005\times$
现值时接受但 $\alpha$ 限 $5\times10^{-3}$（:3401--3417）；
\fld{stagnation\_count}$\ge3$ 且不可行度 $>\max(10^{-3},10\cdot\fld{feas\_tol})$
时强制跳过信赖尝试直进恢复（:3273--3275）；
接受后 $x$ 重夹：受约束列 $10^{-3}\times$width、未受约束列 $10^{-8}\times$width
（:3629--3641）；\textbf{每步接受后 $\lambda$ 清零}（:3643，旧路径不带对偶续传）；
动态参数：恢复接受后 $\rho\times1.2/\delta\times1.5/w_{obj}\times0.9$；
停滞 $\ge4$ 时 $\rho\times2/\delta\times5/w_{obj}\times0.8$；进展好（$<0.8\times$）
时 $\rho\times0.95$ 回落；不可行度 $<10^{-2}$ 时 $w_{obj}\times1.15$ 上限 1.0
（:3650--3664）；外层退出：近可行阈值 $\fld{max\_viol}\le\max(10^{-4},50\cdot\fld{feas\_tol})$
时也减 $\mu$（:3673--3681）；恢复阶段限幅 $0.10\times$width、$\infty$-范数 0.12，
梯度回退限 0.08（:3429, 3547）；
失败 status 诊断格式 \texttt{[max|P|=…, max|Q|=…, max|DC|=…, max|Conv|=…, max(h+)=…,
worst\_stat=varname:val]}（:3780--3784）。

\subsection{参数表}
ACOPFOptions 全字段见第~\ref{sec:options}~章表~\ref{tab:acopfopts}；
$\le0$ 的数值字段在 ac\_opf.cpp:2713--2748 重置为默认。
与 IPMOptions 的映射：\fld{max\_iter} = inner$\times$outer、
\fld{tol\_complementarity} = \fld{barrier\_mu\_min}、
\fld{alpha\_max} = \fld{interior\_fraction}（ac\_opf.cpp:2292--2299）。

\subsection{验证与校核}
\begin{itemize}
  \item \srcpath{tests/test\_opf\_solver\_backends.cpp}：case300\_acdc parity 收敛
        （:79--105）；Ipopt 选项透传（:107--142）；暖启动目标差 $\le10^{-5}$（:144--170）；
        混合微网 DC 参考 $|v_{dc}-v_{set}|<2\times10^{-6}$（:172--195）；
        裸 DC\_V 拒绝（:244--265）；case2000\_acdc $<160$ 迭代（:303--335）；
        case2383wp 目标缩放回归（:481--494）；case30 后端语义与
        Parity 成本 $\le$ ED$\times1.001+1$（:498--553）；case2869pegase $<100$ 迭代
        稀疏后端（:577--599）；case6515rte 不发散（:601--618）。
  \item 经济调度兜底：case30 ``economic-dispatch'' 状态断言（test\_opf\_solver\_backends.cpp:498--535）。
  \item 雅可比诊断：实际 FD 诊断 \srcpath{check\_ac\_opf\_jacobian\_fd} 目前无调用方
        （死代码，见本章诚实口径注）；DC slack KCL Jacobian FD $<10^{-6}$
        （test\_opf\_solver\_backends.cpp:267--301）。
\end{itemize}

% 第 3 章：Parity 全空间 NLP 建模
\section{Parity 全空间 NLP 建模}\label{sec:parity}

\subsection{功能定位与入口}
\compmeta{opf::parity（ParityProblem）}{\srcpath{include/hacdcpf/optimal\_power\_flow/formulation.hpp}}
Parity formulation 是供原生 IPM 与 Ipopt 共用的“可审查”全空间 NLP
（formulation.hpp:11--15）：
\[
\min_x\ f(x)\quad \mathrm{s.t.}\quad g(x)=0,\quad h(x)\le 0,\quad x_{min}\le x\le x_{max}.
\]
装配入口 \srcpath{parity::build\_problem}（src/optimal\_power\_flow/parity\_formulation.cpp:108）。
它是\textbf{唯一同时优化 AC/DC/VSC/DCDC/储能/可切负荷/可再生/柔性负荷/能量路由器的路径}；
混合算例由顶层自动强制启用。

\subsection{变量布局}
\srcpath{VarIndex}（formulation.hpp:66--87；装配 parity\_formulation.cpp:338--406）：
\begin{align*}
x=[\,&\theta;\ V_m;\ P_g;\ Q_g;\ V_{dc};\ P_{ac};\ Q_{ac};\ P_{dc};\ \Delta P_d;\ \Delta Q_d;\
P_{ren};\ Q_{ren};\\
&P_{stor};\ Q_{stor};\ P_{storDC};\ P_{dcdc};\ P_{flex};\ ER_P;\ ER_Q].
\end{align*}
纯 AC 与混合统一；空家族长度 0 但偏移有效。
等式行序（\srcpath{ConstraintIndex}，formulation.hpp:103--122；计数 :408--467）：
$[gP_{ac};\ gQ_{ac};\ gP_{dc};\ g_{vsc};\ g_{dcdc}(=0);\ g_{er};\ g_{dc\_ref}]$
（$n_{dcdc\_bal}$ 仍分配但评估器留空，实现怪癖见 11c 节技术笔记）。
非线性不等式行序：from 端 $|S|^2$、to 端 $|S|^2$、VSC 容量圆、DC 支路 $P^2$、
VSC $I_{ac}^2$、$m$ 上限、$m$ 下限、DC/DC 占空比（每受限台 2 行线性行）。

\subsection{目标函数}
\srcpath{objective}（parity\_formulation.cpp:1138--1240；梯度/Hessian :1242--1329）：
\begin{itemize}
  \item \textbf{Economic}：
        \begin{align*}
        f={}&\sum_g(c_2P_g^2+c_1P_g+c_0)+10^{-3}\sum_c(Q_{ac}-q_{set})^2
        +\mathrm{VOLL}\sum(\Delta P_d+\Delta Q_d)\\
        &{}+\sum \text{固定DER线性成本}
        +\sum_{ren}\big[c_1P+c_{curt}(P_{rated}-P)\big];
        \end{align*}
        无功设定点正则化常数 \texttt{kReactiveSetpointRegularization}$=10^{-3}$（:28，
        注释说明动机：消除平坦方向）；VOLL 自动值
        $\max(100,\ \text{负荷}\ \fld{cost\_mw},\ 1.1\times\max\text{边际成本})$（:64--83）；
  \item \textbf{VoltageDeviation}：$w_v\sum_i(V_i-V^*)^2$；
  \item \textbf{ActiveLoss}：$w_l\,(\sum P_g+P_{ren}+P_{stor}+P_{storDC}-P_{flex}+\Delta P_d)\,S_b$
        （物理损耗 = 发电 $-$ 服务负荷，:1155--1171）；
  \item \textbf{VoltageDeviationAndLoss}：组合目标（:1142--1147，
        RPO Combined 内层使用，reactive\_power\_opt.cpp:344--345）；
  \item \textbf{固定 DER 线性成本是常数项}（不影响极小解，仅影响报告值）：
        AC/DC static gen、不可削减可再生、不可控 PV 的 $c_1\max(0,P\cdot\fld{scaling})$，
        DC PV 阵列用 \fld{p\_set\_mw}（:1200--1216）；
  \item \textbf{可再生削减成本}：仅当 \fld{cost\_curtail\_mwh}$>0$ 且 $p_{rated}-p>0$
        时收取，梯度 $(c_1-\max(0,c_{curt}))\,S_b$（:1231--1233, 1322--1323）；
        注意削减量对 $p_{rated}$ 计量而上界是 $\min(\text{额定},\text{可用出力})$——
        可用$<$额定时差额恒计入成本；PV 系统只有 $c_1$ 无削减价（:1234--1237）；
  \item 非经济目标加 $10^{-6}$ 经济 tie-break（:1177--1183）。
\end{itemize}

\subsection{等式约束}
装配 :1335--1572。

\paragraph{固定注入台账（\fld{p\_fixed\_inj}/\fld{q\_fixed\_inj}，:502--564）}
AC 侧固定注入组成：StaticGenerator（$\times$scaling）、不可削减 RenewableGen、
不可控 PVSystem（经 \srcpath{compute\_pv\_power\_mw} 功率曲线，:533）、
VPP（pcc 出力）、并网 Microgrid（仅 \fld{p\_exchange}）、非 InTransit 的
MobileStorage。DC 侧平衡另有组成：\fld{dc\_loads} 存在时逐负荷
\srcpath{effective\_load\_p\_mw}、dc\_static\_generators、dc\_pv\_arrays 注入
（:1453--1488）。

AC 平衡（ZIP 负荷，\fld{scale\_p}/\fld{scale\_q}
$=1/\max(|P_d|,1)$ 预缩放，:1431--1440）：
\begin{align*}
gP_i={}&P_i^{net}-P_{g,i}-P_{fixed,i}-P_{ren,i}-P_{stor,i}\\
&{}+P_{d,i}(V)-\Delta P_{d,i}+P_{flex,i}-P_{ac,i}=0
\end{align*}
$Q$ 行同构。DC 平衡：
\[
gP_k=p_{d,k}+V_k\sum_m G_{km}V_m-P_{conv\_dc,k}=0;
\]
DC/DC 折叠进 DC 行：输入侧 $+p_{in}$，输出侧 $-p_{out}$，
$p_{out}=p_{in}\eta$（正向）或 $p_{in}/\eta$（反向），加 $I^2R$ 损耗
$r_{eq}p_{in}^2/v_{dc,in}^2$（:1497--1514）。
VSC 耦合：$P_{ac}+P_{dc}+a+bI_{ac}+cI_{ac}^2=0$（:1516--1530，$\varepsilon=\fld{eps\_iac}=10^{-6}$）。
能量路由器：端口功率为决策变量，每路由器一条\textbf{无损}守恒行 $\sum P_{port}=0$
（:1548--1566；注释明确 \fld{loss\_percent} 细化 deferred，:1563--1564）。
DC\_V 参考电压行 $V_{dc,k}-v_{set}=0$（:1567--1571）；
无导电支路的 DC 母线通过 $\pm10^{-6}$ 界锚定（:613--627）。

\subsection{不等式约束}
$h(x)\le0$ 约定，:1824--1925：
\begin{align*}
\text{AC 支路两端：}\quad &P_f^2+Q_f^2-S_{max}^2\le0\quad(\pi\text{ 模型，含 tap 与 \fld{shift\_deg}，:1836--1861})\\
\text{VSC 容量圆：}\quad &P_{ac}^2+Q_{ac}^2-S_{max}^2\le0\quad(:1862--1867)\\
\text{DC 支路：}\quad &P_{km}^2-P_{max}^2\le0,\quad P=gV_k(V_k-V_m)\quad(:1869--1879)\\
\text{VSC 电流限：}\quad &P_{ac}^2+Q_{ac}^2-i_{ac,max}^2V_m^2\le0\quad(:1881--1890)\\
\text{调制度（线性行）：}\quad &m=\frac{V_mV_{n,ac}}{K_mV_{dc}V_{n,dc}}\ \text{上下限}\quad(:1891--1908)\\
\text{DC/DC 占空比（线性行）：}\quad &\text{见下方全拓扑公式表}\ (:1797--1818,\ 1909--1924)
\end{align*}
注意口径：AC 支路限为 $|S|$（视在功率）双端圆约束；DC 支路限为 $|P|$。

\paragraph{换流器容量圆半径 \srcpath{converter\_smax\_pu}（:85--101）}
优先 \fld{conv.p\_rated\_mw}，无效时回退
$\max(|p_{max}|,|p_{min}|,|q_{max}|,|q_{min}|,10^{-6})$；
\fld{enforce\_converter\_capacity}=false 时返回 $10^{12}$ 使圆恒不激活。

\paragraph{准入条件（未声明限值的设备不占行，:435--464）}
$I_{ac}$ 行需 \fld{isfinite(i\_ac\_max\_pu)}$>0$；$m$ 行需
\fld{k\_m\_modulation}$>0\ \land\ v_{n,ac}>0\ \land\ v_{n,dc}>0\ \land\ m_{max}>0$
（$m_{min}$ 行同理）；占空比行需 topology$\ne$Generic 且 $0\le d_{min}<d_{max}\le1$；
端口额定 kV 取 \fld{vn\_in/out\_kv}，缺省回退母线 base\_kv 再回退 1.0（:1791--1795）。

\paragraph{DC/DC 占空比行全拓扑公式（\srcpath{dcdc\_duty\_rows}，:1797--1818）}
每受限台 2 行 $h=a_{in}V_{in}+a_{out}V_{out}\le0$（上限行 $D\le d_{max}$、
下限行 $D\ge d_{min}$），$n$=\fld{n\_ratio}（默认 1）：
\begin{center}
\footnotesize
\begin{tabular}{@{}llll@{}}
\toprule
拓扑 & 电压律 & $D\le d_{max}$ 行 $(a_{in},a_{out})$ & $D\ge d_{min}$ 行\\
\midrule
Buck & $V_{out}=DV_{in}$ & $(-d_{max}Vn_{in},\ Vn_{out})$ & $(d_{min}Vn_{in},\ -Vn_{out})$\\
Boost & $V_{out}=V_{in}/(1-D)$ & $(-Vn_{in},\ (1-d_{max})Vn_{out})$ & $(Vn_{in},\ -(1-d_{min})Vn_{out})$\\
BuckBoost & $V_{out}=\tfrac{D}{1-D}V_{in}$ & $(-d_{max}Vn_{in},\ (1-d_{max})Vn_{out})$ & $(d_{min}Vn_{in},\ -(1-d_{min})Vn_{out})$\\
Isolated & $V_{out}=nDV_{in}$ & $(-d_{max}nVn_{in},\ Vn_{out})$ & $(d_{min}nVn_{in},\ -Vn_{out})$\\
\bottomrule
\end{tabular}
\end{center}

\subsection{变量界与初值}
界（:569--754）：$\theta\in[-\pi,\pi]$；$V_m$ 母线限；$P_g/Q_g$ 机组限
（退役机组钉死 $\pm10^{-8}$）；储能钉在调度值 $\pm10^{-4}/S_b$
（无跨时段 SOC 模型的诚实说明，:674--717）；\textbf{例外}：
\fld{cap\_charging\_strategy}==``static'' 时 $P$ 钉在 $-$\fld{cap\_static\_charging\_mw}、
$Q$ 钉 0（$\pm10^{-4}/S_b$ 容差带），即承载力静态 ESS 完全退出优化；
其余储能钉调度值 \fld{p\_mw}$\pm$eps 且同时受 $[p_{min},p_{max}]$ 夹（:679--717）；
可再生上限 $=\min(\text{额定},\text{可用出力})$
（:651--672）；柔性负荷 $[p-\fld{flex\_down},\ p+\fld{flex\_up}]$（:730--734）；
DC/DC 界 $p_{max}=\fld{pmax\_mw}>10^{-9}\,?\,\fld{pmax\_mw}:\fld{sn\_mva}$，
$p_{min}\ge p_{max}$ 时取 $-p_{max}$（:719--727）。

初值链：母线设定点 $\to$ 零出力机组按 $P^{max}$ 比例分摊缺额（:940--966）
$\to$ DC 潮流暖启动（稀疏 $B'$ 求解，:757--857）$\to$ $Q_g$ 暖启动（:861--917）
$\to$ 切负荷吸收残差（:1038--1095）$\to$ $J_gJ_g^{\mathsf T}$ 最小二乘可行性修正
（阻尼 0.8，:1097--1126）$\to$ 1\% 内点钳制。
DC/DC 暖启动 \fld{p\_ref} 是输出参考，经 $\eta$（clamp $[0.01,1]$）换算成输入参考
（:1004--1012）；ER 端口注入在 AC 行缩放之后补乘 \fld{scale\_p}/\fld{scale\_q}
（:1545--1556）。

\subsection{解析导数}
解析等式 Jacobian（:1574--1780）、不等式 Jacobian（:1927--2071）、
稀疏 Lagrangian Hessian $\nabla^2L=\nabla^2f+\sum\lambda\nabla^2g+\sum\nu\nabla^2h$
（:2087--2560，含 ZIP 二阶项、VSC 损耗全二阶、DC/DC $I^2R$ 二阶、支路 $|S|^2$ 全二阶）。

\paragraph{ZIP 系数与二阶导}
系数：有组件负荷时用 SolverData 逐母线加权 \fld{bus\_zip\_*}，否则全局
\fld{zip\_pw}/\fld{zip\_qw}（默认恒功率）（:133--160）；
Hessian 中 ZIP 二阶项 $2w_2p_d$（P 行）、$2w_2q_d$（Q 行）加在 $V_m$ 对角
（:2235--2239），且 AC 平衡行曲率乘 $\lambda\cdot\fld{scale\_p}/\fld{scale\_q}$
（:2187--2188）。

\paragraph{Hessian 稀疏结构与关键二阶项}
下三角 triplet 双边写入、去重求和装配；$|\lambda|$ 或 $|\nu|<10^{-14}$ 的行整块跳过
（:2189 等）；调制/占空比行线性、零曲率，仅用游标推进保持对偶对齐（:2548--2553）。
DC/DC $I^2R$ 项（$g_{out}$ 中 $+rp_{in}^2/v^2$）：
$\partial^2/\partial p^2=2r/v^2$、$\partial^2/\partial p\partial v=-4rp/v^3$、
$\partial^2/\partial v^2=6rp^2/v^4$（:2319--2351，$v$ 取 $\max(v_{in},0.1)$）；
$I_{ac}$ 限行 $\partial^2h/\partial V_m^2=-2i_{max}^2$（:2546）。

正确性校核：FD 诊断 \srcpath{check\_ac\_opf\_jacobian\_fd}（死代码）与测试中的
FD 断言（见第~\ref{sec:validation}~章）。

\subsection{参数表（ParityOptions）}
定义 formulation.hpp:35--49。
\begin{paramtable}
\fld{load\_shedding} & -- & bool & true/false & true & 启用可切负荷变量 $\Delta P_d,\Delta Q_d$\\
\fld{voll} & VOLL & \$/MWh & $\ge100$ & 0（自动） & 切负荷价值；0 = $\max(100,1.1\times\max$边际成本$)$\\
\fld{eps\_iac} & $\varepsilon$ & pu & $10^{-8}$--$10^{-4}$ & $10^{-6}$ & VSC 电流数值正则化\\
\fld{objective} & -- & 枚举 & 四档 & Economic & Economic / VoltageDeviation / ActiveLoss / VoltageDeviationAndLoss\\
\fld{voltage\_target\_pu} & $V^*$ & pu & 0.95--1.05 & 1.0 & 电压偏差目标\\
\fld{voltage\_deviation\_weight} & $w_v$ & -- & $\ge0$ & 1.0 & 电压偏差目标权重\\
\fld{active\_loss\_weight} & $w_l$ & -- & $\ge0$ & 1.0 & 网损目标权重\\
\fld{enforce\_branch\_limits} & -- & bool & -- & true & AC 支路 $|S|$ 双端约束\\
\fld{enforce\_converter\_capacity} & -- & bool & -- & true & VSC 容量圆\\
\fld{enforce\_converter\_current\_limits} & -- & bool & -- & true & VSC $I_{ac}$ 限\\
\fld{enforce\_converter\_modulation\_limits} & -- & bool & -- & true & 调制度 $m$ 与 DC/DC 占空比行\\
\end{paramtable}

\subsection{验证与校核}
\begin{itemize}
  \item parity 成本不劣于经济调度 $\times1.001+1$（test\_opf\_solver\_backends.cpp:537--553）；
  \item 夜间纯 external-grid 算例：储能钉调度值 WithinAbs $2\times10^{-4}$、
        DC/DC $I^2R$ 公式复核、AC 供需平衡 $\pm10^{-3}$
        （test\_nighttime\_opf.cpp:46--96）；
  \item DC slack KCL Jacobian FD $<10^{-6}$（test\_opf\_solver\_backends.cpp:267--301）；
  \item 结果中的 \fld{converter\_model\_scope} 按实际启用约束族拼装标签
        （ac\_opf.cpp:2254--2272）。
\end{itemize}

% 第 4 章：原生原对偶内点法（Parity IPM）
\section{原生原对偶内点法（Parity IPM）}\label{sec:ipm}

\subsection{功能定位与入口}
\compmeta{opf::parity::solve\_primal\_dual\_ipm}{\srcpath{include/hacdcpf/optimal\_power\_flow/native\_ipm\_solver.hpp}}
实现 src/optimal\_power\_flow/parity\_ipm.cpp:852--2149。对第~\ref{sec:parity}~章的
全空间 NLP 求 KKT 点：原变量 $x$、等式对偶 $\lambda$、不等式对偶 $\mu$、松弛 $z>0$，
障碍子问题
\[
\min\ f(x)-\gamma\sum_j\log z_j\quad\mathrm{s.t.}\quad g(x)=0,\quad h(x)+z=0,\quad x_{min}\le x\le x_{max}.
\]

\subsection{初始化与目标缩放}
不等式装配 \srcpath{assemble\_inequalities}（:799--844）：非线性行 + 有限上下界行。
松弛初始化 $z=\max(r_h,1)$ 或 $\max(-r_h,10^{-2})$，乘子 $\mu=\max(1/z,10^{-2})$
（:924--941）。目标缩放
\[
\fld{obj\_scale}=\min\bigl(1,\ 100/\lVert\nabla f\rVert_\infty\bigr)\quad(:957--987),
\]
治理大电网成本梯度淹没平衡行的问题（case2383wp 回归，见验证节）。

\subsection{收敛度量（归一化）}
:1014--1027：
\begin{align*}
\mathrm{feascond}&=\frac{\max(\lVert g\rVert_\infty,\ \max h^+)}{1+\max(\lVert x\rVert_\infty,\lVert z\rVert_\infty)},
&\mathrm{gradcond}&=\frac{\lVert\nabla f+J_g^{\mathsf T}\lambda+J_h^{\mathsf T}\mu\rVert_\infty}{1+\max(\lVert\mu\rVert,\lVert\lambda\rVert)},\\
\mathrm{compcond}&=\frac{z^{\mathsf T}\mu/n_{iq}}{1+\lVert x\rVert_\infty},
&\fld{comp\_tol}&=\max(10\,\mathrm{tol},\ 2\times10^{-3})\ (:1036).
\end{align*}
另保留未归一化的 \fld{raw\_feascond}。

\subsection{Newton 系统两种形式}
\srcpath{KktForm}（:167--199），env \texttt{HACDCPF\_OPF\_KKT\_FORM}，默认 condensed：
\begin{itemize}
  \item \textbf{Condensed}：$W=L_{xx}+J_h^{\mathsf T}\Sigma J_h$，$\Sigma=\mathrm{diag}(\mu/z)$
        （:1176--1189）；KKT
        $\begin{bmatrix}W+\delta I & J_g^{\mathsf T}\\ J_g & -\delta I\end{bmatrix}$（:443--481）；
        递进正则化 $\{0,10^{-8},10^{-6},10^{-4}\}\times\lVert W\rVert$（:1467--1473）。
        RHS 装配（:1461--1544）：预估 $N_{aff}=L_x+J_h^{\mathsf T}(\mu\circ r_h/z)$；
        校正 $N_{corr}=L_x+J_h^{\mathsf T}\big((\mu\circ r_h+\gamma e-\mathrm{cross})/z\big)$；
        恢复 $dz=-(r_h+z)-J_h\,dx$，
        $d\mu=-\big(\mu+\Sigma dz+(\mathrm{cross}-\gamma e)/z\big)$；
  \item \textbf{Augmented}：
        $\begin{bmatrix}L_{xx}+\delta_W I & J_g^{\mathsf T} & J_h^{\mathsf T}\\
        J_g & -\delta_C I & 0\\ J_h & 0 & -ZM^{-1}\end{bmatrix}$（:562--629）；
        Wächter--Biegler 惯性校正：MUMPS 负主元数须 $=n_{eq}+n_{iq}$，否则 $\delta_W$
        从 $10^{-8}\lVert L_{xx}\rVert$ 起步 $\times8$ 增至 $10^{-2}\lVert L_{xx}\rVert$ 上限
        （:1281--1301）；$\delta_C$ 默认 $10^{-8}\lVert L_{xx}\rVert$（:1232--1239，
        可由 \texttt{HACDCPF\_OPF\_DELTA\_C} 绝对值覆盖）。
        \textbf{$\delta_W$ 暖启动与 KLU/MUMPS 不对称}（:1271--1319）：
        仅 MUMPS 步用 $\delta_W=\max(10^{-8}\lVert L_{xx}\rVert,\ 0.5\,\delta_{W,prev})$
        （Ipopt $\kappa_W^-$ 规则）；KLU 步恒 $\delta_W=0$（惯性不可报，且防止
        $\delta_W$ 泄漏污染 KLU 轨迹，实测 case13659 +44 迭代）；
        惯性循环最多 8 次、第 8 次封顶接受；\texttt{HACDCPF\_OPF\_INERTIA\_GATE}=N
        门控的周期 MUMPS 复检后\textbf{总是恢复 KLU}。
\end{itemize}

\subsection{线性代数后端与数值健壮性}
后端优先级 \srcpath{backend\_order}（:268--327）：结构感知——纯 AC 先 KLU
（BTF 结构实测最快），混合 AC/DC 先 MUMPS（需 $LDL^{\mathsf T}$ 惯性）；
UMFPACK/Eigen SparseLU 为升级链；env \texttt{HACDCPF\_OPF\_LINEAR\_SOLVER} 可钉死；
Windows 上 KKT 维数 $\le1024$ 自动走 dense PartialPivLU（:1069--1078）。
健壮性措施：对称 Ruiz 均衡 5 次扫描、每次求解冻结（:415--438, 503--517）；
2 步迭代精化带发散守卫（:655--671）；Eigen SparseLU 残差超
$10^{-6}\lVert rhs\rVert$ 标记 \fld{solve\_degraded} 触发后端升级（:677--683）。

\subsection{步长计算与全局化}
\begin{itemize}
  \item \textbf{Mehrotra 预估-校正}：预估 $\gamma=0$，$\sigma=(\mu_{aff}/\mu)^3$
        （:1500--1507）；校正 cross 项限幅 $\pm0.1\gamma$（:1513--1518）；
        Gondzio 附加校正 $\le2$ 次（:1554--1591）；
        \textbf{校正器再中心化重试}（:1375--1450）：校正方向若 ftb 步长 $<10^{-9}$
        （$\kappa\sim1/\mu$ 爆炸征兆），则 $\gamma_{corr}=\max(\gamma,\ 0.5\,\mu_{cur})$、
        cross 清零重解，$\le2$ 次。
        \textbf{ABO 动态 $\mu$ 策略}（可选，env \texttt{HACDCPF\_OPF\_MU\_STRATEGY}，
        :1343--1359, 1454--1456）：$\mu^+=\mu^2/\mu_0$（$\theta_S$ 更新，$\gamma_q=1$，
        $b=\mu_0$），$\gamma=\mathrm{clamp}\big(\mu+\alpha_{aff}(\mu^+-\mu),\ 10^{-14},\ 10^4\big)$，
        $\alpha_{aff}=\min(\alpha_{p,aff},\alpha_{d,aff},1)$；
        $\mu_{dyn}/\mu_0$ 首迭代初始化为当前 $\mu_{cur}$，校正步后
        $\mu_{dyn}=\gamma_{corr}$（含再中心化）；
  \item \textbf{$(\theta,\varphi)$ 滤波线搜索}（Wächter--Biegler §3.2）：
        常数 $\gamma_\theta=\gamma_\varphi=10^{-5}$、$\Delta=1$、$s_\theta=1.1$、
        $\eta_\varphi=10^{-8}$、$\theta_{min}=10^{-4}$、
        $\theta_{max}=10^4\max(1,\text{初始不可行度})$（:1625--1641）；
        domination 记忆 \fld{filter\_history}（:1731--1776）；
        $\varphi_\gamma=\tilde f-\gamma\sum\log z$（:1669--1681）；
        仅当 KKT 维数 $\ge512$ 或 env \texttt{HACDCPF\_OPF\_FILTER\_ALL} 启用（:1646--1653）；
        另有 legacy 三轴接受路径（feas/stationarity/complementarity 进展，
        带对偶轴 $\times100$ 防爆守卫，:1707--1729）；
  \item \textbf{其他}：二阶校正 SOC（仅 augmented，精确触发：bt==0 且
        $\theta_{trial}>\theta_k$ 且 meq$>$0，rhs $=-(\alpha r_g+r_g(x+\alpha dx))$，
        :1800--1874）、中心性恢复重试（:1890--1907）、恢复阶段 $\rho I$
        Gauss--Newton 可行性步（见下节）、
        空步规则 $\min(\alpha_p,\alpha_d)<10^{-7}$ 判方向不可用（:1600--1608）。
\end{itemize}
另：\fld{obj\_scale} 的 Hessian 修正 $(\fld{obj\_scale}-1)\cdot h_{diag}$ 加在
$L_{xx}$ 对角（:973--987）。

\subsection{终止、认证与同伦切线}
\begin{itemize}
  \item \textbf{best-iterate 跟踪}（:1993--2005）与\textbf{可接受收敛}
        （:2038--2074）：feas/raw/grad $<100\times$tol 且 comp $<$ \fld{comp\_tol}，
        best-vs-final 择优，status 区分 ``best iterate''/``final iterate''；
  \item \textbf{目标停滞提前停止}（:2007--2026）：$\fld{raw\_feas}<10^{-4}\ \land\
        \mathrm{comp}<\fld{comp\_tol}$ 且 $|\Delta obj|\le10^{-6}(1+|obj|)$
        连续 20 接受步，status=``objective-stagnant (dual-polish certification pending)''；
  \item \textbf{恢复阶段}（:1909--1962）：$(1,1)$ 块 $\rho I$，
        $\rho=10^{-8}\lVert L_{xx}\rVert$，$\delta_C=10^{-8}\lVert L_{xx}\rVert$，
        尾部 rhs 用 $\gamma_{corr}$；$\le8$ 次 $\alpha$ 减半，接受条件 $\theta$ 改善
        $(1-10^{-4}\alpha)$ 且对偶轴 $\times100$ 有界；
  \item \textbf{对偶最小二乘抛光}（:2076--2126）：解
        $\begin{bmatrix}I & J_g^{\mathsf T} & J_h^{\mathsf T}\\ J_g & 0 & 0\\
        J_h & 0 & 0\end{bmatrix}(r,\lambda,\mu)=(-\tilde f\nabla f,\ 0,\ 0)$
        （$z=0$、$\mu=e$ 的增广 KKT，$\delta_C=10^{-8}$），$\mu$ 裁 $\ge0$；
        采用规则 $\fld{stat\_ls}\le\fld{gradcond}$ 或未收敛即采用；
        认证 status=``converged (dual least-squares certified)''；dense 路径跳过。
\end{itemize}
\srcpath{homotopy\_tangent}（:2154--2257）：一次惯性控制 KKT 求解
$K_{aug}\,(dx/dt,d\lambda/dt,d\mu/dt)=(-\nabla f/t,\ 0,\ 0)$，$dz/dt=-J_h\,dx/dt$，
供目标同伦 Davidenko 预测（第~\ref{sec:acopf}~章）。

\subsection{参数表（IPMOptions）}
定义 native\_ipm\_solver.hpp:17--38。
\begin{paramtable}
\fld{max\_iter} & -- & 次 & 100--1000 & 400 & IPM 最大迭代数（AC OPF 映射为 inner$\times$outer）\\
\fld{tol\_primal} & -- & -- & $10^{-8}$--$10^{-4}$ & $10^{-6}$ & 原始可行容差\\
\fld{tol\_dual} & -- & -- & $10^{-8}$--$10^{-4}$ & $10^{-6}$ & 对偶（stationarity）容差\\
\fld{tol\_complementarity} & -- & -- & $10^{-8}$--$10^{-4}$ & $10^{-6}$ & 互补容差（映射自 \fld{barrier\_mu\_min}）\\
\fld{regularization} & $\delta$ & -- & $10^{-10}$--$10^{-4}$ & $10^{-8}$ & KKT 递进正则化起步\\
\fld{alpha\_max} & $\alpha_{max}$ & -- & 0.9--0.999 & 0.95 & 最大步长（映射自 \fld{interior\_fraction}）\\
\fld{verbose} & -- & bool & -- & false & 迭代日志\\
\fld{primal\_start} 等四项 & -- & 向量 & -- & 空 & 暖启动 $x/\lambda/\mu/z$\\
\end{paramtable}
IPMResult 字段（converged、三项残差及初始值、factorization/linear\_solve/accepted/
rejected/backend\_escalations 计数、\fld{linear\_solver} 后端名等）见第~\ref{sec:options}~章。

\subsection{验证与校核}
\begin{itemize}
  \item case300\_acdc：accepted$>0$ 且 $\le$iterations、rejected$>0$
        （test\_opf\_solver\_backends.cpp:79--105）；
  \item case2383wp 目标缩放回归：修复前 stationarity$\approx0.45$ 顶格，现收敛（:481--494）；
  \item case2869pegase 稀疏 KKT $<100$ 迭代（:577--599）；case6515rte 保持有限不发散（:601--618）；
  \item dense KKT 后端（\texttt{HACDCPF\_OPF\_LINEAR\_SOLVER=dense}）解 case9/case30（:555--573）；
  \item 大系统诊断与 filter 回溯修复过程记录于 \srcpath{docs/large\_hybrid\_ipm\_diagnostics.md}
        （$\ge512$ 维 KKT 启用残差 filter 回溯，0.5 回溯最多 14 次）。
\end{itemize}

% 直流潮流近似最优潮流（DC OPF）
% 本章文件由 \input 引入，不含导言区。
\section{直流潮流近似最优潮流（DC OPF）}\label{sec:dcopf}

DC OPF 是平台 OPF 家族中的线性化快速通道：把交流网络按直流潮流近似
（无损、无电压幅值变量、无无功方程）化为线性规划（LP）或凸二次规划（QP），
用于经济调度、节点边际电价（LMP）估计、可靠性评估中的大规模故障状态筛选等
对速度敏感、对电压/无功精度不敏感的场景。求解器本体为
\srcpath{src/optimal\_power\_flow/dc\_opf.cpp}（1394 行），公共入口声明于
\srcpath{include/hacdcpf/optimal\_power\_flow/dc\_opf\_solver.hpp}。与 AC OPF
（第\ref{sec:acopfnative}章、第\ref{sec:parityipm}章）不同，本求解器只覆盖
交流侧有功--相角子问题：直流网络、VSC/DC-DC/LCC 换流器物理一律不建模
（详见本章\ref{subsec:dcopf-dcgrid}节）。

\subsection{功能定位与公共接口}

\compmeta{solve\_dc\_opf / check\_dc\_opf\_feasibility}{\srcpath{include/hacdcpf/optimal\_power\_flow/dc\_opf\_solver.hpp}}

\subsubsection*{入口函数}
公共接口仅两个函数
（\srcpath{include/hacdcpf/optimal\_power\_flow/dc\_opf\_solver.hpp}:17--25）：
\begin{align*}
&\text{\srcpath{solve\_dc\_opf}(\fld{sys}, \fld{opts})} \;\longrightarrow\;
  \text{\fld{DCOPFResult}}\\
&\text{\srcpath{check\_dc\_opf\_feasibility}(\fld{sys}, \fld{result}, \fld{tol})}
  \;\longrightarrow\; (\text{feasible},\ \text{max\_violation},\ \text{message})
\end{align*}
前者组装并求解 DC OPF；后者对一份已求解结果做独立的可行性复查
（发电机界、支路界、含松弛节点在内的逐节点功率平衡），容差 \fld{tol} 默认
$10^{-6}$，判据按 $\fld{tol}\times S_{\mathrm{base}}$ 折算为 MW
（\srcpath{src/optimal\_power\_flow/dc\_opf.cpp}:1238--1392，函数
\srcpath{check\_dc\_opf\_feasibility}；判据在 :1390--1391）。公共门面
\srcpath{src/api/hacdcpf.cpp}:1892--1894 仅是透传封装，不附加语义。

\subsubsection*{内部装配结构}
建模产物集中在匿名命名空间的结构体 \fld{DCOPFFormulation}
（\srcpath{src/optimal\_power\_flow/dc\_opf.cpp}:66--115）：成员 \fld{lp} 与
\fld{qp} 分别承载线性化成本 LP 模型与真二次成本 QP 模型（约束结构共用），
\fld{gen\_map}/\fld{branch\_map} 记录活跃发电机/支路到原始组件向量的映射，
\fld{bus\_map} 记录母线稳定编号（\fld{ACBus::index}）到 0-based 位置的映射
（构建于 :35--49，函数 \srcpath{build\_bus\_map}；发现重复母线编号时抛出
\fld{std::invalid\_argument}，:40--46）。基准功率取
$\fld{base\_mva}=\max(\fld{sys.ac.base\_mva},\,1.0)$，默认 100~MVA
（:114、:128）。全章未注明处功率变量均为 $S_{\mathrm{base}}$ 标幺值。

\subsection{DC 线性化假设}\label{subsec:dcopf-assumptions}

代码中的近似假设以注释与结果声明为准，逐条如下。

\begin{enumerate}
  \item \textbf{忽略电阻、以电纳 $b=1/x$ 传输}。建模注释原文：
        ``In DC power flow: $P_{f,ij}=(\theta_i-\theta_j)/x_{ij}$ (ignore r,
        assume $x \gg r$)''（\srcpath{src/optimal\_power\_flow/dc\_opf.cpp}:330，
        函数 \srcpath{build\_dc\_opf\_lp} 的注释块 :326--335）。支路模型只读取
        \fld{Branch::x\_pu}，不读取电阻、分接头变比、移相角与充电电纳
        （:407--409、:438--440）。
  \item \textbf{无电压幅值变量}。决策向量 $x=[\theta;\ P_g;\ P_f;\ \Delta P_d;\
        \lambda]$ 中只有相角 $\theta$（变量布局注释 :73--81），结果结构
        \fld{DCOPFResult} 也只有 \fld{va} 没有 \fld{vm}
        （\srcpath{include/hacdcpf/optimal\_power\_flow/opf\_result.hpp}:186--190），
        即经典 DC 近似中 $V_m\equiv 1$~pu 的内涵在代码里体现为“根本不存在
        电压幅值未知量”。
  \item \textbf{无无功方程}。负荷侧仅聚合有功 \fld{pd\_mw}/\fld{p\_mw}
        （:264--297），结果中没有任何无功字段；\fld{Generator::cost\_*} 之外的
        \fld{qmin\_mvar}/\fld{qmax\_mvar} 不进入模型。
  \item \textbf{无损}。结果对象固定写入模型局限声明：``DC OPF uses a lossless
        linearized AC network and does not model AC voltage magnitude, reactive
        power, or converter physics.''（:683--684，函数
        \srcpath{extract\_dc\_opf\_result}）。
  \item \textbf{小角度假设未显式检验}。非松弛节点相角仅受宽界
        $-\pi\le\theta_i\le\pi$ 约束（:210--217），代码不做
        $|\theta_i-\theta_j|$ 越限告警；小角度是 DC 近似的理论前提而非运行期
        校验项，角度张角大的工况误差由使用者自担。
\end{enumerate}

\subsection{数学模型}\label{subsec:dcopf-model}

\subsubsection*{决策变量与布局}
设活跃母线 $n_b$、活跃发电机 $n_g$、活跃支路 $n_l$（均按 \fld{in\_service}
过滤后计数，:139--155），决策向量分段布局为
（:73--81 注释，偏移量赋值 :157--163）
\begin{align*}
x=[\ \underbrace{\theta}_{n_b}\ ;\ \underbrace{P_g}_{n_g}\ ;\
\underbrace{P_f}_{n_l\ \text{（可选）}}\ ;\
\underbrace{\Delta P_d}_{n_b\ \text{（可选）}}\ ;\
\underbrace{\lambda}_{\text{PWL 断点}}\ ]
\end{align*}
（\srcpath{src/optimal\_power\_flow/dc\_opf.cpp}:157--178，函数
\srcpath{build\_dc\_opf\_lp}）。其中 $P_f$ 变量仅当
\fld{include\_branch\_limits}=\texttt{true} 时引入（:158--159），$\Delta P_d$
削减变量仅当 \fld{load\_shedding}=\texttt{true} 时每母线一个（:160），
$\lambda$ 为凸二次成本发电机（$\fld{cost\_c2}>10^{-12}$ 且
$\fld{pmax\_mw}>\fld{pmin\_mw}$，:170）的分段线性凸组合权重，段数取
$\mathrm{clamp}(\fld{pwl\_segments},1,1000)$（:166）。

\subsubsection*{目标函数}
目标为发电成本最小化，存在 LP 与 QP 两种口径，由后端分派决定实际生效者
（见\ref{subsec:dcopf-solve}节）。

\textbf{（a）LP 口径。} 无线性化段的发电机按边际成本计价；成本数据全零时
取默认边际成本 1（成本单位/MWh），避免目标恒零
（\srcpath{src/optimal\_power\_flow/dc\_opf.cpp}:189--198）。由于 $P_g$ 以
标幺值进入 LP，成本系数乘以 $S_{\mathrm{base}}$ 使报告目标值与等式行对偶
保持 MW 工程单位（注释 :199--202）：
\begin{align*}
\min\ \sum_{k\in\mathcal{G}_{\mathrm{lin}}} c_{1,k}\,S_{\mathrm{base}}\,P_{g,k}
\;+\;\sum_{i} \mathrm{VOLL}\cdot S_{\mathrm{base}}\,\Delta P_{d,i}
\;+\;\sum_{k\in\mathcal{G}_{\mathrm{pwl}}}\sum_{p}
f_k(\hat P_{k,p})\,\lambda_{k,p}
\end{align*}
（:186--202、:321--322、:243--246）。凸二次成本的发电机以 $\lambda$ 凸组合
插值逼近：断点
\begin{align*}
\hat P_{k,p} &= \fld{pmin\_mw}_k + \alpha_p\bigl(\fld{pmax\_mw}_k-\fld{pmin\_mw}_k\bigr),
& \alpha_p &= \frac{p}{N_{\mathrm{seg}}},\quad p=0,\dots,N_{\mathrm{seg}}\\
f_k(\hat P_{k,p}) &= c_{2,k}\hat P_{k,p}^{2}+c_{1,k}\hat P_{k,p}+c_{0,k}
\end{align*}
（:238--246、:483--490），$\lambda_{k,p}\in[0,1]$（:243）。注释说明：凸性加
最小化使下包络自动选取相邻断点，无需二进制或 SOS2 变量（:231--232）。

\textbf{（b）QP 口径。} 真二次成本写成标准形
$\min \tfrac12 x^{\top}Qx+c^{\top}x$（\srcpath{src/optimal\_power\_flow/dc\_opf.cpp}:514--566，
函数 \srcpath{build\_dc\_opf\_qp}）：
\begin{align*}
Q_{P_{g,k},P_{g,k}} &= 2\,c_{2,k}\,S_{\mathrm{base}}^{2},
& c_{P_{g,k}} &= c_{1,k}\,S_{\mathrm{base}}
\end{align*}
（标幺换算推导见注释 :529--534；赋值 :543--556）。$c_2$、$c_1$ 同时为零时
同样回退默认边际成本 1（:552--554）；削减惩罚从 LP 模型原样拷贝（:558--561）。
常数项 $c_0$ 不进 QP 矩阵，仅在报告目标值时补加
（:571--587，函数 \srcpath{compute\_qp\_objective}；$c_0$ 累加在 :581--584）。

\subsubsection*{节点功率平衡}
所有母线（含松弛节点）各有一行平衡方程
（\srcpath{src/optimal\_power\_flow/dc\_opf.cpp}:347--359）。显式 $P_f$ 口径下
\begin{align*}
\sum_{g\in\mathcal{G}(i)} P_{g,g} + \Delta P_{d,i}
- \sum_{l:\,\mathrm{from}(l)=i} P_{f,l} + \sum_{l:\,\mathrm{to}(l)=i} P_{f,l}
= P_{d,i}
\end{align*}
（发电机系数 $+1$：:361--369；削减变量系数 $+1$ 的符号推导注释 :371--376；
$P_f$ 在 from 侧取 $-1$、to 侧取 $+1$：:378--394；右端项 $P_d$：:355--359）。
节点净负荷 $P_{d,i}$ 聚合两类需求并扣除固定注入：母线自带 \fld{pd\_mw}
（:264--267）加 \fld{Load::p\_mw}$\times$\fld{scaling}（:268--273），减去
\fld{StaticGenerator}（:274--279）、\fld{RenewableGen}（:280--285）、
\fld{PVSystem}（:286--291）与 \fld{Storage::p\_mw}（:292--297）的固定出力，
全部折算标幺。

无显式 $P_f$ 变量时（\fld{include\_branch\_limits}=\texttt{false}），支路潮流
直接代入平衡行，形成经典 $B\theta$ 形式（:417--469）：
\begin{align*}
\sum_{j} B_{ij}\,\theta_j &= P_{g,i}+\Delta P_{d,i}-P_{d,i},
& B_{ii} &= \textstyle\sum_{j} b_{ij},\quad B_{ij}=-b_{ij},\quad b_{ij}=1/x_{ij}
\end{align*}
（$B$ 矩阵装配 :424--446；写入等式行 :454--468）。该路径曾漏掉整个松弛节点
平衡行，现仅跳过 $\theta$ 列中松弛节点的零值项而保留其平衡行（BUG-3 修复
注释 :448--453；对角元 :457--460；非对角元 :462--468）。

\subsubsection*{支路功率定义与输电界}
每条活跃支路一行潮流定义方程（\srcpath{src/optimal\_power\_flow/dc\_opf.cpp}:396--416）：
\begin{align*}
P_{f,k} - b_k\,(\theta_{\mathrm{from}(k)}-\theta_{\mathrm{to}(k)}) = 0,
\qquad b_k = 1/x_k
\end{align*}
其中 $x_k$ 取 \fld{Branch::x\_pu}，$|x_k|<10^{-12}$ 时兜底为 $10^{-6}$ 防除零
（:407--409；$B\theta$ 路径同兜底 :438--440；结果回算同兜底 :739--740）。
输电能力以变量盒界实现而非显式不等式（:499--504 注释）：
\begin{align*}
|P_{f,k}| \;\le\; \fld{rate\_a\_mva}_k\cdot\fld{branch\_limit\_margin}\,/\,S_{\mathrm{base}}
\end{align*}
（:249--259；\fld{rate\_a\_mva}$<10^{-9}$ 视为不限，界放宽至
\fld{kHugeBound}$=10^{20}$，:255、:30）。注意仅长期额定 \fld{rate\_a\_mva}
生效，\fld{rate\_b}/\fld{rate\_c} 不进入模型。

\subsubsection*{发电机出力界与相角参考}
发电机界直接折算标幺盒界（:219--229）：
\begin{align*}
\fld{pmin\_mw}_k/S_{\mathrm{base}} \;\le\; P_{g,k} \;\le\;
\fld{pmax\_mw}_k/S_{\mathrm{base}}
\end{align*}
并以 $\pm 10^{20}$ 截断防无界（:225--227）。相角参考节点由
\srcpath{find\_slack\_bus} 选取：第一台 \fld{BusType::SLACK} 型母线，没有则
退化为第 0 号母线（:54--61）；其相角以盒界 $[0,0]$ 钉死，其余节点自由于
$[-\pi,\pi]$（:210--217）。$B\theta$ 路径中松弛列因 $\theta_{\mathrm{slack}}
=0$ 而省略（:448--468）。

\subsubsection*{负荷削减变量与 VOLL}
开启 \fld{load\_shedding} 时每母线一个削减变量
（\srcpath{src/optimal\_power\_flow/dc\_opf.cpp}:299--324）：
\begin{align*}
0 \;\le\; \Delta P_{d,i} \;\le\; \max\!\bigl(P_{d,i},\ 0.01/S_{\mathrm{base}}\bigr)
\end{align*}
（:318--320；只能削减正需求）。失负荷价值 \fld{voll}$\le 0$ 时自动取
（:305--316）
\begin{align*}
\mathrm{VOLL}_{\mathrm{eff}} = \max\Bigl(100,\;
1.1\max_{k}\bigl(c_{1,k}+2c_{2,k}\,\fld{pmax\_mw}_k\bigr)\Bigr)
\end{align*}
即略高于最贵发电机的边际成本上限，保证“能发不削”的字典序。

\subsubsection*{外部电网升格与固定注入}
在求解入口，每台在役 \fld{ExternalGrid} 被升格为一台挂在原母线上的松弛发电机，
界 $\pm 10^{5}$~MW/MVAr，成本系数照搬 \fld{ExternalGrid::cost\_c2/c1/c0}
（\srcpath{src/optimal\_power\_flow/dc\_opf.cpp}:836--855，函数
\srcpath{solve\_dc\_opf}）；求解后其出力拆回
\fld{external\_grid\_p\_mw}，并从 \fld{pg\_mw} 中裁掉（:856--866，lambda
\srcpath{separate\_external\_grid\_dispatch}）。

\subsection{求解流程与后端分派}\label{subsec:dcopf-solve}

\subsubsection*{求解主流程}
\srcpath{solve\_dc\_opf}（\srcpath{src/optimal\_power\_flow/dc\_opf.cpp}:808--1233）
的总体流程如图~\ref{fig:dcopf-flow}：先做 LCC 拒绝与规范投影，再做图论孤岛
预检，然后一次性装配 LP 约束结构并派生 QP 模型，最后按后端枚举分派求解并
回收对偶。

\begin{figure}[H]
\centering
\begin{tikzpicture}[
  node distance=0.55cm,
  box/.style={draw, rectangle, align=center, inner sep=3.5pt, font=\small},
  dec/.style={draw, diamond, aspect=2.1, align=center, inner sep=1.5pt, font=\small},
  arr/.style={-{Stealth[length=2.2mm]}, thick}]
\node[box] (in) {输入 \fld{HybridPowerSystem}};
\node[dec, below=of in] (lcc) {含在役 LCC？};
\node[box, right=1.6cm of lcc] (rej) {拒绝求解\\ \fld{dc-opf:rejected-lcc}};
\node[box, below=of lcc] (proj) {规范投影（开关合并母线）\\ + 外网升格为松弛发电机};
\node[dec, below=of proj] (isl) {孤岛预检\\ 存在无松弛孤岛？};
\node[box, right=1.6cm of isl] (shed) {预削减孤岛负荷\\ 无有效孤岛则报不可行};
\node[box, below=of isl] (build) {装配 LP 约束结构\\ 派生 QP 成本模型};
\node[box, below=of build] (disp) {后端分派（回退链）\\ Gurobi / LCQP / HiGHS / 对偶单纯形};
\node[box, below=of disp] (dual) {支撑 LP 回收对偶（LMP）\\ \fld{compute\_lmp} 门控};
\node[box, below=of dual] (out) {结果归因：反投影回原始母线/支路\\ 填写 \fld{model\_limitations} 等诚实字段};
\draw[arr] (in) -- (lcc);
\draw[arr] (lcc) -- node[above, font=\small]{是} (rej);
\draw[arr] (lcc) -- node[right, font=\small]{否} (proj);
\draw[arr] (proj) -- (isl);
\draw[arr] (isl) -- node[above, font=\small]{是} (shed);
\draw[arr] (isl) -- node[right, font=\small]{否/剪枝后} (build);
\draw[arr] (shed.south) |- (build.east);
\draw[arr] (build) -- (disp);
\draw[arr] (disp) -- (dual);
\draw[arr] (dual) -- (out);
\end{tikzpicture}
\caption{DC OPF 求解主流程（对应
\srcpath{src/optimal\_power\_flow/dc\_opf.cpp}:808--1233）}
\label{fig:dcopf-flow}
\end{figure}

\textbf{LCC 拒绝}：任何在役 \fld{LCCConverter} 都使求解直接返回，状态写明
“this AC-only linear formulation does not model LCC AC/DC coupling”，并提示改用
ParityIPM 或 Ipopt 的 AC OPF（:813--825；见第\ref{sec:parityipm}章）。

\textbf{规范投影}：与 AC OPF/潮流共用同一份 canonical 拓扑——闭合开关与
断路器合并母线（\fld{strip\_dead\_islands}=\texttt{false}，:826--832）。求解后
\fld{va}/\fld{lmp} 按密集量、\fld{load\_shedding\_mw} 按广延量反投影回原始
母线，\fld{pf\_mw} 按 \fld{branch\_orig\_to\_proj} 映射回原始支路（:867--897）。
失败解不回灌投影，以免用尺寸异常掩盖原始不可行性（注释 :869--871）。

\textbf{孤岛预检}：用图模块 \srcpath{build\_power\_system\_graph} 与
\srcpath{analyze\_topology} 在昂贵的 LP 装配前发现 \fld{IsolatedLoad}/
\fld{NoSlack} 孤岛（:921--932）。死岛负荷（母线 \fld{pd\_mw} 与
\fld{ac.loads} 两源相加，口径与建模一致，P1a 注释 :933--935）计入直接削减
（:936--949）；全部孤岛均无效时直接返回
\texttt{"Infeasible: no island with slack bus"}（:950--963）；否则把死岛上的
负荷、发电机、静态发电机、新能源、光伏、储能与相连支路在本地副本中置为
退出再建模（:964--995）。

\textbf{装配}：\srcpath{build\_dc\_opf\_lp} 构建约束结构（:1002--1012 调用；
重复母线编号等拓扑错误在此捕获并转状态，:1004--1012）；空系统或无发电机
直接返回 \texttt{"Empty system or no generators"}（:1014--1019）；随后
\srcpath{build\_dc\_opf\_qp} 派生真二次成本模型（:1021--1022）。

\subsubsection*{后端选择与回退链}
后端枚举 \fld{DCOPFSolverBackend} 取值为 \fld{Auto}/\fld{Native}/
\fld{NativeQP}/\fld{HiGHS}/\fld{Gurobi}
（\srcpath{include/hacdcpf/optimal\_power\_flow/opf\_options.hpp}:128--134）。
各候选通道（均定义于 \srcpath{src/optimal\_power\_flow/dc\_opf.cpp}）：
\begin{itemize}
  \item \textbf{Gurobi QP}（:1064--1074，lambda \srcpath{try\_gurobi\_qp}）：
        经 engine 适配器 \fld{GurobiAdapter} 求解 QP；仅在
        \fld{gurobi.available()} 探测到安装/许可证时尝试，否则记
        \fld{unavailable}；
  \item \textbf{NativeLCQP}（:1047--1061，lambda \srcpath{try\_native\_qp}）：
        自研线性约束 QP 内点求解器（MIPSolvers engine 的
        \srcpath{lcqp\_solver}，经转发头
        \srcpath{include/hacdcpf/engine/kernel/ipm/lcqp\_solver.hpp} 引入），
        求解真二次成本 QP；
  \item \textbf{HiGHS 通道}（:1082--1111，lambda \srcpath{try\_highs}）：
        HiGHS 本身无原生 QP 接口，先经 \fld{SolverEngine} 以
        \fld{preferred\_solver}=\fld{"HiGHS"} 路由一次 QP（实际求解者可能是
        Gurobi 或 LCQP，\fld{solver\_name} 如实反映，注释 :1076--1081）；QP
        失败后回退 \fld{HighsAdapter::solve\_lp} 求 LP，此时目标与对偶只反映
        线性成本模型，求解器名记为 \fld{"HiGHS-LP(QP-fallback)"}
        （:1100--1108）；
  \item \textbf{NativeDualSimplex}（:1118--1137，lambda
        \srcpath{try\_native\_simplex}）：自研对偶单纯形
        \srcpath{solve\_lp\_with\_basis}（同样来自 MIPSolvers engine，经转发头
        \srcpath{include/hacdcpf/engine/kernel/lp\_kernel/dual\_simplex.hpp}），
        求解线性化成本 LP。
\end{itemize}
回退顺序（:1140--1190）：\fld{Auto} 依次试 Gurobi-QP → NativeLCQP → HiGHS
通道 → NativeDualSimplex（:1176--1189）；\fld{Gurobi} 同序但从 Gurobi 起
（:1149--1160）；\fld{NativeQP} 失败仅回退 NativeDualSimplex（:1141--1147）；
\fld{HiGHS} 不可用时直接失败返回（:1162--1170）；\fld{Native} 只用对偶单纯形
（:1172--1174）。每一步尝试都记入 \fld{solver\_chain}，格式
\fld{<backend>:<ok|fail(detail)>}，供事后审计真实回退路径（:1034--1040、
:1200--1205）。

\textbf{目标语义消歧}：LP 回退成功时必须把 \fld{use\_qp} 复位，否则会用 QP
口径重算 LP 可行点，造成目标/对偶语义不一致（注释 :1042--1046、
:1113--1117）。最终按收敛路径重算或保留目标值并填写
\fld{objective\_model}：QP 收敛为 \fld{"QP"}（真二次成本含 $c_0$，由
\srcpath{compute\_qp\_objective} 重算，:1216--1219），否则为 \fld{"LP-PWL"}
或 \fld{"LP"}，并以 \fld{pwl\_segments\_effective} 记录实际分段数（QP 时置 0，
:1220--1223）。

\subsubsection*{对偶恢复与支撑 LP}
QP/LP 主解的对偶若不完整（\srcpath{has\_complete\_duals} 校验行对偶与盒对偶
数量，:589--595），且 \fld{compute\_lmp}=\texttt{true} 时，追加一次支撑 LP：
以原 LP 模型（QP 情形把成本换成梯度 $Qx+c$）用对偶单纯形重解，取回全部行
对偶与盒对偶（:635--665，函数
\srcpath{populate\_missing\_duals\_from\_supporting\_lp}；QP 梯度替换 :643--645；
门控 :1195--1198）。代码注释如实声明该通道的已知缺陷：自研单纯形对
有界变量 LP 的盒对偶恢复不正确，\fld{include\_branch\_limits}=\texttt{true}
且支路界真正起作用时 \fld{branch\_mu\_lower}/\fld{branch\_mu\_upper}
\textbf{不可靠}；LMP（节点平衡行对偶）仍有效可用于经济调度
（NOTE 注释块 :621--634）。因此结果中 \fld{branch\_mu\_valid} 在 KKT 级提取
实现之前\textbf{恒为 false}（:680--682、:781--785）。

\subsubsection*{与 HiGHS/AML 建模层的分派关系}
本求解器只经 engine 适配器调用 HiGHS/Gurobi（转发头
\srcpath{include/hacdcpf/engine/solver/external/adapters.hpp}），自身不经过
AML 建模层。平台另有一套基于 MIPSolvers AML 的 DC OPF 实现
\srcpath{power\_models::solve\_dc\_opf}
（\srcpath{src/power\_models/dc\_opf\_builder.cpp}:13），以集合/参数/变量的
代数建模语言重构同一数学模型并交由 AML 后端求解，用于交叉验证；其建模细节
见第\ref{sec:powermodels}章。两套实现相互独立：本章所述的孤岛预检、外网
升格、负荷削减与支撑 LP 对偶恢复只存在于引擎版。

\subsection{直流电网与换流器覆盖情况}\label{subsec:dcopf-dcgrid}

\textbf{不覆盖。} 具体口径：
\begin{itemize}
  \item 直流母线、直流支路、VSC 换流器、DC/DC 变换器、能量路由器均不进入
        模型——建模仅遍历 \fld{sys.ac.*} 容器（
        \srcpath{src/optimal\_power\_flow/dc\_opf.cpp}:125--127），结果结构也
        无任何直流侧字段
        （\srcpath{include/hacdcpf/optimal\_power\_flow/opf\_result.hpp}:186--236）；
  \item 结果固定声明 \fld{model\_scope}=\fld{"dc-opf:linear-no-converter-model"}
        （:678--679），\fld{converter\_model\_scope} 各换流器标志位恒为
        \texttt{false}（注释见
        \srcpath{include/hacdcpf/optimal\_power\_flow/opf\_result.hpp}:232--234）；
  \item 含在役 LCC 的系统被整体拒绝而非近似处理（:813--825）；
  \item 直流侧负荷由此形成已知盲区：可靠性评估调用方对此有显式告警
        （“DC 负荷中断不计入 OPF，EENS/LOLE 可能偏乐观”，
        \srcpath{src/reliability/reliability\_assessment.cpp}:1047--1056）。
\end{itemize}
需要交直流联合优化的场景应使用 Parity IPM 全空间 OPF
（第\ref{sec:parityipm}章）或 AML 层的 acdcopf builder
（第\ref{sec:powermodels}章）。

\subsection{NativeLCQP 的结构化初值}

对每个在役 AC 支路连通分量 $\mathcal C$，先在发电上下界内调整 $P_g$，不足
部分再由允许的 $\Delta P_d$ 补足，使
\begin{equation}
\sum_{i\in\mathcal C}(P_{g,i}+\Delta P_{d,i}-P_{d,i})=0.
\end{equation}
每个分量选择一个参考角并删除其行列，求约化 Laplacian
$B_{\mathcal C,r}\theta_r=p_{\mathcal C,r}$；随后由
$P_{f,ij}=(\theta_i-\theta_j)/x_{ij}$ 回填显式支路流，并补齐 PWL convexity/
interpolation 变量。因此连通分量功率和是 Laplacian 可解性的兼容条件，而非
额外优化问题。生成点经过变量盒界与 $\|A_{eq}x_0-b_{eq}\|_\infty$ 审计后，
只通过 \fld{QPModel::x0} 交给 NativeLCQP；其他后端不消费该初值，构建失败则
精确回到确定性冷初始化。该步骤复杂度为图遍历线性工作加各分量一次约化
Laplacian 稀疏分解，不运行非线性 Phase I。

作为 ACOPF warm-start-only 子问题时，
\fld{phase\_one\_time\_limit\_ms} 从 \srcpath{solve\_dc\_opf} 入口开始计时，
因此 canonical projection、formulation、上述结构初值、symbolic analyze 与
numeric iteration 都计入总预算。NativeLCQP 在迭代/分解边界协作式检查截止；
一个已经开始的稀疏分解不可中断，允许完成后报告
\fld{phase\_one\_budget\_overshoot\_ms}。超时迭代仍须通过显式原坐标残差门槛
才可作为 warm start，且不声明 DCOPF 最优；有限预算耗尽后不再运行 simplex
fallback。

\subsection{选项：DCOPFOptions}

\compmeta{DCOPFOptions}{\srcpath{include/hacdcpf/optimal\_power\_flow/opf\_options.hpp}}

全部字段如下（定义 :139--160）。

\begin{paramtable}
\fld{feasibility\_tol} & $\varepsilon_{\mathrm{f}}$ & — & $10^{-9}$--$10^{-4}$（经验值） & $10^{-6}$ & 原始可行性容差；传入 LCQP（\fld{tol\_primal}）与对偶单纯形\\
\fld{optimality\_tol} & $\varepsilon_{\mathrm{o}}$ & — & $10^{-9}$--$10^{-4}$（经验值） & $10^{-6}$ & 对偶/最优性容差；传入 LCQP（\fld{tol\_dual}/\fld{tol\_gap}）与对偶单纯形\\
\fld{max\_iterations} & — & — & $10^{2}$--$10^{5}$（经验值） & 10000 & 单纯形/LCQP 最大迭代数\\
\fld{verbose} & — & — & true/false & false & spdlog 打印规模与收敛信息（:1024--1028、:1225--1229）\\
\fld{solver} & — & — & Auto/Native/\allowbreak NativeQP/\allowbreak HiGHS/\allowbreak Gurobi & Auto & 后端选择与回退链入口（:1140--1190）\\
\fld{pwl\_segments} & $N_{\mathrm{seg}}$ & — & 1--1000（clamp，:166） & 4 & 凸二次成本的 LP 分段线性插值段数；QP 路径不使用\\
\fld{include\_branch\_limits} & — & — & true/false & true & 是否引入显式 $P_f$ 变量及 \fld{rate\_a\_mva} 盒界；false 时退回纯 $B\theta$ 形式且无支路约束\\
\fld{branch\_limit\_margin} & — & — & 0.9--1.0（经验值） & 1.0 & 支路界缩放系数，乘在 \fld{rate\_a\_mva} 上（:256）\\
\fld{load\_shedding} & — & — & true/false & true & 是否引入每母线削减变量 $\Delta P_d$\\
\fld{voll} & VOLL & —/MWh & $\ge 0$ & 0.0 & 失负荷价值；$\le 0$ 时按 $\max(100,\,1.1\times$最高边际成本$)$ 自动取值（:305--316）\\
\fld{compute\_lmp} & — & — & true/false & true & 主解后是否追加支撑 LP 回收对偶；false 时 \fld{lmp} 清空并注明原因（:1195--1198、:1206--1211）\\
\fld{structural\_warm\_start} & — & — & true/false & true & 为 NativeLCQP 构造逐分量功率配平与约化 Laplacian 初值；LP/外部后端忽略\\
\fld{compact\_quadratic\_model} & — & — & true/false & false & Phase-I QP 省略仅供 LP-PWL 使用的插值变量与行\\
\fld{accept\_phase\_one\_iterate} & — & — & true/false & false & 允许非最优 NativeLCQP 末迭代经显式残差门槛成为 warm-start-only 结果\\
\fld{phase\_one\_iterate\_tolerance} & $\varepsilon_I$ & pu & $\ge0$ & 0.1 & warm-start-only 末迭代的 equality/box 最大残差门槛\\
\fld{phase\_one\_time\_limit\_ms} & $T_I$ & ms & $\le0$ 或 $>0$ & 0 & warm-start-only 端到端协作式墙钟预算；$\le0$ 不限时\\
\end{paramtable}

\subsection{结果字段：DCOPFResult}

\compmeta{DCOPFResult}{\srcpath{include/hacdcpf/optimal\_power\_flow/opf\_result.hpp}}

全部字段如下（定义 :186--236）。向量字段的索引空间：\fld{va}/\fld{lmp}/
\fld{load\_shedding\_mw} 按原始母线向量位置（经反投影）；\fld{pg\_mw} 按
\fld{sys.ac.generators} 位置；\fld{pf\_mw}/\fld{branch\_mu\_*} 按
\fld{sys.ac.branches} 位置；\fld{external\_grid\_p\_mw} 按
\fld{sys.ac.external\_grids} 位置。

\begin{paramtable}
\fld{va} & $\theta_i$ & rad & $(-\pi,\pi)$ & 空 & 母线电压相角（\fld{BusType::SLACK} 处恒 0）\\
\fld{pg\_mw} & $P_g$ & MW & \fld{pmin\_mw}--\fld{pmax\_mw} & 空 & 发电机有功调度（由标幺乘 $S_{\mathrm{base}}$ 还原，:707--713）\\
\fld{external\_grid\_p\_mw} & $P_{\mathrm{ext}}$ & MW & $\pm 10^{5}$ & 空 & 外网（升格发电机）有功（:856--866）\\
\fld{pf\_mw} & $P_f$ & MW & $\pm$\fld{rate\_a\_mva} & 空 & 支路有功，from$\to$to 为正；无显式 $P_f$ 时由 $(\theta_f-\theta_t)/x$ 回算（:715--744）\\
\fld{converged} & — & — & true/false & false & 后端报告成功即 true（:691）\\
\fld{iterations} & — & — & — & 0 & 最终后端迭代数\\
\fld{objective} & $f$ & —/h & — & 0.0 & 目标值；口径见 \fld{objective\_model}（QP 含 $c_0$，LP 为线性/PWL 成本）\\
\fld{status} & — & — & — & 空 & 后端状态串或拒绝/失败原因（如 \texttt{"HiGHS not available"}）\\
\fld{solver\_name} & — & — & — & 空 & 实际收敛后端名（可能含 \fld{(QP-fallback)} 后缀）\\
\fld{runtime\_sec} & — & s & — & 0.0 & 含投影与预检的端到端耗时（:1192--1193）\\
\fld{structural\_warm\_start\_*} & — & pu/计数/布尔 & — & inf/0/false & 请求、构建、实际消费、分量/分解数、原坐标 equality residual 与状态原因\\
\fld{solver\_initial\_primal\_residual} & — & 缩放坐标 & $\ge0$ & inf & NativeLCQP 第 0 次迭代的 primal residual；含严格正 bound-slack 残差\\
\fld{phase\_one\_warm\_start\_only} & — & — & true/false & false & 末迭代只获准供上层 warm start 使用，不是 DCOPF 最优证书\\
\fld{phase\_one\_budget\_exhausted} & — & — & true/false & false & formulation 或 LCQP 阶段达到有限墙钟预算\\
\fld{phase\_one\_budget\_overshoot\_ms} & — & ms & $\ge0$ & 0 & 不可中断在途工作造成的预算超时\\
\fld{native\_qp\_symbolic\_analyze\_calls} & — & 次 & 0--1 & 0 & 本次 NativeLCQP KKT 稀疏模式分析次数；每轮 numeric factorization 复用该模式\\
\fld{lmp} & $\lambda_i$ & —/MWh & $\ge 0$（无罚项时） & 空 & 节点边际电价：平衡行对偶除以 $S_{\mathrm{base}}$（:746--759）\\
\fld{branch\_mu\_lower} & $\mu^{-}$ & —/MWh & — & 空 & 支路下界盒对偶/$S_{\mathrm{base}}$；当前\textbf{不可靠}（:769--779）\\
\fld{branch\_mu\_upper} & $\mu^{+}$ & —/MWh & — & 空 & 支路上界盒对偶/$S_{\mathrm{base}}$；当前\textbf{不可靠}（:769--779）\\
\fld{branch\_mu\_valid} & — & — & true/false & false & 恒 false，直至 KKT 级盒对偶提取实现（:680--682、:785）\\
\fld{branch\_mu\_validity\_reason} & — & — & — & 空 & 上述无效性的人类可读说明\\
\fld{lmp\_valid} & — & — & true/false & false & 行对偶证书完整时为 true（:760--767）\\
\fld{lmp\_validity\_reason} & — & — & — & 空 & LMP 来源/缺失原因说明\\
\fld{load\_shedding\_mw} & $\Delta P_d$ & MW & $\ge 0$ & 空 & 逐母线削减量（含孤岛预削减；$<10^{-6}$ 数值噪声清零，:794）\\
\fld{total\_load\_shedding\_mw} & — & MW & $\ge 0$ & 0.0 & 削减总量\\
\fld{solver\_chain} & — & — & — & 空 & 依序记录各后端尝试结果，供审计回退路径（:1034--1040）\\
\fld{objective\_model} & — & — & QP/LP-PWL/LP & 空 & 报告目标值对应的成本口径（:1213--1223）\\
\fld{pwl\_segments\_effective} & — & — & 0--1000 & 0 & LP-PWL 实际段数；0 表示 QP 真二次成本\\
\fld{converter\_model\_scope} & — & — & — & 全 false & 换流器建模覆盖声明（DC OPF 不建模任何换流器）\\
\fld{model\_limitations} & — & — & — & 空 & 本次求解的模型局限声明（无损线性化、储能单时段等，:683--688）\\
\end{paramtable}

\subsection{典型使用场景}

\begin{itemize}
  \item \textbf{可靠性评估（FMEA/状态评估）}：每个故障状态以开启负荷削减的
        DC OPF 评估最小负荷中断；VOLL 被强制设为远高于发电成本，实现“先削减
        最小化、后成本”的字典序
        （\srcpath{src/reliability/reliability\_assessment.cpp}:1058 调用，
        选项设置 :1223--1228；AC-only 适用范围声明 :782--789）。
  \item \textbf{EV--交通耦合联合优化}：按时段把 EV 充电负荷注入系统后逐时段
        求解 DC OPF，累加发电成本并取 LMP 构成联合社会福利目标
        （\srcpath{src/ev\_power\_traffic/joint\_social\_welfare.cpp}:258--263；
        CTM 变体 \srcpath{src/ev\_power\_traffic/ctm\_joint\_welfare.cpp}:251--255）。
  \item \textbf{SPPT 语义保持验证}：MR3d 用 DC OPF 节点电价校验
        rich$\to$canonical 投影前后的语义不变性
        （\srcpath{src/sppt/metamorphic.cpp}:219--225）。
  \item \textbf{HTTP 后端}：v1 分析接口在 \fld{analysis}=\fld{"optimal\_power\_flow"}
        且 \fld{solver}=\fld{"dc"} 时调用本求解器并序列化结果
        （\srcpath{src/server/runtime\_api\_v1.cpp}:864--877，函数
        \srcpath{execute\_analysis}）。
  \item \textbf{公共门面}：\srcpath{src/api/hacdcpf.cpp}:1892--1894 透传封装。
\end{itemize}
如实说明两处“不调用”：\textbf{time\_series} 多时段/年度流水线挂接的是
AC OPF（\srcpath{src/time\_series/time\_series\_pf.cpp}:24 引入
\srcpath{ac\_opf\_solver.hpp}），不使用本求解器；\textbf{market} 模块的
SCUC/SCED 为其自研 LP 体系（\srcpath{src/market/market\_simulation.cpp}:2455
附近），亦不复用 \srcpath{solve\_dc\_opf}。

\subsection{验证与校核}

\begin{itemize}
  \item \textbf{MATPOWER 标准算例}：case9 收敛且目标有限、\fld{pg\_mw} 尺寸
        正确（\srcpath{tests/test\_matpower\_cases.cpp}:618--624）；case14 调度
        逐机检界（同文件 :626--632 起）。
  \item \textbf{AC/DC/PF 三方交叉验证}：小算例集上 DC OPF 收敛后以 AC 潮流
        回代其调度检验可行性；case57 上“DC 调度$\to$AC 潮流”不保证收敛，
        仅告警不断言（\srcpath{tests/test\_acopf\_dcopf\_crossval.cpp}:250--303，
        豁免逻辑 :273--281）。
  \item \textbf{可行性自检}：\srcpath{check\_dc\_opf\_feasibility} 对 case14
        解复检通过（同文件 :308--346）。
  \item \textbf{孤岛回归}：孤立负荷岛逐母线削减记账（:348--392）、死岛发电
        机/支路剪枝（:394--444）、松弛节点平衡行纳入复检（:446--485）。
  \item \textbf{调度质量}：case30 上与 AC OPF（NativeAC）对比，$\sum|\Delta
        P_g|/\sum P_g$ 容忍阈值为 1.0——注释明言“DC OPF is a crude
        approximation”（同文件 :490--548，阈值 :536--547）。
  \item \textbf{LMP 口径}：无支路界时 LMP 非负且不超过最高边际成本的数量级
        （防 $S_{\mathrm{base}}$ 换算 bug，同文件 :703--741）；无拥塞时全网
        LMP 一致（:748--774）。
\end{itemize}
验证体系的总体组织见第\ref{sec:verification}章。

\subsection{边界与局限}

\begin{enumerate}
  \item \textbf{DC 近似的固有盲区}：无损、无电压幅值、无无功（结果
        \fld{model\_limitations} 首条，\srcpath{src/optimal\_power\_flow/dc\_opf.cpp}:683--684）。
        电压越限、无功不足、网损引起的附加阻塞均不可见；重载、大角度张角或
        低 $x/r$ 比配电网场景下误差显著，代码不做任何运行期误差估计。
  \item \textbf{支路模型简化}：仅取 \fld{x\_pu}，忽略电阻、分接头、移相角与
        充电电纳（:407--409）；限值只用 \fld{rate\_a\_mva}（:254）；零电抗支路
        以 $x=10^{-6}$ 兜底而非合并（:407--409、:438--440），可能引入数值
        病态。
  \item \textbf{支路拥塞对偶不可用}：\fld{branch\_mu\_valid} 恒 false
        （:785）；需要拥塞租金分解的研究只能使用 LMP 总量，不能拆分
        能量/拥塞分量（NOTE 注释 :621--634）。
  \item \textbf{LMP 的适用范围}：LMP 取自支撑 LP 平衡行对偶，标幺换算已按
        $1/S_{\mathrm{base}}$ 还原（:746--759）；\fld{include\_branch\_limits}
        =\texttt{true} 时 LMP 数值仍有效，但注意测试约定 LMP 专项用例是在
        \fld{include\_branch\_limits}=\texttt{false} 下跑的
        （\srcpath{tests/test\_acopf\_dcopf\_crossval.cpp}:700--701 注释）。
  \item \textbf{无直流网络与换流器}：见本章\ref{subsec:dcopf-dcgrid}节；
        直流负荷对结果完全不可见。
  \item \textbf{储能为固定注入}：\fld{Storage::p\_mw} 作为不可调注入扣除，
        无跨时段 SOC 动态（:685--688 如实声明）。
  \item \textbf{外网升格为宽界松弛发电机}：$\pm 10^{5}$~MW（:846--847），
        多外网并存时各外网出力由成本排序决定，不模拟实际互联调度协议。
  \item \textbf{后端可用性依赖环境}：Gurobi 需本机安装与许可证
        （:1066--1073）；\fld{HiGHS} 后端在库不可用时直接失败
        （:1162--1170）；QP 经 \fld{SolverEngine} 路由时实际求解者以
        \fld{solver\_name}/\fld{solver\_chain} 为准（:1076--1081）。
  \item \textbf{LP-PWL 插值误差}：凸二次成本走 LP 路径时为分段线性近似，
        精度由 \fld{pwl\_segments}（默认 4）控制；报告口径以
        \fld{objective\_model}/\fld{pwl\_segments\_effective} 区分
        （:1220--1223）。
  \item \textbf{相角界为宽界}：$[-\pi,\pi]$ 仅防无界（:210--217），不构成
        静态安全（角差）约束；需要角差约束的研究应改用
        AC OPF 或 AML 建模层（第\ref{sec:powermodels}章）。
\end{enumerate}

% 第 6 章：无功优化（RPO）
\section{无功优化（RPO）}\label{sec:rpo}

\subsection{功能定位与入口}
\compmeta{opf::solve\_rpo}{\srcpath{include/hacdcpf/optimal\_power\_flow/reactive\_power\_opt.hpp}}
实现 src/optimal\_power\_flow/reactive\_power\_opt.cpp:621--1386。
MINLP 结构：连续变量（$V_m$、$Q_g$，嵌在 parity OPF 内）+ 整数变量
（OLTC 档位 $t\in[t_{min},t_{max}]$、可切并联 $s\in[0,n_{steps}]$）
（reactive\_power\_opt.hpp:293--304）。配套入口
\srcpath{inspect\_rpo\_controls}（离散控制清单含排除原因）与
\srcpath{apply\_rpo\_discrete\_solution}（动作回写）。

\subsection{数学模型}
档位-变比：$\mathrm{ratio}=1+(t-t_{neutral})\cdot\fld{step\_pct}/100$
（reactive\_power\_opt.cpp:80--82），经 \fld{source\_branch\_idx} 联动更新 AC 支路 tap
（:84--106）。目标 \fld{RPOObjective}（hpp:13--17）：
MinVoltageDeviation $\sum(V-V^*)^2$ / MinActiveLoss / Combined，
由内层 OPF 解出后从潮流账重算（\srcpath{rpo\_objective} :269--285；
\srcpath{compute\_loss} 全源-全荷台账 :189--266）。

\paragraph{内层 NLP 目标解耦（RPOInnerObjective，hpp:21--34）}
\begin{itemize}
  \item \texttt{MatchRPO}：内层直接优化 RPO 目标，机组无功参与调压（中小电网默认）；
  \item \texttt{LossEconomic}：单位边际成本经济调度（min 总发电 $\equiv$ min 网损，
        理论论证 :687--704），大电网稳健；MinActiveLoss 恒用 LossEconomic；
        $\ge100$ 母线先做经济种子，$\ge$\fld{robust\_inner\_min\_buses}=1000 直接
        LossEconomic 起步（:720--724）。
\end{itemize}
OLTC 准入条件：投运 + \fld{tap\_max}$>$\fld{tap\_min} + \fld{tap\_step\_percent}$>0$
+ 当前档位在铭牌内；MATPOWER 固定变比不伪造为 OLTC；默认档位窗口 $\pm2$
（GUI 显式默认 \fld{max\_tap\_move}=2，run\_gui\_server.cpp:18370/18426）。

\paragraph{网损台账 \srcpath{compute\_loss}（:189--266）}
源侧 = pg + external\_grid + pren + pstor + 固定 DER（static gen、不可削减可再生、
不可控 PV 功率曲线、DC SG、DC DCSG(\fld{p\_set})、DC PV 阵列、移动储能、VPP）
$-$ microgrid 交换；\textbf{VSC/DC-DC/能量路由器为内部转移不计入}；
荷侧 = bus.pd + loads + 不对称负荷（pa+pb+pc）+ pflex（结果空则
\fld{flexible\_loads.p\_mw}）+ 充电站（\fld{p\_total\_kw}$>$0 用之，否则
max$\times$utilization$\times$simultaneity/1000）+ 电机（$s_n\cos\varphi/\eta$）
+ DC bus.pd + DC loads $-$ $\max(0,\text{切负荷})$。

\paragraph{并联设备建模}
排除原因全集：\texttt{not\_in\_service} / \texttt{not\_switchable} /
\texttt{insufficient\_step\_range}（$n_{steps}\le1$）/ \texttt{zero\_or\_invalid\_step\_size}
/ \texttt{current\_step\_out\_of\_range}（:118--127）；
动作回写 \fld{current\_step}$=s$、\fld{bs\_mvar}$=$\fld{bs\_per\_step}$\cdot s$（:302--307）；
ShuntResult 报告 \fld{bs\_mvar\_before/after}（:1362--1371）。

\paragraph{控制清单 \srcpath{inspect\_rpo\_controls}（:512--601）}
\fld{selected\_for\_optimization} = adjustable $\land$
（\fld{restrict\_tap\_indices}$\to$在 enabled 列表）$\land$ \fld{max\_tap\_move}$\ne0$；
优化窗口 \fld{optimization\_tap\_min/max} $=\mathrm{clamp}(\fld{tap\_pos}\mp\fld{max\_tap\_move},\ \text{铭牌})$；
\fld{electrical\_tap\_base} $=$ \fld{branch.tap}/\fld{ratio\_current}，
\fld{electrical\_tap\_min/max} $=$ base$\times$\fld{ratio\_min/max}（:560--575）；
字段全集（hv/lv\_bus、tap\_side、position\_count 等）见 reactive\_power\_opt.hpp:181--208。

\subsection{算法流程（离散坐标搜索）}
文件头注释 :1--20；主流程 :621--1358：
\begin{enumerate}
  \item \textbf{Phase 0 灵敏度排序}：每变量 $\pm1$ 扰动按 $|\Delta obj|$ 降序（:1006--1083）；
  \item \textbf{Phase 1 离散三分搜索}（近似单峰，$O(\log R)$），保持 Gauss--Seidel 串行（:1085--1159）；
  \item \textbf{Phase 2 坐标下降} $\le3$ 轮（:1162--1224）；
  \item \textbf{Phase 3 成对 $\pm1$ 邻域抛光}（:1227--1284）。
\end{enumerate}
每个候选 = 一次完整 AC OPF；解缓存 \fld{eval\_cache}、就地修改避免深拷贝；
并行 Jacobi 波次（\srcpath{evaluate\_batch\_parallel}，:405--486）——
动态 work-stealing（共享 atomic 索引）、每线程一次深拷贝复用、最优候选在互斥锁下
连结果一起留存、incumbent 免二次求解推进、不可行候选记 $10^{30}$；
混合 AC/DC 强制单线程（MUMPS 进程级互斥，:937--941）；
接受准则：内层 converged 或实测容差（viol $\le$ \fld{ipm\_tol} 且 stat $\le$
\fld{stationarity\_tol}）（:376--380）；纯 AC 用 \fld{search\_ipm\_iter} 快速失败，
混合保留 Ipopt 第二腿（:904--911）。局部敏感度面永不用于剪枝（:1164--1166）。

\paragraph{搜索细节}
Phase 1 三分搜索整数网格更新：$f_1<f_2\to hi=m_2-1$，否则 $lo=m_1+1$，余窗穷举
（:1104--1159）；Phase 2 循环改善判据 $obj < incumbent - \fld{gap\_tol}$（:1200, 1217）；
无可调离散设备时直接返回 baseline（converged，
``Feasible OPF point (no adjustable discrete devices)''，\fld{nodes\_explored}=1，
:843--865）。

\paragraph{内层求解配置 \srcpath{solve\_opf\_inplace}（:312--403）}
\fld{max\_outer\_iterations}=1；目标映射 Economic（LossEconomic）/
VoltageDeviationAndLoss（Combined）/VoltageDeviation；
第二腿按 $\fld{quality}=\max(viol,stat)$ 择优保留（:392--397）；
经济种子：$\ge100$ 母线触发，ParityIPM、
\fld{max\_inner}$=\mathrm{clamp}(\fld{max\_ipm\_iter},200,800)$、
\fld{stationarity\_tol}$=\min(\fld{rpo\_tol},10^{-6})$、无 fallback；
$\ge$\fld{robust\_inner\_min\_buses} 直接同伦起步、跳过注定失败的直接种子（:737--782）；
\srcpath{apply\_unit\_costs} 一次性改写 $c_1=1/c_2=c_0=0$ 保证暖链同源（:705--711）。
相位计时可由 \texttt{HACDCPF\_RPO\_PROF}=1 开启（:42--53）。

\subsection{诚实口径与结果结构}
RPOResult 定义 \srcpath{reactive\_power\_opt.hpp:245--291}（逐字段表见第~\ref{sec:options}~章）：
\fld{algorithm}=\fld{discrete\_coordinate\_search\_with\_ac\_opf}、
\fld{globally\_certified}=false、\fld{optimality\_gap\_available}=false、\fld{gap}=1.0
（hpp:248--255）；离散动作报告于 \fld{taps}/\fld{shunts}
（\texttt{TapResult}/\texttt{ShuntResult} 向量，hpp:275--276）；
状态文案区分 time-limit/evaluation-limit/local optimum（:1328--1333）；
\fld{bc\_stats} 仅为兼容信封，注释明确“非 B\&B 全局 gap 证书”（hpp:284--286）——
字段映射 \fld{cuts\_added}=灵敏度平面数、\fld{lp\_solves}=\fld{nlp\_solves}（:1374--1383）；
基线容差接受会附加 ``[baseline accepted at measured tolerance: ...]''（:829--834）。
即：只返回可行局部最优点，不报虚假 $\mathrm{gap}=0$，不声称全局 MINLP 最优。

\paragraph{终止状态标志（hpp:256--257）}
两个布尔标志默认 false：\fld{terminated\_by\_time\_limit}
（$=$ \fld{time\_limit\_sec}$>0$ 且耗时 $\ge$ 时限，:1298--1299）与
\fld{terminated\_by\_evaluation\_limit}（$=$ \fld{nlp\_solves} $\ge$ \fld{max\_nodes}，:1300）；
状态串 ``Feasible incumbent at time limit'' / ``Feasible incumbent at evaluation limit'' /
``Feasible local optimum (heuristic search)'' 的判定复用这两个标志
（:1328--1333，语义不变、判定单点化）。

\paragraph{\fld{model\_limitations}（hpp:258）}
string 向量、默认空；``无全局 MINLP 证书''口径现以结构化字段对外暴露，追加点：
\begin{itemize}
  \item 基线 OPF 失败路径：离散坐标/邻域搜索无全局 MINLP 最优性证书（:821--823）；
  \item 无可调离散设备早退路径：结果即连续 OPF 基线而非全局 MINLP 证书（:858--859）；
  \item LossEconomic 内层目标与非 MinActiveLoss 外层目标不匹配：
        早退路径（:860--863）与正常路径（:1311--1314）各追加一条说明；
  \item 正常路径无条件条目（可行局部 incumbent、无全局证书，:1301--1302），
        以及时限截断（:1303--1305）/ 评估上限截断（:1307--1309）各一条。
\end{itemize}

\subsection{参数表（RPOOptions）}
定义 reactive\_power\_opt.hpp:37--155。
\begin{paramtable}
\fld{objective} & -- & 枚举 & 三档 & Min\-Voltage\-Deviation & RPO 目标\\
\fld{inner\_nlp\_objective} & -- & 枚举 & 两档 & Match\-RPO & 内层 NLP 目标解耦\\
\fld{robust\_inner\_on\_failure} & -- & bool & -- & true & 失败时稳健内层重试\\
\fld{robust\_inner\_min\_buses} & -- & 母线 & -- & 1000 & 直接 LossEconomic 阈值\\
\fld{vdev\_weight}/\fld{loss\_weight} & $w_v,w_l$ & -- & $\ge0$ & 1.0/1.0 & Combined 目标权重\\
\fld{v\_target} & $V^*$ & pu & 0.95--1.05 & 1.0 & 电压目标\\
\fld{max\_nodes} & -- & 个 & -- & 50000 & 离散评估上限\\
\fld{time\_limit\_sec} & -- & s & -- & 120 & 时限\\
\fld{gap\_tol} & -- & -- & -- & $10^{-4}$ & 兼容保留\\
\fld{int\_tol} & -- & -- & -- & $10^{-5}$ & 兼容遗留（未用）\\
\fld{num\_threads} & -- & 线程 & 0=auto & 0 & 并行波次线程数\\
\fld{branching}/\fld{node\_sel} & -- & -- & -- & -- & 遗留未用\\
\fld{max\_ipm\_iter} & -- & 次 & -- & 2000 & 内层 IPM 迭代上限\\
\fld{search\_ipm\_iter} & -- & 次 & -- & 200 & 搜索期快速失败迭代上限\\
\fld{ipm\_tol} & -- & -- & -- & $10^{-6}$ & 内层可行性容差\\
\fld{stationarity\_tol} & -- & -- & -- & $10^{-3}$ & 内层 stationarity 容差\\
\fld{inner\_solver\_backend} & -- & 枚举 & -- & ParityIPM & 内层求解器\\
\fld{seed\_homotopy\_on\_failure} & -- & bool & -- & true & 失败时目标同伦种子\\
\fld{max\_tap\_move} & -- & 档 & -- & -1（全程） & 档位窗口；GUI 显式默认 2\\
\fld{restrict\_tap\_indices} + \fld{enabled\_tap\_indices} & -- & -- & -- & false & 限定可调 OLTC 子集\\
\fld{verbose} & -- & bool & -- & false & 透传内层 AC OPF 迭代日志（:364）\\
\fld{enforce\_*}（四项） & -- & bool & -- & true & 内层约束族开关\\
\end{paramtable}

\subsection{验证与校核}
\begin{itemize}
  \item case9 诚实本地证书：\fld{globally\_certified}=false、
        \fld{optimality\_gap\_available}=false、\fld{model\_limitations} 非空（:33）、
        \fld{gap}=1.0、$0<\mathrm{loss}<50$~MW（test\_reactive\_power\_opt.cpp:16--42）；
  \item 并联动作与目标不劣化（:43--81）；case300 seeded 求解 \fld{nlp\_solves}=3、
        \fld{max\_tap\_move}=2（:83--109）；
  \item 控制清单排除原因（\texttt{fixed\_or\_invalid\_tap\_range} 等）与
        \fld{max\_tap\_move}=1 窗口（:111--181）；
  \item MATPOWER 固定变比不被虚构为 OLTC（:183--210）；
        \srcpath{apply\_rpo\_discrete\_solution} 联动支路 tap=$1.03\times1.025$（:212--241）；
  \item 交叉验证文档 \srcpath{docs/reactive\_power\_optimization\_validation.md}：
        独立 OPF 复算 + PF 回放，判据目标相对差 $\le\max(10^{-3},10\epsilon_d)$、
        电压差 $\le\max(5\times10^{-4},10\epsilon_p)$~pu；
        大规模实测（case2869+20 OLTC+10 并联）：基线 74 迭代/25~s 收敛到参考值 133999，
        暖启动链 2--7 迭代、均值 $\approx0.46$~s/次；MinActiveLoss 网损 1561.8$\to$1558.8~MW；
  \item GUI E2E：\srcpath{tests/e2e/rpo\_gui\_e2e.mjs}（ctest 目标 \texttt{rpo\_gui\_e2e}，
        tests/CMakeLists.txt:502--507，标签含 \texttt{cross\_validation}）；
  \item 已知限制：LossEconomic 内层不直接为电压最优（2869 上 vdev 0.1392$\to$0.1401
        基本未降）；3000+ 母线完整搜索受时限约束。
\end{itemize}

% 第 7 章：三相混合 OPF
\section{三相混合 OPF（活跃研发）}\label{sec:threephase}

\subsection{功能定位与入口}
\compmeta{opf::phase\_hybrid::solve\_three\_phase\_hybrid\_opf}{\srcpath{include/hacdcpf/optimal\_power\_flow/three\_phase\_hybrid\_opf.hpp}}
入口（three\_phase\_hybrid\_opf.hpp:186；实现
\srcpath{src/optimal\_power\_flow/three\_phase\_hybrid\_opf.cpp:2328--2933}）。
输入是低层 \texttt{ThreePhaseHybridOPFCase}（相节点 $Y_{ac}$、$G_{dc}$、负荷、限值，
hpp:66--92），非 rich 模型。配套：
\srcpath{make\_three\_phase\_hybrid\_pf\_case}（OPF 运行点 $\to$ PF 复核，:2935--3011）、
序列/分支参数化驱动（:3013--3124）。
\textbf{README 声明：活跃研发/论文验证性质，接口与产物格式仍可能调整。}

\subsection{数学模型（直角坐标）}
变量（Layout，:31--51，装配 :482--499）：
\begin{align*}
x=[\,&e;\ f\ (\text{全相节点电压实/虚部});\ p_g;\ q_g;\ u_{dc};\\
&p_{ac}^{phase};\ q_{ac}^{phase};\ gfm\_e;\ gfm\_f;\ p_{dc}].
\end{align*}

\paragraph{等式（:533--645）}
每相节点：$\mathrm{Re}/\mathrm{Im}\big(V\odot\overline{(YV+I_{fixed})}\big)+P/Q_{load}-P_g/Q_g-P_{ac}/Q_{ac}=0$；
DC：$u_k(G_{dc}u)_k+p_{d,k}-\sum p_{dc}=0$；
换流器耦合 $p_{dc}+\eta\sum_{phase}p_{ac}=0$（效率单参数，:574--576）；
DC 参考电压、AC 参考；换流器控制行按模式：
\begin{itemize}
  \item \textbf{EqualPhasePower}：各相 $p_{ac}/q_{ac}$ 相等（:621--637）；
  \item \textbf{GridFollowingPLL}（正序电流）：$s_\varphi v_a=\mathrm{rot}_\varphi\,v_\varphi s_a$（:605--620）；
  \item \textbf{GridFormingDroop}（虚拟阻抗 Norton）：
        $\overline{e_\varphi}\,v-|v|^2-\overline{z}\,s=0$（:584--604）。
\end{itemize}

\paragraph{不等式（:871--913）}
$v_{min}^2\le|v|^2\le v_{max}^2$（每相节点）；
VUF $|v_2|^2\le\fld{vuf\_max}^2|v_1|^2$（默认 \fld{vuf\_max}=0.03，hpp:81）；
换流器 $\sum_{phase}(p_{ac}^2+q_{ac}^2)\le s_{max}^2$；
每相电流 $p_{ac}^2+q_{ac}^2\le I_{max}^2|v_{phase}|^2$。

\paragraph{目标（:1039--1055）}
\[
f=\sum_g(c_2P_g^2+c_1P_g)+10^{-4}q_g^2+10^{-3}(p_{ac}^2+q_{ac}^2)+10^{-4}p_{dc}^2
\quad(\text{小正则项}).
\]

\paragraph{变量界与常数（:439--444, 1671--1710）}
效率 clamp $[0.01,1]$（:439）；相电流上限缺省 $s_{max}/\sqrt{n_{phases}}$（:441--444）；
界：$e,f\in\pm1.2$、$gfm\_e/f\in\pm1.5$、$p_{ac}/q_{ac}/p_{dc}\in\pm s_{max}$
（unity PF 时 $q_{ac}$ 钉 0）。

\paragraph{图降阶变体 GraphReduced（:362--424）}
稀疏 Kron 消元 + 恢复算子（\srcpath{include/hacdcpf/graph/sparse\_kron\_reduction.hpp:16--42}）：
\begin{itemize}
  \item \textbf{保留集} = 非零负荷/固定电流节点 $\cup$ 发电机节点 $\cup$ 换流器节点
        $\cup$ 参考节点，其余被动节点消元（:363--378）；
  \item \fld{SparseKronOptions}：\fld{max\_front}=12、\fld{max\_nnz\_ratio}=2.0、
        \fld{pivot\_tolerance}$=10^{-12}$、\fld{drop\_tolerance}$=10^{-13}$
        （最小 front 排序，超 cap 拒绝消元）；
  \item 恢复算子 $v_{full}=R\,v_{model}$（\srcpath{sparse\_kron\_recovery\_operator}），
        \fld{recovery\_rows} 每全节点一行；空行即报
        ``reference-disconnected passive components''（:417--424）；
  \item \textbf{不等式恒在全节点空间评估}（$2\times$full\_n 电压行），经 $R$ 回链到
        模型变量（:871--913，\srcpath{append\_voltage\_norm\_gradient} :976--985）；
        \fld{eliminated\_phase\_nodes} = \fld{recovery\_steps.size()}（:2780--2781）；
  \item \fld{reduction} 证书（reduced / retained / recovery\_steps / original\_size /
        nnz / max\_front / elapsed\_ms / error）原样进结果。
\end{itemize}

\subsection{算法与收敛认证}
后端：\texttt{Ipopt}（engine NLPModel 适配，acceptable=$10\times$tolerance，
:2542--2553）或 \texttt{NativeIPM}（MIPSolvers NativeIPMAdapter，filter 全局化，
\fld{tol\_primal}$=\min(\fld{tolerance},10^{-7})$、\fld{scale\_problem}=false，
:2559--2689, 2640--2648）；求解后对偶缺失时（Ipopt 腿）在解点重跑对偶初始化，
$\fld{recovered\_dual\_residual}\le10^{-6}$ 才采用（:2693--2708）。

\paragraph{对偶初始化 \srcpath{initialize\_primal\_dual\_start}（:1760--2069）}
门控 $\fld{primal\_neighborhood}\le10^{-5}$；slack 重建 $\max(-h,2\times10^{-10})$
（有种子）/ $\max(-h,10^{-2})$；$\mu=\mathrm{barrier}/\mathrm{slack}$，
barrier$=\max(\mu_{min},\mu_{init})$，有种子时
$\mu_{init}=\mathrm{clamp}(\bar s^{\mathsf T}\bar\mu/n,\ \mu_{min},\ 0.1)$；
等式对偶优先\textbf{基划分}路径——$\fld{nvar}-\fld{neq}==n_c$ 时每换流器首相
$q_{ac}$ 列标记 \fld{equality\_free\_columns}，解 $J_B^{\mathsf T}\lambda=-r_B$
（:1749--1756, 1874--1944）；否则列缩放最小二乘（$n\le2000$ 用 COD 阈值 $10^{-14}$，
否则 SparseQR），部分种子走 Schur 恢复（:1974--2048）。

\paragraph{原始可行性恢复 \srcpath{restore\_primal\_feasibility}（:2071--2152）}
$\le30$ 次、行缩放 $J$、COD（阈值 $10^{-12}$）或 $JJ^{\mathsf T}+10^{-10}I$ 的
SimplicialLDLT、24 次减半线搜索带界夹、容差 $10^{-7}$。

\paragraph{约束 oracle（:2340--2502）}
初始强制集 = 种子行 $\cup$ \{满足 $h\ge-\max(m_0,m_a)$ 的行（$m_0$=\fld{oracle\_initial\_margin}，
$m_a$=\fld{oracle\_activation\_margin}）\} $\cup$ \textbf{全部换流器行}
（row $\ge2\times$full\_n$+n_{3ph}$，:2369--2376）；
每轮暖链（primal + equality\_dual + inequality dual/slack 头 nineq 截断）且
round$>$0 强制关 \fld{warm\_start\_with\_ipopt}（:2404--2418）；
\fld{oracle\_probe\_iterations}$>$0 时先预测轮（限迭代）再精确校正轮
（added==0 触发 \fld{exact\_corrector\_required}）；added==0 即停；
轮尽返回 converged=false ``enrichment round limit reached''；
trace 字段（round / enforced / added / iterations / converged / runtime\_ms /
max\_full\_inequality）；\fld{max\_omitted\_inequality} 报告被省略行的最大违反（诚实）。

收敛门 \fld{audited\_kkt}：primal/dual/comp 残差与电压/换流器违反全 $\le10^{-6}$
才算 converged（:2918--2924）。

\paragraph{初始化链（:1260--1632）}
AC Newton $\le20$ 次（tol $10^{-10}$，线搜索要求恢复电压 $\in(0.5,1.3)$），
失败回退 Z-bus 同伦：20 加载步 $\times\le80$ 次、松弛 $\alpha$ 初 0.65、
电压坍缩守卫 $|v|>0.2$、收敛 $10^{-11}$；随后 3 轮``投影+AC PF''交替（:1602--1607）；
发电机出力由网络解回算并夹到 $[p_{min},p_{max}]$（:1609--1632）。

\paragraph{换流器控制投影与动态平衡恢复}
投影 \srcpath{project\_converter\_control}（:1516--1579）：GFM 解
$\bar e=\big(\overline{z}\,S_{total}+\sum|v_\varphi|^2\big)/\big(\sum v_\varphi\overline{\mathrm{rot}_\varphi}\big)$，
再 $i_\varphi=(\mathrm{rot}_\varphi e-v_\varphi)/z$、$s_\varphi=v_\varphi\bar\imath_\varphi$；
GFL $\bar\imath_a=S_{total}/(3v_1)$、$s_\varphi=v_\varphi\,\mathrm{rot}_\varphi\,\bar\imath_a$。
恢复 \srcpath{recover\_dynamic\_equilibria}（:2234--2324）：dq 功率态 = 三相总功率/3；
PLL $\theta=\arg(v_1)$、积分器 $=-k_pv_q/k_i$、残差
$\max(|i_0|,|i_2|,|k_pv_q+k_i\xi|)$；GFM 内电势=$e_a$、$v_{ref}=\mathrm{mean}|v_\varphi|$、
电压积分器 $=(|e_a|-v_t)/k_{int}$、残差
$\max_\varphi|\mathrm{rot}_\varphi e-v_\varphi-zi_\varphi|$。

\fld{verify\_derivatives}（:2783--2851）：Jacobian/Hessian 与中心差分对照，
步长 $10^{-5}$、相对误差 $/\max(1,|fd|)$，Hessian 用合成
$\lambda_i=0.01((i\bmod 7)-3)$、$\nu_i=0.01((i\bmod 5)+1)$，误差记入结果。

\paragraph{结果结构（ThreePhaseHybridOPFResult，hpp:138--184）}
顶层字段：\fld{converged}/\allowbreak\fld{status}/\allowbreak\fld{solver}/\allowbreak\fld{iterations}/\allowbreak
\fld{objective}/\allowbreak\fld{runtime\_ms}（hpp:139--142, 154--155）；
规模统计字段为 \fld{variables}/\allowbreak\fld{equalities}/\allowbreak\fld{inequalities}
（hpp:143--145，无 num\_ 前缀）；
\fld{verify\_derivatives} 开启时 FD 对照误差记入
\fld{max\_equality\_jacobian\_error}/\allowbreak\fld{max\_inequality\_jacobian\_error}/\allowbreak
\fld{max\_lagrangian\_hessian\_error}（hpp:169--171，默认 0）。
逐字段表见第~\ref{sec:options}~章，本节不重复维护。

\subsection{OPF$\to$PF 回放与序列/分支驱动}
\paragraph{\srcpath{make\_three\_phase\_hybrid\_pf\_case}（:2935--3011）}
前置校验（converged 且尺寸匹配，否则 throw）；换流器 \fld{p\_set}/\fld{q\_set}
$=\sum$ 相功率；\fld{v\_dc\_set} 取 dc\_reference 否则 \fld{v\_dc\_start}；
控制模式映射：GFL $\to$ GridFollowingVdcQ（该端子是 DC 参考时）/ GridFollowingPQ，
GFM $\to$ GridFormingVdc / GridFormingVoltage 并透传
\fld{internal\_voltage\_positive}，Equal $\to$ EqualPhaseVdcQ / EqualPhasePQ。

\paragraph{序列/分支参数化驱动（:3013--3124）}
固定拓扑校验（$y_{ac}/g_{dc}$ 维数与 nnz、设备数一致，否则 throw）；
ModelData 一次性构建后逐案例浅拷贝 + 换 source/fixed\_current；
暖链：primal + equality\_dual + inequality dual/slack 头 nineq +
\fld{oracle\_seed\_rows}=上一案例的 \fld{enforced\_inequality\_rows} +
关 \fld{warm\_start\_with\_ipopt}；sequence 把模型构建耗时计入首个结果；
branches 要求 \fld{certified\_base.converged} 否则 throw。

\subsection{参数表（ThreePhaseHybridOPFOptions）}
定义 three\_phase\_hybrid\_opf.hpp:104--126。
\begin{paramtable}
\fld{variant} & -- & 枚举 & Full / GraphReduced & Full & 模型变体\\
\fld{backend} & -- & 枚举 & Ipopt / NativeIPM & Ipopt & 求解后端\\
\fld{reduction\_options} & -- & 结构 & -- & 默认 & SparseKron：max\_front=12、max\_nnz\_ratio=2.0、pivot/drop 容差\\
\fld{max\_iterations} & -- & 次 & -- & 300 & 迭代上限\\
\fld{tolerance} & -- & -- & -- & $10^{-7}$ & 收敛容差\\
\fld{warm\_start\_with\_ipopt} & -- & bool & -- & false & Ipopt 可行性暖解\\
\fld{verify\_derivatives} & -- & bool & -- & false & FD 导数对照\\
\fld{verbose} & -- & bool & -- & false & 初始化/oracle 诊断输出至 stderr，并透传 IPM 日志（:359, 2334, 2648）\\
\fld{use\_constraint\_oracle} & -- & bool & -- & false & 活动集 oracle\\
\fld{oracle\_initial\_margin} & -- & -- & -- & $10^{-3}$ & oracle 初始裕度\\
\fld{oracle\_activation\_margin} & -- & -- & -- & $10^{-5}$ & oracle 激活裕度\\
\fld{oracle\_probe\_iterations} & -- & 次 & -- & 0 & 预测轮迭代上限；0=禁用预测/校正两段式\\
\fld{oracle\_max\_rounds} & -- & 轮 & -- & 8 & 富集轮次上限\\
\fld{oracle\_seed\_rows} & -- & 行号 & -- & 空 & 种子行（越界即 throw，:2362--2368）\\
\fld{enforced\_inequality\_rows} & -- & 行号 & -- & 空 & 高级覆盖：空=全部行启用；校验有序/无重/在界\\
\fld{primal\_start} 等四项 & -- & 向量 & -- & 空 & 暖启动\\
\end{paramtable}

\subsection{验证与校核}
算例：3 母线$\times$3 相 + 2 DC 节点手工算例（test\_three\_phase\_hybrid\_opf.cpp，427 行）：
\begin{itemize}
  \item Full vs GraphReduced 同解：目标相对差 $\le10^{-6}$、电压差 $<10^{-6}$、
        VUF $\le\fld{vuf\_max}+10^{-7}$（:105--138）；
  \item GFL-PLL / GFM-droop 动态平衡：电流 VUF $<10^{-7}$、负载率 $\le1+10^{-7}$、
        平衡残差 $<10^{-7}$；解析 Jacobian/Hessian 误差 $<10^{-5}$，
        并用 dynamics 设备模型独立复核导数 $<10^{-7}$（:140--248）；
  \item 独立三相混合 PF 复现 OPF 运行点：电压差 $<2\times10^{-6}$（PF tol $10^{-9}$，:250--283）；
  \item Native Full 认证 graph-reduced 解的 primal-dual 迁移（:285--367）；
  \item 约束 oracle 保解且省略行 $<-$activation\_margin（:369--427）；
  \item 基准工具：\srcpath{tools/phase\_hybrid\_opf\_benchmark.cpp}（Full vs GraphReduced vs
        continuation 目标/电压误差与加速比，GFL/GFM 灵敏度场景矩阵）、
        \srcpath{tools/phase\_graph\_reduction\_benchmark.cpp}。
\end{itemize}
GUI 适配范围（README:395）：精确覆盖恒功率星形相负荷/纯电阻 DC 支路/直接 VSC；
Delta/ZIP、DCDC、能量路由器显式拒绝；AC 支路热限与 VSC 调制比尚未进入相域模型，
经 \fld{scope}/\fld{model\_limitations} 标注。

% 第 8 章：power_models AML 求解器无关建模层
\section{power\_models：AML 求解器无关建模层}\label{sec:powermodels}

\subsection{定位（诚实口径）}
\compmeta{hacdcpf::power\_models}{\srcpath{include/hacdcpf/power\_models/}}
五个 builder（acopf/acdcopf/dc\_opf/lindistflow/scuc）基于兄弟仓库 MIPSolvers 的
AML（代数建模层），经 \srcpath{include/hacdcpf/aml/aml.hpp} 转发
（\texttt{namespace hacdcpf::aml = mipsolvers::aml}）。
\textbf{这是独立/备用的公式化层，并非生产 OPF 主链路}：全仓 grep 显示
\srcpath{power\_models::solve\_*} 在生产 \texttt{src/} 中无任何调用方，仅被测试调用；
生产 OPF 走 \srcpath{src/optimal\_power\_flow/}（第~\ref{sec:acopf}--\ref{sec:dcopf}~章）。
该定位声明现已\textbf{形式化为结果字段}：\fld{AMLBuildResult} 默认携带
\fld{model\_scope}=``experimental-aml-builder:not-production-runtime'' 与一条
\fld{model\_limitations}（AML builder 是实验性形式化路径，不被生产
\texttt{solve\_ac\_opf}/\texttt{solve\_dc\_opf} 运行时使用，
aml\_build\_result.hpp:17--19），各 builder 结果按此口径进一步细化（见后文各节）。
\textbf{命名冲突警示}：\fld{power\_models::DCOPFResult} 与 \fld{opf::DCOPFResult}
同名不同物（dc\_opf\_model\_builder.hpp:7--8, 46）。

\subsection{AML 建模原语}
\srcpath{MIPSolvers/include/mipsolvers/aml/model.hpp}：
集合（\fld{add\_set}/\fld{add\_ordered\_set}）、参数、变量
（\fld{add\_var}，\fld{VarType::Continuous/Binary}，逐键 \fld{set\_lb/set\_ub/fix}）、
目标（\fld{minimize} 重载 LinearExpr/QuadExpr/NonlinearExpr）、
非线性表达式 DAG（\fld{nl\_var/nl\_mul/nl\_sq/nl\_cos/nl\_sin} 等，arena 分配）、
约束（线性 \fld{add\_constraint}、索引族 \fld{add\_constraints}、非线性
\fld{add\_nl\_constraint}）、求解 \fld{solve(SolveOptions)}、辅助
\fld{check\_bounds}/\fld{set\_nlp\_x0}/\fld{write\_lp}。
\fld{SolveOptions}：\fld{solver\_name}（"" = 自动）、\fld{time\_limit\_sec}、
\fld{mip\_gap\_tol}=1e-4、\fld{allow\_fallback}。
默认分派链（MIPSolvers/src/engine/strategy/dispatcher.cpp:61--81）：
LP $\to$ NativeIPMLP/NativePDLP/NativeLCQP/Gurobi/HiGHS；
QP $\to$ Gurobi/NativeLCQP；NLP $\to$ Ipopt/NativeIPM/NativeNLP；
MILP $\to$ HiGHS/Gurobi/NativeBranchAndCut；MINLP $\to$ Scip/NativeBranchAndCut。
\fld{SolveResult}：TerminationStatus/PrimalStatus/DualStatus、
\fld{objective\_value}/\fld{optimality\_gap}、\fld{var\_value}/\fld{dual}/
\fld{reduced\_cost}、IIS 数据。

\subsection{AC OPF builder（AML 非线性）}
头文件 \srcpath{ac\_pf\_model\_builder.hpp}；实现 src/power\_models/acopf\_builder.cpp（557 行）。
数据转换 \srcpath{to\_acopf\_data}（:31--158）；求解 \srcpath{solve\_acopf}（:163--555），
默认 \fld{solver\_name}=``NativeIPM'' 且 \fld{allow\_fallback}=false（:172--178，
注释：NativeNLP 缺线搜索，病态时可能堆损坏，永不回退）。

\paragraph{数学模型}
变量 $V_m\in[v_{min},v_{max}]$（默认 [0.9,1.1]）、$V_a\in[-\pi,\pi]$、
$P_g\in[p_g^{min},p_g^{max}]$、$Q_g\in[q_g^{min},q_g^{max}]$（pu，:204--211）。
目标（:234--257）：
\[
\min\sum_g\big(a_2P_g^2+a_1P_g+a_0\big),\qquad a_2=c_2S_b^2,\ a_1=c_1S_b,\ a_0=c_0.
\]
$\pi$ 支路导纳（branch\_admittance.hpp:42--77，复变比 $\tau e^{-j\varphi}$，
$y_s=1/(r+jx)$）：
\[
Y_{ff}=\frac{y_s+jb_c/2}{\tau^2},\quad Y_{tt}=y_s+\tfrac{jb_c}{2},\quad
Y_{ft}=-\frac{y_se^{+j\varphi}}{\tau},\quad Y_{tf}=-\frac{y_se^{-j\varphi}}{\tau}.
\]
支路潮流（极坐标完整 AC，非二阶锥松弛，:341--407）：
\[
P_f=G_{ff}V_f^2+V_fV_t\big(G_{ft}\cos\theta_{ft}+B_{ft}\sin\theta_{ft}\big),\quad
Q_f=-B_{ff}V_f^2+V_fV_t\big(G_{ft}\sin\theta_{ft}-B_{ft}\cos\theta_{ft}\big),
\]
$P_t/Q_t$ 同构（$\theta_{ft}=\theta_f-\theta_t$）。
节点平衡（:294--414，注意并联符号：有功 $-G_sV^2$、无功 $+B_sV^2$）：
\[
\sum_{g\in i}P_g-P_{d,i}-G_{s,i}V_i^2-\sum_{\ell\ni i}P_\ell=0,\qquad
\sum_{g\in i}Q_g-Q_{d,i}+B_{s,i}V_i^2-\sum_{\ell\ni i}Q_\ell=0.
\]
参考角 $V_a[slack]=0$ 以等式而非 fix 实现（:417--423，保持 IPM 可行）；
热限仅 \fld{rate\_a\_pu}$>0$ 的支路两端 $P^2+Q^2\le S_{max}^2$（:425--443）。
暖启动 x0 顺序 = 全部 $V_m$、$V_a$、$P_g$、$Q_g$（:259--269）。

\paragraph{数据转换要点（to\_acopf\_data）}
仅 \fld{in\_service} 设备；\fld{pg\_max}$=\max(p_{max},p_{min})$ 防反转（:74）；
非有限 Q 界回退 $\mp10^4$（:75--78）；external grid 仅当 $c_2\ne0\lor c_1\ne0$
时作为可调度 slack 电源加入（注释 :95--97：让 IPM Hessian 在无常规机组时非奇异）；
零阻抗支路（$r^2+x^2<10^{-12}$）静默跳过（:24,128--129，依赖 projection 层合并）；
无 ref 母线时强制 \fld{buses[0].is\_ref}（:148--155）。

\paragraph{结果与诚实口径}
\srcpath{AMLBuildResult}（aml\_build\_result.hpp:16--38）：诚实口径字段
\fld{model\_scope}（默认 ``experimental-aml-builder:not-production-runtime''）与
\fld{model\_limitations}（默认一条：AML builder 是实验性形式化路径，不被生产
\texttt{solve\_ac\_opf}/\texttt{solve\_dc\_opf} 运行时使用）（:17--19）；
\fld{solve\_result}、
\fld{obj\_per\_h}（\$/h，用 MW 重算核验而非 solver 原始目标，:466--476）、
违例指标 \fld{max\_p/q/thermal\_viol\_pu} 与 \fld{max\_dc\_p\_viol\_pu}、字符串 id 键的结果映射
（\fld{vm\_pu}/\fld{va\_rad}/\fld{pg\_mw}/\fld{qg\_mvar}/\fld{pf\_mw}/\fld{qf\_mw}，
DC 侧映射 \fld{vdc\_pu}/\fld{pac\_mw}/\fld{qac\_mvar}/\fld{pdc\_mw}）。
无负荷削减变量——不可行时整个问题不可行；external grid 成本项是 Hessian 正则化手段，
非市场行为模型。

\subsection{混合 AC/DC OPF builder}
头文件 \srcpath{hybrid\_opf\_model\_builder.hpp}；实现 src/power\_models/acdcopf\_builder.cpp
（766 行）。\fld{ACDCOPFData} = \fld{ACOPFData} + dc\_buses/dc\_branches/converters/
dcdc\_converters。求解 \srcpath{solve\_acdcopf}（:232--764），同样锁 NativeIPM 禁回退。

\paragraph{数学模型}
新增变量 $V_{dc,k}\in[v_{dc}^{min},v_{dc}^{max}]$（默认 [0.9,1.1]）、换流器
$P_{ac,c},Q_{ac,c}$（正 = 注入 AC 母线）、DC/DC 输出 $P_{out}$（:269--319）。
目标仅 AC 机组成本（:322--336）。DC 支路潮流（对 $V_{dc}$ 线性，:508--524）：
$P_{ij}=(V_{dc,i}-V_{dc,j})/r_{ij}$。
VSC 耦合（无独立等式，直接注入 DC 平衡，:527--566）：
\[
I_{ac,c}=\frac{\sqrt{P_{ac,c}^2+Q_{ac,c}^2+\varepsilon^2}}{V_{m,a}}\ (\varepsilon=10^{-6}),\qquad
P_{dc,c}=-P_{ac,c}-\big(a_c+b_cI_{ac,c}+c_cI_{ac,c}^2\big),
\]
$a=\fld{loss\_mw}/S_b$、$b=\fld{loss\_percent}/100$、$c=1-\eta$（:154--156）。
DC/DC：输出母线 $+P_{out}$，输入母线 $-P_{in}$（:583--592）：
\[
P_{in}=\frac{P_{out}}{\eta}+\frac{r_{eq}\,(P_{out}/\eta)^2}{V_{dc,in}^2},
\]
其中 $r_{eq}=\fld{r\_eq\_pu}\ge0$ 为等值串联电阻（hpp:63，double 默认 $0.0$；
\srcpath{to\_acdcopf\_data} 经 $\max(0,\fld{r\_eq\_pu})$ 保留而不再丢弃，:179），
$r_{eq}=0$ 时退化为纯效率模型，$r_{eq}>0$ 时追加的
$r_{eq}(P_{out}/\eta)^2/V_{dc,in}^2$ 为非线性 $I^2R$ 损耗项（:586--591）。
按控制模式加一条等式——
Power：$P_{out}=p_{ref}$；Droop：$P_{out}-k_{droop}(V_{dc,out}-v_{ref})=p_{ref}$；
Voltage：$V_{dc,out}=v_{ref}$（输出母线已是 VSC Vdc 松弛则跳过）（:594--605）。
DC 母线平衡 $\mathrm{Pdc\_acc}_k=0$（:608--612）；Vdc 松弛 $V_{dc}[k]=V_{dc,set}$（:614--621）。
Vdc 松弛标记经共享解析器 \srcpath{resolve\_device\_control\_role} 判定（:112--126），
与潮流/短路/暂态口径一致。

\paragraph{诚实口径}
\fld{ACDCOPFBuilderResult} 自带 \fld{model\_scope}=``experimental-aml-hybrid-acdc-opf''
与两条 \fld{model\_limitations}（hybrid\_opf\_model\_builder.hpp:83--86）：
本层为实验路径，生产混合 OPF 用 parity 形式化；DC/DC 仅建模\textbf{正向传输}的
smooth 效率+$I^2R$ 模型，不表示双向符号切换。
DC/DC 的 $I^2R$ 损耗\textbf{现已建模}（见上式；本手册早前“省略 $I^2R$、eta 仅捕获
一阶损耗”的声明已过时），DC 平衡审计改用与约束完全相同的效率+$I^2R$ 约定
（:745--758，$V_{dc,in}$ 取下限 $0.1$ 防除零），修复了审计残差与约束模型不一致、
把正常 $I^2R$ 损耗误报为平衡违规的问题——\fld{max\_dc\_p\_viol\_pu} 不再因
DC/DC 项虚高。DC 支路仍无容量约束；孤立 DC 母线被静默剔除且其负荷不被服务
（:52--56）。

\subsection{DC OPF builder（AML LP）}
头文件 \srcpath{dc\_opf\_model\_builder.hpp}；实现 src/power\_models/dc\_opf\_builder.cpp
（205 行）。输入 \fld{DCOPFData}：bus\_ids、slack\_bus、GenTuple
\{id,bus,pmin,pmax,cost\}、BranchTuple\{from,to,susceptance,max\_flow\}、demands。
变量 $p_g\in[p_{min},p_{max}]$~MW；$\theta_i\in[-10^6,10^6]$——\textbf{刻意宽角界}
（:78--81：紧界会遮蔽线路潮流约束对偶、阻塞时给出错误 LMP）；slack 钉 0。
目标 $\min\sum_gc_gp_g$（线性 \$/(MWh)）。
节点平衡（:97--129）：$\sum_{g\in i}p_g-D_i-\sum_{\ell}b_\ell(\theta_i-\theta_j)=0$，
支路潮流 $P_\ell=b_\ell(\theta_{from}-\theta_{to})$；
潮流限 $-P_{max}\le b_\ell(\theta_f-\theta_t)\le P_{max}$（:131--161）。
后端 auto 走 LP 链。结果：\fld{gen\_dispatch\_MW}、\fld{voltage\_angle\_rad}、
\fld{branch\_flow\_MW}（键 ``from->to''）、\fld{lmp\_per\_MWh}
（平衡约束对偶；:189--199）。诚实口径：纯 LP 线性成本、无损耗、无 V/Q、无切负荷；
\textbf{无直接单元测试}。

\subsection{LinDistFlow builder（配网线性 LP）}
头文件 \srcpath{lindistflow\_builder.hpp}（:3--31 自带数学文档）；实现
\srcpath{src/power\_models/lindistflow\_builder.cpp}（186 行）。适用辐射状配电网。
变量：支路 $P_{ij},Q_{ij}$、母线电压平方 $v_i\in[v_{min}^2,v_{max}^2]$
（默认 $0.81$--$1.21$ pu$^2$）。
约束：根电压固定 $v_{root}=\fld{root\_v\_sq\_pu}$；热限为\textbf{分离盒式近似}
$|P_{ij}|\le S_{max}$ 且 $|Q_{ij}|\le S_{max}$（hpp:22 标注 approximate）；
节点平衡 $\sum P_{ij}-P_{k(i),i}=-p_{load,i}$（辐射网每非根母线恰一条入支路）；
LinDistFlow 电压降（Baran \& Wu 线性化，忽略 $I^2r$ 损耗项，:133--148）：
\[
v_j=v_i-2\,(r_{ij}P_{ij}+x_{ij}Q_{ij}).
\]
目标 $\min\sum_{i\ne root}(const-v_i)$——电压调节代理目标（hpp:27--28）。
结果 \fld{voltage\_sq\_pu}（\textbf{单位 pu$^2$}）、\fld{branch\_p\_mw}/\fld{branch\_q\_mvar}。
诚实口径：头文件注释现已与实现一致——平衡方程右端为 $-p_{load,i}$（hpp:15--16，
实现亦只含负荷项，cpp:114--115），并新增 Scope 说明：仅固定负荷，
不建 DG、可控源或切负荷变量（hpp:30--31）；无直接单元测试。

\subsection{SCUC builder（MILP）}
头文件 \srcpath{scuc\_builder.hpp}；实现 src/power\_models/scuc\_builder.cpp（310 行）。
\textbf{非网络 SCUC}：只有系统级功率平衡，无潮流/网络安全约束。
变量 $u_{g,t}\in\{0,1\}$、$p_{g,t}\ge0$、$s_{g,t}\ge0$（启动指示连续松弛，
靠目标启动成本迫使取 0/1，:71）。
目标（:76--91）：$\min\sum_g\sum_t(c_gp_{g,t}+S_gs_{g,t}+C_gu_{g,t})$。
约束：平衡 $\sum_gp_{g,t}=D_t$；容量 $p_g^{min}u\le p\le p_g^{max}u$；
爬坡 $-R^{dn}\le p_{g,t}-p_{g,t-1}\le R^{up}$；启动逻辑 $s_{g,t}\ge u_{g,t}-u_{g,t-1}$；
最小开/停机（Beta，窗口末端截断）：
\[
\sum_{j=t}^{t+W-1}u_{g,j}\ge W(u_{g,t}-u_{g,t-1}),\qquad
\sum_{j=t}^{t+W-1}(1-u_{g,j})\ge W(u_{g,t-1}-u_{g,t}),\quad W=\min(T_{up/dn},\ n_T-t);
\]
旋转备用 $\sum_gp_g^{max}u_{g,t}\ge D_t+R_t$（仅当提供）。
后端 auto 走 MILP 链（HiGHS / Gurobi / NativeBranchAndCut）。
结果 \fld{commitment}/\fld{dispatch\_MW}/\fld{startup\_event}/\fld{price\_per\_MWh}
（平衡对偶，仅 LP 松弛/MIP LP 可提供，:296--304），另含诚实口径字段
\fld{model\_scope}=``experimental-aml-scuc:system-balance-no-network''、两条
\fld{model\_limitations}（只系统级平衡、无网络约束；MILP 平衡对偶不是认证出清价，
需固定组合后重跑 LP 定价），以及 \fld{price\_valid}=false 与
\fld{price\_validity\_reason}=``MILP dual prices are unavailable or economically
uncertified.''——先前“MILP 下价格对偶不可得”的声明已形式化为这两个字段
（scuc\_builder.hpp:86--92）。诚实口径：无网络约束、无 N-1；
min-up/down 与 reserve 为 Beta 扩展；无直接单元测试
（市场 SCUC 测试走 \texttt{src/market/} 自有实现，与此 builder 无关）。

\subsection{求解后端汇总}
\begin{center}
\footnotesize
\begin{tabular}{@{}lllll@{}}
\toprule
Builder & 问题类 & 默认 solver & allow\_fallback & 实际链\\
\midrule
AC OPF & NLP & NativeIPM & \textbf{false} & 仅 NativeIPM\\
混合 AC/DC OPF & NLP & NativeIPM & \textbf{false} & 仅 NativeIPM\\
DC OPF & LP & auto & true & NativeIPMLP$\to$PDLP$\to$LCQP$\to$Gurobi$\to$HiGHS\\
LinDistFlow & LP & auto & true & 同 LP 链\\
SCUC & MILP & auto & true & HiGHS$\to$Gurobi$\to$NativeBranchAndCut\\
\bottomrule
\end{tabular}
\end{center}

\subsection{验证与校核}
\textbf{仅两个测试目标直接覆盖 AML 层}：
\begin{itemize}
  \item \srcpath{tests/test\_hacdcpf.cpp:234--261}：孤立 DC 母线被剔除
        （\fld{dc\_buses.size()}==2、支路剔除）；
  \item \srcpath{tests/test\_converter\_coordination.cpp:1312--1590}：控制角色解析器
        遍历 7 种 VSC 模式；Vdc 松弛选择与解析器一致；AC 参考与 grid-forming 解析一致；
        DC/DC 三控制模式——算例设 \fld{r\_eq\_pu}=0.01（:1541），断言数据转换保留
        \fld{r\_eq\_pu}（:1555）；Voltage：\fld{vdc\_pu}[``DC2'']$\approx1.03$
        （margin $10^{-3}$）、Power：\fld{pdcdc\_mw}$\approx p_{ref}$（margin $10^{-2}$）
        且 \fld{max\_dc\_p\_viol\_pu}$<10^{-4}$（:1575，验证含 $I^2R$ 的审计口径）、
        Droop：\fld{pdcdc\_mw}$\approx p_{ref}+k(V_{dc}-v_{ref})S_b$（margin $2\times10^{-2}$）。
\end{itemize}
AC OPF / DC OPF / LinDistFlow / SCUC 的 AML 版本均无直接测试；
\srcpath{tests/test\_acopf\_dcopf\_crossval.cpp} 验证的是生产 \texttt{opf::} 求解器
而非本层（:20--21），引用时须区分。

% 第 9 章：选项、结果与诊断参考
\section{选项、结果与诊断参考}\label{sec:options}

\subsection{总约定：单位与索引空间}
\begin{itemize}
  \item \textbf{相角一律弧度}：\fld{ACOPFResult::va} 与 \fld{DCOPFResult::va} 均以
        \textbf{弧度}存储——parity IPM 路径直接取解向量（ac\_opf.cpp:2481）、
        经济调度兜底拷贝潮流结果（潮流结果同为弧度，ac\_opf.cpp:1800）、
        DC OPF 取 $\theta$ 变量（dc\_opf.cpp:702--705）；RPO 的
        \fld{va\_before}/\fld{va\_after} 同约定。``\texttt{va\_deg}'' 字样仅为历史命名。
        度（$^\circ$）换算只发生在展示层：GUI/HTTP 出口 $\times180/\pi$
        （run\_gui\_server.cpp:14016；OPF 母线结果 JSON 直接以 \texttt{va\_rad} 键输出，:6853）。
  \item \textbf{索引空间}：结果向量一律回投影到 \textbf{authored 组件序}（原始
        \fld{.index}），内部 canonical 位置不对外；跨域母线经 domain-qualified 映射。
\end{itemize}

\subsection{ACOPFOptions 全字段表}
定义 opf\_options.hpp:49--123；$\le0$ 的数值字段在 ac\_opf.cpp:2713--2748 重置为默认。
\label{tab:acopfopts}
\begin{paramtable}
\fld{ac\_solver\_backend} & -- & 枚举 & 四档 & Auto & Auto / ParityIPM / Ipopt / EconomicDispatch；分派点 ac\_opf.cpp:2937--2956\\
\fld{objective} & -- & 枚举 & 四档 & Economic & Economic / VoltageDeviation / ActiveLoss / VoltageDeviationAndLoss\\
\fld{voltage\_target\_pu} & $V^*$ & pu & 0.95--1.05 & 1.0 & 电压偏差目标\\
\fld{voltage\_deviation\_weight} & $w_v$ & -- & $\ge0$ & 1.0 & 电压偏差权重\\
\fld{active\_loss\_weight} & $w_l$ & -- & $\ge0$ & 1.0 & 网损权重\\
\fld{max\_inner\_iterations} & -- & 次 & 20--200 & 80 & 内层迭代上限\\
\fld{max\_outer\_iterations} & -- & 次 & 1--20 & 8 & 外层（障碍）迭代上限\\
\fld{max\_line\_search\_steps} & -- & 步 & -- & 20 & 线搜索步上限\\
\fld{feasibility\_tol} & -- & -- & $10^{-8}$--$10^{-4}$ & $10^{-6}$ & 可行性容差\\
\fld{stationarity\_tol} & -- & -- & $10^{-6}$--$10^{-2}$ & $10^{-6}$ & 平稳性容差\\
\fld{barrier\_mu0} & $\mu_0$ & -- & -- & $10^{-2}$ & 初始障碍参数\\
\fld{barrier\_mu\_reduction} & -- & -- & 0.1--0.5 & 0.2 & $\mu$ 缩减因子\\
\fld{barrier\_mu\_min} & $\mu_{min}$ & -- & -- & $10^{-8}$ & $\mu$ 下限（映射为 IPM 互补容差）\\
\fld{merit\_penalty} & -- & -- & 1--100 & 10.0 & merit 罚系数基数\\
\fld{regularization} & $\delta$ & -- & -- & $10^{-6}$ & KKT 正则化\\
\fld{step\_backoff} & -- & -- & 0.3--0.8 & 0.5 & 步长回退因子\\
\fld{interior\_fraction} & $\alpha_{max}$ & -- & 0.9--0.999 & 0.995 & fraction-to-boundary\\
\fld{ac\_eval\_threads} & -- & 线程 & 0=auto & 1 & AC 求值线程数\\
\fld{enable\_primal\_dual} & -- & bool & -- & false & \textbf{遗留兼容开关（仍生效）}：\fld{ac\_solver\_backend} 分派时被后端枚举改写（ac\_opf.cpp:2937--2956）；与 \fld{use\_parity\_ipm} 同真时进入 parity 全空间路径（:2967）\\
\fld{use\_parity\_ipm} & -- & bool & -- & false & \textbf{遗留兼容开关（仍生效）}，同上；混合 AC/DC 系统在未选 EconomicDispatch 时自动置位（ac\_opf.cpp:2961--2964）\\
\fld{allow\_fallback} & -- & bool & -- & true & 未收敛时 Ipopt 兜底\\
\fld{verbose} & -- & bool & -- & false & 迭代日志\\
\fld{warm\_start} & -- & 指针 & -- & nullptr & 暖启动状态（x/$\lambda$/$\mu$/z）\\
\fld{ac\_pf\_warm\_start} & -- & bool & -- & false & 先跑 AC PF 取电压初值\\
\fld{compute\_homotopy\_tangent} & -- & bool & -- & false & 输出 KKT 同伦切线\\
\fld{homotopy\_t} & $t$ & -- & (0,1] & 1.0 & 同伦参数\\
\fld{objective\_homotopy} & -- & bool & -- & false & 目标同伦两阶段求解（实现 ac\_opf.cpp:2969 起）\\
\fld{homotopy\_dt0} & $dt_0$ & -- & 0.01--0.5 & 0.10 & 同伦初始步长（$\le0$ 回退 0.10）\\
\fld{enforce\_branch\_limits} & -- & bool & -- & true & AC 支路 $|S|$ 约束\\
\fld{enforce\_converter\_capacity} & -- & bool & -- & true & VSC 容量圆\\
\fld{enforce\_converter\_current\_limits} & -- & bool & -- & true & VSC 电流限\\
\fld{enforce\_converter\_modulation\_limits} & -- & bool & -- & true & 调制度/占空比行\\
\end{paramtable}
DCOPFOptions、RPOOptions、ParityOptions、IPMOptions、ThreePhaseHybridOPFOptions
分别见第~\ref{sec:dcopf}、\ref{sec:rpo}、\ref{sec:parity}、\ref{sec:ipm}、\ref{sec:threephase}~章参数表。

\subsection{ACOPFResult 结构}
定义 opf\_result.hpp:72--164。\fld{va} 为\textbf{弧度}（见本章总约定）。索引空间：
全部向量回投影到 \textbf{authored 组件序}（原始 \fld{.index}）。回投影语义
（ac\_opf.cpp:45--76）：\fld{vm}/\fld{va}/\fld{lmp\_p}/
\fld{lmp\_q} 按 \fld{BusVectorSemantics::Intensive}、\fld{dpd}/\fld{dqd} 按
\texttt{Extensive} 反投影；\fld{vdc} 按投影证书 \fld{n\_authored\_dc\_buses} 截断
内部 DC 母线。组件映射族（\texttt{ComponentRef\{original\_index, source\_type\}}，
ac\_opf.cpp:2514--2598）：\fld{gen\_map}、\fld{ren\_map}（source\_type 0=RenewableGen/
1=PVSystem）、\fld{stor\_map}（0=AC/1=DC，\fld{pstor} 向量 AC 在前 DC 在后）、
\fld{dcdc\_map}、\fld{flex\_map}、\fld{er\_port\_map}（source\_type=端口 \fld{.index}）。
\begin{itemize}
  \item \textbf{状态}：\fld{vm}/\fld{va}（弧度，AC 母线）、\fld{vdc}（DC 母线）；
  \item \textbf{调度}：\fld{pg\_mw}/\fld{qg\_mvar}（仅 authored 机组）、
        \fld{external\_grid\_p\_mw}/\fld{external\_grid\_q\_mvar}、
        \fld{dpd\_mw}/\fld{dqd\_mvar}（切负荷）、
        \fld{pac\_mw}/\fld{qac\_mvar}（VSC）、\fld{pren\_mw}/\fld{qren\_mvar}、
        \fld{pstor\_mw}/\fld{qstor\_mvar}（AC 在前 DC 在后）、\fld{pdcdc\_mw}、
        \fld{pflex\_mw}、\fld{er\_port\_p\_mw}/\fld{er\_port\_q\_mvar}；
  \item \textbf{价格}：\fld{lmp\_p}/\fld{lmp\_q}（\$/(MWh)，$\lambda\cdot scale/S_b$，
        ac\_opf.cpp:2603--2615；Ipopt 路径无对偶、LMP 提取跳过，:2105--2108）；
        \fld{lmp\_valid}（bool，默认 false）与 \fld{lmp\_validity\_reason}（:99--102）——
        仅当节点电价来自产生该原始解的同一形式的认证等式乘子时为 true；
  \item \textbf{续传状态}：\fld{ipm\_primal\_state}/\fld{ipm\_equality\_dual\_state}/
        \fld{ipm\_inequality\_dual\_state}/\fld{ipm\_slack\_state}（暖启动，:108--111）、
        \fld{ipm\_tangent\_primal}/\fld{ipm\_tangent\_slack}/
        \fld{ipm\_tangent\_equality\_dual}/\fld{ipm\_tangent\_inequality\_dual}
        （同伦切线，:119--122）；
  \item \textbf{收敛与诊断}：\fld{converged}/\fld{iterations}/\fld{outer\_iterations}/
        \fld{objective}/\fld{max\_constraint\_violation}/\fld{max\_stationarity}/
        \fld{status}/\fld{solver\_path}\{Unknown,NativeAC,ParityIPM\}（:135--142）；
        \fld{profiling}（\texttt{ACOPFProfiling}，逐字段见下表）；
  \item \textbf{诚实口径}：\fld{infeasibility\_hints}（:147）、\fld{audit}
        （\texttt{OpfAudit}，逐字段见下表）、
        \fld{converter\_model\_scope}（按实际启用约束族拼装，:154）、
        \fld{model\_limitations}（\texttt{vector<string>}，:156--159；本次求解路径的
        能力边界声明——被省略的物理、近似或不可用的证书，供 API/GUI 展示）、
        \fld{ac\_switch\_flows}/\fld{ac\_circuit\_breaker\_flows}（API 层补算，:162--163）。
\end{itemize}

\subsubsection*{OpfAudit（opf\_result.hpp:21--34）}
由 \srcpath{verify\_opf\_result}（src/api/hacdcpf.cpp:1348--1451）填充的独立
事后可行性审计；成员函数 \texttt{feasible()} 在 \fld{violations} 为空时返回 true（:33）。
\begin{paramtable}
\fld{audited} & -- & bool & -- & false & \srcpath{verify\_opf\_result} 运行后置 true\\
\fld{max\_power\_balance\_violation\_mw} & -- & MW & -- & 0 & 逐母线 $\max|P_{gen}-P_{load}-P_{loss}|$\\
\fld{max\_voltage\_limit\_violation\_pu} & -- & pu & -- & 0 & 越 $[V_{min},V_{max}]$ 最大电压偏差\\
\fld{max\_branch\_limit\_violation\_pu} & -- & pu & -- & 0 & 超 \fld{rate\_a\_mva} 最大支路负载\\
\fld{max\_gen\_limit\_violation\_mw} & -- & MW & -- & 0 & 越 $[P_{min},P_{max}]$ 最大机组出力偏差\\
\fld{objective\_recomputed} & -- & \$/h & -- & 0 & 由 \fld{pg\_mw}/\fld{qg\_mvar} 重算的成本\\
\fld{objective\_reported} & -- & \$/h & -- & 0 & 求解器报告的成本\\
\fld{objective\_discrepancy\_pct} & -- & \% & -- & 0 & $|\mathrm{recomputed}-\mathrm{reported}|/|\mathrm{reported}|\times100$\\
\fld{violations} & -- & 串表 & -- & 空 & 人类可读违规描述；空即 \texttt{feasible()}=true\\
\end{paramtable}

\subsubsection*{ACOPFProfiling（opf\_result.hpp:47--67）}
\begin{paramtable}
\fld{linear\_solver\_backend} & -- & 串 & -- & 空 & 实际选用的 KKT 线性后端名\\
\fld{analyze\_calls} & -- & 次 & -- & 0 & 符号分析调用数（:49）\\
\fld{factorization\_calls} & -- & 次 & -- & 0 & 数值分解次数\\
\fld{linear\_solve\_calls} & -- & 次 & -- & 0 & 三角回代次数\\
\fld{total\_iterations} & -- & 次 & -- & 0 & 总迭代数（:52）\\
\fld{accepted\_steps} / \fld{rejected\_steps} & -- & 次 & -- & 0 & 被接受/拒绝步数\\
\fld{backend\_escalations} & -- & 次 & -- & 0 & 线性后端升级次数\\
\fld{scaling\_rebuilds} & -- & 次 & -- & 0 & Ruiz 均衡重建次数\\
\fld{warm\_start\_used} & -- & bool & -- & false & 本次是否用了暖启动\\
\fld{initial\_primal\_residual} / \fld{initial\_dual\_residual} & -- & -- & -- & 0 & 初始原始/对偶残差\\
\fld{max\_ac\_p\_balance\_residual\_pu} & -- & pu & -- & 0 & AC 有功平衡残差分族（:60）\\
\fld{max\_ac\_q\_balance\_residual\_pu} & -- & pu & -- & 0 & AC 无功平衡残差分族（:61）\\
\fld{max\_dc\_balance\_residual\_pu} & -- & pu & -- & 0 & DC 平衡残差分族（:62）\\
\fld{max\_converter\_balance\_residual\_pu} & -- & pu & -- & 0 & 换流器平衡残差分族（:63）\\
\fld{max\_other\_equality\_residual\_pu} & -- & pu & -- & 0 & 其余等式残差分族（:64）\\
\fld{max\_nonlinear\_inequality\_violation\_pu} & -- & pu & -- & 0 & 非线性不等式最大违反（:65）\\
\fld{final\_barrier\_mu} & $\mu$ & -- & -- & 0 & 最终障碍参数（:66）\\
\end{paramtable}
残差分族在写入前除回 \fld{scale\_p}/\fld{scale\_q} 还原物理口径
（ac\_opf.cpp:2641--2646）。

\subsection{DCOPFResult 结构}
定义 opf\_result.hpp:181--231。\fld{va} 为\textbf{弧度}（见本章总约定）；
全部向量回投影到 authored 序。
\begin{paramtable}
\fld{va} & $\theta$ & rad & -- & 0 & AC 母线相角（弧度；取自解向量 dc\_opf.cpp:702--705）\\
\fld{pg\_mw} & $P_g$ & MW & -- & 0 & 机组有功调度（authored 序）\\
\fld{external\_grid\_p\_mw} & -- & MW & -- & 0 & 外部电网注入\\
\fld{pf\_mw} & $P_f$ & MW & -- & 0 & 支路首端有功\\
\fld{converged} & -- & bool & -- & false & 收敛标志\\
\fld{iterations} & -- & 次 & -- & 0 & 迭代数\\
\fld{objective} & -- & \$/h & -- & 0 & 目标值（语义见 \fld{objective\_model}）\\
\fld{status} / \fld{solver\_name} & -- & 串 & -- & 空 & 状态串 / 实际求解器名\\
\fld{runtime\_sec} & -- & s & -- & 0 & 求解耗时\\
\fld{lmp} & $\lambda$ & \$/(MWh) & -- & 空 & 节点电价；\fld{compute\_lmp}=false 时清空（dc\_opf.cpp:1190--1194）\\
\fld{branch\_mu\_lower} / \fld{branch\_mu\_upper} & $\mu$ & \$/(MWh) & -- & 空 & 支路拥塞对偶；原生单纯形不能正确处理有界 $P_f$ 变量，支路限生效时\textbf{不可作工程解释}（dc\_opf.cpp:621--634, 680, 783）\\
\fld{branch\_mu\_valid} & -- & bool & -- & false & \textbf{恒 false}，直到实现 KKT 基提取（见第~\ref{sec:dcopf}~章）\\
\fld{branch\_mu\_validity\_reason} & -- & 串 & -- & 空 & 上述判定的原因串（:201）\\
\fld{lmp\_valid} & -- & bool & -- & false & \fld{lmp} 来自所报原始解的等式行对偶时为 true（:205）\\
\fld{lmp\_validity\_reason} & -- & 串 & -- & 空 & 原因串（:206）；注释明确：平衡行对偶认证\textbf{不蕴含}支路拥塞对偶有效\\
\fld{load\_shedding\_mw} & -- & MW & -- & 空 & 逐母线切负荷\\
\fld{total\_load\_shedding\_mw} & -- & MW & -- & 0 & 切负荷总量（:209）\\
\fld{solver\_chain} & -- & 串表 & -- & 空 & 实际走过的后端链，``\texttt{<backend>:<status>}'' 形式（:211--215）\\
\fld{objective\_model} & -- & 串 & QP/LP/LP-PWL & 空 & 目标成本模型：``QP''（真二次）、``LP''（保守线性化）、``LP-PWL''（分段线性，dc\_opf.cpp:1200--1206）\\
\fld{pwl\_segments\_effective} & -- & 段 & 0--1000 & 0 & 二次成本 LP 近似实际生效段数；QP 后端恒 0（dc\_opf.cpp:1203），LP 回退报实际段数（:1206）（:225）\\
\fld{converter\_model\_scope} & -- & 结构 & -- & 全 false & DC OPF 不建模换流器，所有换流器标志保持 false（:227--229）\\
\fld{model\_limitations} & -- & 串表 & -- & 空 & 本次求解能力边界声明（:230）\\
\end{paramtable}

\subsection{RPOResult 结构}
定义 reactive\_power\_opt.hpp:245--291。\fld{va\_before}/\fld{va\_after} 为\textbf{弧度}
（见本章总约定）。离散动作存于 \fld{taps}/\fld{shunts} 子结构向量（逐字段见下两表）。
\begin{paramtable}
\fld{converged} & -- & bool & -- & false & 收敛标志\\
\fld{objective} & -- & -- & -- & 0 & RPO 目标值\\
\fld{gap} & -- & -- & -- & 1.0 & 最优间隙；无全局证书时保持 1.0（诚实口径）\\
\fld{nodes\_explored} & -- & 个 & -- & 0 & 已探索离散节点数\\
\fld{nlp\_solves} & -- & 次 & -- & 0 & 内层连续 OPF 求解次数\\
\fld{runtime\_sec} & -- & s & -- & 0 & 求解耗时\\
\fld{status} & -- & 串 & -- & 空 & 区分 time-limit / evaluation-limit / local optimum\\
\fld{algorithm} & -- & 串 & -- & 见默认 & ``\fld{discrete\_coordinate\_search\_with\_ac\_opf}''（:253）\\
\fld{globally\_certified} & -- & bool & -- & false & 恒 false：只保证可行局部最优（:254）\\
\fld{optimality\_gap\_available} & -- & bool & -- & false & 恒 false（:255）\\
\fld{terminated\_by\_time\_limit} & -- & bool & -- & false & 因时限终止（:256）\\
\fld{terminated\_by\_evaluation\_limit} & -- & bool & -- & false & 因评估数上限终止（:257）\\
\fld{model\_limitations} & -- & 串表 & -- & 空 & 本次求解能力边界声明（:258）\\
\fld{effective\_inner\_nlp\_objective} & -- & 枚举 & 两档 & MatchRPO & 实际生效的内层目标（MatchRPO / LossEconomic，:259--260）\\
\fld{vm\_before} / \fld{vm\_after} & $V_m$ & pu & -- & 空 & 优化前/后母线电压幅值\\
\fld{va\_before} / \fld{va\_after} & $\theta$ & rad & -- & 空 & 优化前/后母线相角（弧度）\\
\fld{pg\_mw\_before} / \fld{pg\_mw\_after} & $P_g$ & MW & -- & 空 & 优化前/后机组有功\\
\fld{qg\_mvar\_before} / \fld{qg\_mvar\_after} & $Q_g$ & Mvar & -- & 空 & 优化前/后机组无功\\
\fld{taps} & -- & 表 & -- & 空 & \texttt{vector<TapResult>}，OLTC 离散动作（见下表）\\
\fld{shunts} & -- & 表 & -- & 空 & \texttt{vector<ShuntResult>}，可投切并联动作（见下表）\\
\fld{total\_loss\_before\_mw} / \fld{total\_loss\_after\_mw} & -- & MW & -- & 0 & 优化前/后总网损\\
\fld{max\_vdev\_before} / \fld{max\_vdev\_after} & -- & pu & -- & 0 & 优化前/后最大电压偏差\\
\fld{bc\_stats} & -- & 结构 & -- & -- & 兼容信封：计局部评估/灵敏度平面数，\textbf{非} B\&B 全局 gap 证书（:284--286）\\
\fld{baseline\_opf} / \fld{optimized\_opf} & -- & 结构 & -- & -- & 优化前/后两个完整 \texttt{ACOPFResult}（:289--290），供独立 OPF/PF 回放交叉验证（见第~\ref{sec:validation}~章）\\
\end{paramtable}

\subsubsection*{TapResult（reactive\_power\_opt.hpp:158--167）}
\begin{paramtable}
\fld{trafo\_index} & -- & -- & -- & 0 & 变压器下标（\texttt{ac.transformers\_2w} 位置）\\
\fld{name} & -- & 串 & -- & 空 & 设备名\\
\fld{tap\_before} / \fld{tap\_after} & -- & 档 & -- & 0 & 优化前/后档位\\
\fld{ratio\_before} / \fld{ratio\_after} & -- & -- & -- & 1.0 & 优化前/后变比\\
\fld{electrical\_tap\_before} / \fld{electrical\_tap\_after} & -- & -- & -- & 1.0 & 优化前/后电气抽头值\\
\end{paramtable}

\subsubsection*{ShuntResult（reactive\_power\_opt.hpp:170--177）}
\begin{paramtable}
\fld{shunt\_index} & -- & -- & -- & 0 & 并联设备下标（\texttt{ac.shunts} 位置）\\
\fld{name} & -- & 串 & -- & 空 & 设备名\\
\fld{step\_before} / \fld{step\_after} & -- & 档 & -- & 0 & 优化前/后投切档数\\
\fld{bs\_mvar\_before} / \fld{bs\_mvar\_after} & $B_s$ & Mvar & -- & 0 & 优化前/后电纳（Mvar 口径）\\
\end{paramtable}

\subsection{IPMResult 结构}
定义 native\_ipm\_solver.hpp:34--59（与 parity\_ipm\_solver.hpp 共享）。
\begin{paramtable}
\fld{converged} & -- & bool & -- & false & 收敛标志\\
\fld{iterations} & -- & 次 & -- & 0 & IPM 迭代数\\
\fld{primal\_inf} & -- & -- & -- & 0 & 最终原始不可行度（赋值点 parity\_ipm.cpp:2132）\\
\fld{dual\_inf} & -- & -- & -- & 0 & 最终对偶不可行度（:2133）\\
\fld{complementarity} & -- & -- & -- & 0 & 最终互补间隙（:2134）\\
\fld{initial\_primal\_inf} / \fld{initial\_dual\_inf} / \fld{initial\_complementarity} & -- & -- & -- & 0 & 上述三项的初始值（:40--42）\\
\fld{warm\_start\_used} & -- & bool & -- & false & 是否用了暖启动\\
\fld{factorization\_calls} / \fld{linear\_solve\_calls} & -- & 次 & -- & 0 & 分解/回代次数\\
\fld{accepted\_steps} / \fld{rejected\_steps} & -- & 次 & -- & 0 & 被接受/拒绝步数\\
\fld{symbolic\_analyze\_calls} & -- & 次 & -- & 0 & 符号分析调用数（:48；赋值 parity\_ipm.cpp:2145）\\
\fld{backend\_escalations} / \fld{scaling\_rebuilds} & -- & 次 & -- & 0 & 后端升级 / 均衡重建次数（:49--50）\\
\fld{status} & -- & 串 & -- & 空 & 状态串\\
\fld{linear\_solver} & -- & 串 & -- & 空 & ``\fld{dense\_lu}''/\allowbreak``\fld{sparse\_umfpack}''/\allowbreak``\fld{sparse\_klu}''/\allowbreak``\fld{sparse\_eigen\_lu}''；MUMPS 后端经 default 分支报告为 ``sparse''（parity\_ipm.cpp:2136--2144）\\
\fld{x} & $x$ & -- & -- & 空 & 原始解向量\\
\fld{lambda\_eq} & $\lambda$ & -- & -- & 空 & 等式乘子（行序按 \texttt{Problem::cidx}）\\
\fld{mu} & $\mu$ & -- & -- & 空 & 不等式乘子（:57）\\
\fld{z} & $z$ & -- & -- & 空 & 界乘子\\
\end{paramtable}

\subsection{ThreePhaseHybridOPFResult 结构}
定义 three\_phase\_hybrid\_opf.hpp:138--184；填充 three\_phase\_hybrid\_opf.cpp:2710--2926：
\begin{itemize}
  \item \textbf{顶层}：\fld{converged}/\fld{status}/\fld{solver}/\fld{iterations}/
        \fld{objective}/\fld{runtime\_ms}（:139--142, :154--155）；
  \item \textbf{规模统计}：\fld{variables}/\fld{equalities}/\fld{inequalities}/
        \fld{enforced\_inequalities}（:143--146）、
        \fld{equality\_jacobian\_nonzeros}/\fld{inequality\_jacobian\_nonzeros}
        （:149--150）、\fld{eliminated\_phase\_nodes}（:153）；
  \item \textbf{初始诊断}：\fld{initial\_primal\_residual}/\fld{initial\_dual\_residual}/
        \fld{initial\_worst\_equality}/\fld{initial\_worst\_inequality}（:156--159）；
  \item \textbf{KKT 残差}：\fld{primal\_residual}/\fld{dual\_residual}/
        \fld{complementarity}（:160--162；取 unscaled 优先，cpp:2741--2747）；
  \item \textbf{导数校核}（\fld{verify\_derivatives}=true 时填充）：
        \fld{max\_equality\_jacobian\_error}/\fld{max\_inequality\_jacobian\_error}/
        \fld{max\_lagrangian\_hessian\_error}（:169--171）；
  \item \textbf{物理指标}：\fld{max\_voltage\_violation}、\fld{max\_vuf}、
        \fld{max\_converter\_current\_vuf}、\fld{max\_converter\_current\_loading}、
        \fld{max\_converter\_violation}、\fld{max\_dynamic\_equilibrium\_residual}
        （:163--168）、\fld{max\_omitted\_inequality}（:151--152；默认 $-\infty$，
        oracle 省略行的最大违反，诚实口径）；
  \item \textbf{解}：\fld{primal}/\fld{equality\_dual}/\fld{inequality\_dual}/
        \fld{inequality\_slack}/\fld{full\_voltage}/\fld{dc\_voltage}/
        \fld{converter\_phase\_power\_pu}/\fld{converter\_dc\_power\_pu}（:172--179）；
        对偶扩展：enforced 子集 expand 回 nineq 长度，box 乘子接尾部（cpp:2750--2777）；
  \item \textbf{oracle 与降阶}：\fld{constraint\_oracle\_rounds}/
        \fld{constraint\_oracle\_added\_rows}（:147--148）、
        \fld{constraint\_oracle\_trace}（:181）、\fld{enforced\_inequality\_rows}
        （:180）、\fld{reduction} 证书（:183）、\fld{converter\_dynamic\_equilibria}
        （:182）。
\end{itemize}

\subsection{OPFProblem 与 IOPFSolver 求解器抽象}
\srcpath{opf\_problem.hpp:15--22} 的 \texttt{OPFProblem} 是``做什么/怎么解''解耦
描述器（sys 指针 + \fld{ac\_opts}/\fld{dc\_opts} + \fld{is\_ac}/\fld{is\_dc}）。
消费方已实现：\srcpath{opf\_solver\_interface.hpp} 定义
\fld{OPFSolveResult} = \texttt{std::variant<\allowbreak ACOPFResult,\allowbreak DCOPFResult>}
（:22）、抽象类 \texttt{IOPFSolver}（:24--28）与
\texttt{DefaultOPFSolver} \texttt{final}（:30--42）——\fld{problem.sys} 为空抛
\texttt{std::invalid\_argument}（:33--34）；\fld{is\_ac}/\fld{is\_dc} 必须恰选一个，
否则同样抛异常（:35--37）；按域转发到生产入口 \srcpath{solve\_ac\_opf}/
\srcpath{solve\_dc\_opf}（:38--40），不引入第二条建模路径。
\textbf{当前口径}：已有实现且被测试消费（test\_hacdcdss.cpp:420，断言 DC 变体
\texttt{holds\_alternative<DCOPFResult>}），但生产 \texttt{src/} 内无调用方；
该头文件已纳入公共门面 \srcpath{include/hacdcpf/api/hacdcpf.hpp:12}。

\subsection{诊断与环境变量}
\begin{itemize}
  \item \srcpath{check\_ac\_opf\_jacobian\_fd}（ac\_opf.cpp:3826--4024）：
        解析 Jacobian vs 中心差分，判据 $\max|\mathrm{err}|\le5\times10^{-4}$ 或
        相对误差 $\le5\times10^{-3}$；\textbf{注意}：头文件声明的公开入口
        \texttt{compute\_jacobian\_diagnostics} 无实现，本函数目前无调用方（死代码）；
  \item \srcpath{verify\_opf\_result}（src/api/hacdcpf.cpp:1348--1451）：
        独立审计电压限、机组限、目标值重算，生成 infeasibility hints；
  \item \textbf{环境变量全集}（均为调试/对照开关，非生产配置；已与
        \texttt{src/optimal\_power\_flow/} 下全部 \texttt{getenv} 调用逐一核对）：
        \begin{itemize}\footnotesize
        \item \texttt{HACDCPF\_OPF\_KKT\_FORM}：condensed（默认）/ augmented
              （\srcpath{kkt\_form\_preference}，parity\_ipm.cpp:188--199）；
        \item \texttt{HACDCPF\_OPF\_LINEAR\_SOLVER}：klu / umfpack / eigen / dense / mumps，
              钉死线性后端（\srcpath{linear\_solver\_preference\_env}，:71--85）；
        \item \texttt{HACDCPF\_OPF\_SPARSE\_SOLVER}：上一项的旧别名（:76）；
        \item \texttt{HACDCPF\_OPF\_MU\_STRATEGY}：mehrotra（默认）/ abo（:1097--1104）；
        \item \texttt{HACDCPF\_OPF\_FILTER\_ALL}：滤波线搜索不限 KKT 维数启用（:1647--1653）；
        \item \texttt{HACDCPF\_OPF\_INERTIA\_GATE}=N（默认 0）：周期 MUMPS 惯性复检，
              复检后总恢复 KLU（:1244--1259）；
        \item \texttt{HACDCPF\_OPF\_DELTA\_C}：augmented $(2,2)$ 块 $\delta_C$ 绝对值覆盖（:1233）；
        \item \texttt{HACDCPF\_OPF\_RESCALE\_PER\_ITER}=1：逐迭代重建 Ruiz 均衡
              （默认冻结，:505）；
        \item \texttt{HACDCPF\_OPF\_STRICT\_THETA}=1：诊断用严格 $\theta$ 滤波（:1655）；
        \item \texttt{HACDCPF\_OPF\_NO\_LEGACY}=1：禁用 legacy 三轴接受路径（:1657）；
        \item \texttt{HACDCPF\_OPF\_PROF}=1：进程级累计 profiler，退出时向 stderr 打印
              factorize/solve/assembly 与分后端分解（:126）；
        \item \texttt{HACDCPF\_RPO\_PROF}=1：RPO 相位计时（reactive\_power\_opt.cpp:44）；
        \item \texttt{HACDCPF\_DISABLE\_IPOPT\_OPF}：禁用 Ipopt 兜底腿（ac\_opf.cpp:2119）。
        \end{itemize}
        交叉参照：HTTP 自省端点 \srcpath{opf\_debug\_environment\_json}
        （run\_gui\_server.cpp:6818）枚举其中 10 个 \texttt{HACDCPF\_OPF\_*} 变量的
        当前取值，返回标注 \texttt{debug\_only\_not\_stable\_api}（不含
        \texttt{HACDCPF\_OPF\_PROF}，以及非 \texttt{OPF\_} 前缀的
        \texttt{HACDCPF\_RPO\_PROF} 与 \texttt{HACDCPF\_DISABLE\_IPOPT\_OPF}）。
\end{itemize}

% 第 10 章：OPF 在平台中的应用与接入
\section{OPF 在平台中的应用与接入}\label{sec:applications}

本章盘点调用 OPF 求解器的下游模块与外部接入层（GUI/HTTP、Python 绑定）。
各模块自身的数学模型见其专章/文档，此处只记\textbf{OPF 调用配置}。

\subsection{时序流水线 UC$\to$OPF$\to$PF}
\srcpath{src/time\_series/time\_series\_pf.cpp:4218} 每时步调用
\srcpath{opf::solve\_ac\_opf(sys\_t, step\_opf\_options)}，选项来自
\fld{TimeSeriesPFOptions.opf\_options}（time\_series\_pf.hpp:105），\fld{run\_opf}
开关（:120）。
\begin{itemize}
  \item \textbf{逐时步暖启动链}：上一时步收敛则 \fld{step\_opf\_options.warm\_start =
        \&previous}（:4214--4217）——这是 \fld{ACOPFOptions::warm\_start} 唯一的
        生产级用户；
  \item OPF 不收敛时回退为 UC 出力直接跑 PF（:4224--4235）；
  \item 收敛时把 \fld{pg/qg}、VSC \fld{pac/qac}（置 PQ 模式）、可再生 \fld{pren}
        （经 \fld{ren\_map} 归因）注入系统再跑 PF 校验（:4238--4275）；
        储能保留跨时段调度不回放（:4277 注释）。
\end{itemize}

\subsection{可靠性评估的 DC OPF 配置}
\srcpath{src/reliability/reliability\_assessment.cpp:1057}（\srcpath{evaluate\_state}，
被 MC/FMEA/时序可靠性批量复用，另见 :4268--4275、:4418--4422）调用
\srcpath{opf::solve\_dc\_opf}：
\begin{itemize}
  \item 强制 \fld{load\_shedding}=true、\fld{voll}=\srcpath{reliability\_shedding\_voll}
        （VOLL 压过发电成本，字典序最小切负荷）、\fld{verbose}=false、
        \fld{compute\_lmp}=\textbf{false}（批量路径跳过支撑 LP，:1224--1228、:1701--1705）；
  \item 死岛负荷直接计切负荷，与 OPF \fld{total\_load\_shedding\_mw} 合并为
        EENS/LOLE（:1089--1099）；
  \item \textbf{OPF 不收敛保守处理为存活负荷全切}（:1062--1087）；
  \item 含 DC 负荷时 warn ``AC-only DC-OPF，EENS 偏乐观''（:1052--1054）——
        诚实口径的典型实例。
\end{itemize}

\subsection{EV-交通耦合的 LMP 价格迭代}
\srcpath{src/ev\_power\_traffic/joint\_social\_welfare.cpp:262} 与
\srcpath{ctm\_joint\_welfare.cpp:254}：每时步 \srcpath{opf::solve\_dc\_opf(sys\_k, dcopf)}，
\fld{compute\_lmp}=true、\fld{load\_shedding}=true（:259--261 / :251--253）。
取每时步 bus$\to$LMP 表，按
\[
\pi_{new}=\alpha\,\lambda/1000+(1-\alpha)\,\pi_{old}
\]
更新充电站电价，构成 CTM-DUE/LP 分配 $\leftrightarrow$ DC OPF 的定点价格迭代
（:273--290 / :282--289）；OPF 失败则保留旧价（\fld{all\_opf\_ok}=false）。
\fld{use\_dcopf\_prices}=false 时退化为固定价 + 边际成本估成本。

\subsection{反事实规划与 SPPT 投影证书}
\begin{itemize}
  \item \textbf{反事实规划}（\srcpath{src/analysis/counterfactual\_planning.cpp:144}）：
        \srcpath{hacdcpf::solve\_ac\_opf}，\fld{max\_inner\_iterations}=120、
        \fld{allow\_fallback}=true（:141--143）；收敛且目标有限时
        $\fld{objective}\times\fld{annual\_hours}$ 作为年化经济基线（:145--148），
        否则回退 ``PF 物理损耗 $\times$ 电价 $\times$ 年小时'' 货币化损耗并把口径写入
        \fld{diagnostics[0]}（:149--161）；
  \item \textbf{SPPT MR3d}（\srcpath{src/sppt/metamorphic.cpp:223--225}）：对 rich 系统与
        \srcpath{project\_to\_canonical\_models(sys)} 各跑一次 DC OPF，按 bus index 对齐
        比较两次求解的 LMP——``投影前后 DC OPF 对偶不变性''的可执行证书；
        残差超 tol 或任一侧不收敛即判 fail。
\end{itemize}

\subsection{与市场模块的边界声明}
市场模块（\srcpath{src/market/market\_simulation.cpp}）的 SCUC/SCED/LMP 是
\textbf{自建优化，不复用} \srcpath{opf::solve\_dc\_opf}：
SCUC MILP 经 \srcpath{build\_market\_commitment\_model} + MIPSolvers 引擎
\srcpath{solve\_milp}（:675--705）；SCED 定价 LP 用
\srcpath{engine::solve\_lp\_with\_basis}（:1897、:2017），
LMP 直接取 \fld{solve.constraint\_duals[balance\_row]}/dt（:1248--1251）；
AC-only 守卫：含 DC 母线直接拒绝（:175）；AC 校验用 \srcpath{solve\_power\_flow}。
因此市场 LMP 与 DC OPF LMP \textbf{不同源}，口径不可互换。

\subsection{GUI/HTTP 参数映射}
统一路由 \texttt{POST} \srcpath{/api/session/opf}（tests/run\_gui\_server.cpp:14533）
是 GUI 与 Python 客户端的主入口。

\paragraph{顶层字段}
\begin{center}
\footnotesize
\begin{tabular}{@{}L{3.4cm}L{5.6cm}L{5.4cm}@{}}
\toprule
请求字段 & 取值 / 默认 & 去向\\
\midrule
\fld{solver} & parity（默认）/ ipopt / auto / robust / dispatch / dc &
  backend 选择；robust 触发策略（见下）\\
\fld{network\_model} & balanced\_aggregate（默认）/ balanced\_with\_three\_phase\_validation
  / three\_phase\_hybrid（:14554--14560） & 三档网络模型；第三档禁 dc/dispatch（:14609--14612）\\
\fld{check\_consistency} & false & OPF 点 PF 回放审计（:14562）\\
\fld{constraints.*} & \fld{branch\_limits}/\fld{converter\_capacity}/\fld{converter\_current}/
  \fld{converter\_modulation} 全默认 true（:14564--14567） & \fld{enforce\_*} 四开关
  （:14989--14992）；dc 分支只接 \fld{branch\_limits}（:14872）\\
\bottomrule
\end{tabular}
\end{center}

\paragraph{solver 字符串 $\to$ backend（:14916--14947）}
\texttt{parity}$\to$ParityIPM；\texttt{ipopt}$\to$Ipopt（强制
\fld{objective\_homotopy}=false，:14993--14998）；\texttt{dispatch}$\to$EconomicDispatch；
\texttt{auto}$\to$Auto + primal\_dual + parity + \fld{allow\_fallback}，
混合交直流时 \fld{ac\_pf\_warm\_start}=true；
\texttt{robust}$\to$Auto + \fld{ac\_pf\_warm\_start} + \fld{objective\_homotopy}，
且请求级开关\textbf{不可关闭}（:14968--14974）。

\paragraph{AC \fld{options.*}（:14949--14987，默认来自 ACOPFOptions）}
\fld{max\_inner\_iterations}（非线性 400 / dispatch 120，clamp 1--$10^5$）、
\fld{max\_outer\_iterations}（8，1--$10^4$）、\fld{max\_line\_search\_steps}（20）、
\fld{feasibility\_tol}/\fld{stationarity\_tol}（$10^{-6}$）、\fld{barrier\_mu0}（$10^{-2}$）、
\fld{barrier\_mu\_reduction}（0.2）、\fld{barrier\_mu\_min}（$10^{-8}$）、
\fld{merit\_penalty}（10）、\fld{regularization}（$10^{-6}$）、\fld{step\_backoff}（0.5）、
\fld{interior\_fraction}（0.995）、\fld{ac\_eval\_threads}（1）、\fld{allow\_fallback}（true）、
\fld{ac\_pf\_warm\_start}（false）、\fld{objective\_homotopy}（false）、
\fld{homotopy\_dt0}（0.10）、\fld{verbose}。

\paragraph{DC \fld{options.*}（:14870--14885）}
\fld{feasibility\_tol}（$10^{-6}$）、\fld{max\_iterations}（10000，1--$10^6$）、
\fld{pwl\_segments}（4，1--1000）、\fld{branch\_limit\_margin}（1.0，$10^{-6}$--10）、
\fld{load\_shedding}（true）、\fld{voll}（0.0$\to\ge0$）、\fld{compute\_lmp}（true）、
\fld{verbose}。

\paragraph{\fld{three\_phase.*}（:14613--14645）}
\fld{include\_shunts}（true）、\fld{vuf\_max}（0.03，0--1）、
\fld{variant}（graph\_reduced / full）、\fld{verify\_derivatives}（false）、
\fld{constraint\_oracle}（false）、\fld{oracle\_max\_rounds}（8）、
\fld{reduction\_max\_front}（64）、\fld{reduction\_max\_nnz\_ratio}（10.0）；
auto/robust 失败时 Ipopt fallback（:14650--14656），robust 时
\fld{warm\_start\_with\_ipopt}=true（:14629）。

\paragraph{RPO 路由 \srcpath{/api/session/run\_rpo}（:18406--18448）}
\fld{objective}（voltage / loss / combined）、\fld{vdev\_weight}/\fld{loss\_weight}（1.0）、
\fld{v\_target}（1.0）、\fld{mip\_gap}（0.01）、\fld{time\_limit\_s}（120）、
\fld{max\_evaluations}（50000）、\fld{max\_ipm\_iter}（2000）、\fld{ipm\_tol}（$10^{-6}$）、
\fld{stationarity\_tol}（$10^{-3}$）、\fld{max\_tap\_move}（2，-1--1000）、
\fld{restrict\_tap\_indices}/\fld{enabled\_tap\_indices}、四个约束开关、
\fld{relax\_only}（仅根松弛，max\_nodes=1）。
解后用 \srcpath{apply\_rpo\_discrete\_solution} 重建系统并跑独立
\srcpath{solve\_ac\_opf} 交叉验证（:18453--18500）。

\paragraph{其他路由与 v1 差异}
\srcpath{/api/session/opf\_ac}（:15633，solver 默认 auto，max\_inner=120）、
\srcpath{/api/session/opf\_parity}（:15709，固定 parity 禁 fallback）、
\srcpath{/api/session/opf\_dc}（:15763）、\srcpath{/api/session/rpo\_inputs}（:18358）、
遗留 \srcpath{/api/opf/ac|dc}（:22621/22641）。
v1 作业 API（\srcpath{src/server/runtime\_api\_v1.cpp:825--840}）：
\fld{network\_model} 仅支持 \fld{balanced\_aggregate}（:826--832），
solver 支持 parity / ipopt / auto / dispatch / dc（:751--766）；
序列化含 \fld{solver\_backend}/\fld{solver\_chain}/\fld{objective\_model}/
\fld{branch\_mu\_valid}/\fld{model\_scope}/\fld{validity\_flags}（:725--745）。

\paragraph{参数契约与自省端点（2026-07-22 新增）}
为消除``前端控件 $\leftrightarrow$ 请求字段 $\leftrightarrow$ 后端语义''三方口径漂移，
服务端把 OPF 参数元数据做成了机器可读的单一事实来源：
\begin{itemize}
  \item \srcpath{opf\_parameter\_contract\_json()}（\srcpath{tests/run\_gui\_server.cpp:6740--6816}）：
        返回 \fld{schema}=\texttt{opf\_parameter\_contract\_v1}，逐项列出 32 个参数的
        \fld{key}、\fld{label}（中文标签）、\fld{unit}、\fld{default}、
        \fld{applies\_to}（ac / dc / phase / all 适用路径）与 \fld{description}
        （用途与影响说明，含``QP 后端 \fld{pwl\_segments} 生效值为 0''一类口径注记）；
  \item \srcpath{opf\_debug\_environment\_json()}（:6818--6832）：枚举 10 个
        \texttt{HACDCPF\_OPF\_*} 调试环境变量（KKT\_FORM、LINEAR\_SOLVER、
        SPARSE\_SOLVER、MU\_STRATEGY、INERTIA\_GATE、RESCALE\_PER\_ITER、DELTA\_C、
        FILTER\_ALL、STRICT\_THETA、NO\_LEGACY），只回显当前已设置者，
        接口等级显式标注 \fld{contract}=\texttt{debug\_only\_not\_stable\_api}；
  \item 新路由 \texttt{GET} \srcpath{/api/opf/parameter\_contract}（:14769--14771），
        无会话状态要求，供前端与外部客户端在请求前自省参数语义；
  \item 统一 OPF 端点响应新增三个公共字段（:14838--14840）：
        \fld{options\_requested}（请求 \fld{options} 原文回显）、
        \fld{parameter\_contract\_schema}（契约模式版本）、
        \fld{debug\_environment}（本次求解生效的调试环境快照）。
\end{itemize}

\paragraph{结果回显增强（诚实口径的机器可读化）}
各 OPF 分支在响应中追加``结果能解释到什么程度''的显式声明：
\begin{itemize}
  \item \textbf{DC OPF 分支}（:15154--15168）：\fld{lmp\_valid} /
        \fld{lmp\_validity\_reason}、\fld{branch\_mu\_valid} /
        \fld{branch\_mu\_validity\_reason}（对应支路对偶未认证的既有声明，见第~1~章）、
        \fld{objective\_model}、\fld{pwl\_segments\_effective}；
        \fld{options\_effective} 中区分 \fld{pwl\_segments\_requested} 与生效值；
        附 \fld{model\_limitations} 与 \fld{scope} 块（DC 路径换流器四族约束
        恒 false）；
  \item \textbf{AC OPF 分支}（:15311--15313）：\fld{lmp\_valid}(\fld{\_reason})
        与 \fld{model\_limitations}；
  \item \textbf{三相混合 OPF 分支}（:15049--15054）：\fld{lmp\_valid}=false 并注明
        该路径不导出约束乘子作为 LMP；\fld{model\_limitations} 声明不执行
        AC 支路热限与换流器调制约束；
  \item \textbf{RPO 端点}（:19118--19121）：\fld{terminated\_by\_time\_limit}、
        \fld{terminated\_by\_evaluation\_limit} 与 \fld{model\_limitations}，
        把``达到时限/评估上限''从日志升格为结构化字段。
\end{itemize}

\paragraph{前端映射与结果呈现（web/js/app.js、web/index.html）}
\begin{itemize}
  \item \fld{OPF\_PARAMETER\_INPUTS} 映射表（app.js:4398--4416）：契约参数键
        $\to$ DOM 控件 id；\srcpath{loadOpfParameterContract()}（:4452--4468）
        拉取契约后为控件挂 tooltip（描述 + 单位）并渲染可检索的参数说明表；
        运行 OPF 前自动加载（:4494），参数指南面板打开时亦加载（:16425）；
  \item \textbf{DC 求解器专属请求参数}（app.js:4522--4528）：新增 \fld{opfDcOptions}
        控件组（index.html:506--513，仅 solver=dc 时显示，app.js:11830），
        对应请求字段如下表。
\end{itemize}
\begin{paramtable}
\fld{pwl\_segments} & --- & 段/机组 & 1--1000 & 4 &
  二次成本分段线性区间数（LP 回退后端；QP 后端生效值为 0）\\
\fld{branch\_limit\_margin} & --- & 倍 & $10^{-6}$--10 & 1.0 &
  DC 支路 rate\_a 统一倍率，$<1$ 收紧、$>1$ 放宽\\
\fld{load\_shedding} & --- & 开关 & --- & true &
  为正负荷加削减变量保可行，按 VOLL 计入目标\\
\fld{voll} & --- & 成本/MWh & $\ge 0$ & 0.0 &
  切负荷惩罚；0 表示按最大发电边际成本自动生成\\
\fld{compute\_lmp} & --- & 开关 & --- & true &
  跑支撑 LP 恢复节点功率平衡对偶；不代表支路拥塞 $\mu$ 已认证\\
\end{paramtable}
\begin{itemize}
  \item \textbf{结果渲染}：scope 表新增``节点电价乘子''行（认证 / 未认证 + 原因，
        app.js:4777）与 DC 专属``支路拥塞 $\mu$''行（:4778）；
        \fld{model\_limitations} 渲染为``模型边界''表（:4780--4783）；
        调试环境变量表（:4784--4788，标注``仅调试，不是稳定接口''）；
        新增 \fld{opfParameterResults}``请求值 vs 实际值''逐参数核对表
        （:4791--4840），判定列三档：按请求生效 / 后端调整或路径决定 / 未回显；
  \item DC / 三相混合模式下隐藏不适用的 AC 专属控件（:11828--11835：
        barrier 参数组、外迭代、线搜索、驻点容差），避免``可调但不生效''的误导；
  \item RPO 结果新增``搜索终止''显示（:5297：达到时限 / 达到评估上限 /
        局部邻域完成）与``能力边界''列表（:5304，逐条展示 \fld{model\_limitations}）。
\end{itemize}

\subsection{Python 绑定与能力查询}
\begin{itemize}
  \item 模型层 \srcpath{python/src/hysim/models.py:130--183}：\texttt{OPFSolver}
        （parity / ipopt / auto / dc / dispatch）、\texttt{OPFNetworkModel} 三档枚举、
        \texttt{OPFConstraints} 四开关（全默认 true）、\texttt{OPFRequest}
        （含 \fld{three\_phase}、\fld{check\_consistency}）；
  \item 客户端 \srcpath{client.py:156} \texttt{optimal\_power\_flow()} 映射到
        \srcpath{/api/session/opf}；v1 会话 \srcpath{v1.py:340--349}；
        AI 工具 \srcpath{ai.py:163}（\texttt{hysim\_run\_optimal\_power\_flow}）。
        Python 侧无独立求解，纯 HTTP 包装；
  \item 能力查询 \srcpath{include/hacdcpf/api/solver\_capabilities.hpp}：
        \fld{supports\_ac\_opf}/\fld{supports\_dc\_opf} 恒 true；
        \fld{has\_ipopt}（决定 solver=ipopt 与 auto fallback 可用性）、
        \fld{has\_highs}/\fld{has\_gurobi}/\fld{has\_scip}（影响 DC OPF LP 后端与
        RPO 内层）、\fld{has\_native\_ipm} 恒 true、\fld{supports\_quadratic\_objective}；
        运行时入口 \srcpath{get\_solver\_capabilities()}。
\end{itemize}

% 第 11 章：验证与校核总览
\section{验证与校核总览}\label{sec:validation}

\subsection{三层验证体系}
\begin{enumerate}
  \item \textbf{手算小算例}：\srcpath{external\_data/opf\_hand\_cases/} 4 个带解析解的
        DC OPF 案例（README 给出期望调度/目标/LMP）：
        01 单母线 merit-order（$P_g=[40,20]$，obj=800，LMP$\approx$20）；
        02 双母线阻塞（$P_g=[40,60]$，obj=3400，$\theta_2=-0.04$~rad，LMP=[10,50]，
        明确支路对偶未认证）；03 切负荷（VOLL 自动=100，切 30~MW，obj=4000）；
        04 二次成本 LP 线性化（$c^{lin}=c_1+2c_2P^{max}$，期望 $P_g=[80,0]$、obj=1120；
        真 QP 参考 obj$\approx$815.333）。\textbf{仅供 GUI 手工核对，未接入自动回归。}
  \item \textbf{MATPOWER 标准算例回归}：\srcpath{external\_data/matpower/} 共 85 个
        \texttt{.m} 文件（经 \texttt{ls} 核实），其中 \texttt{contab\_*}（4 个）与
        \texttt{scenarios\_*}（2 个）为扰动/场景辅助文件，标准算例 79 个
        （case4gs 至 case13659pegase、ACTIVSg 系列、RTE/Polish 系列），
        由 test\_opf\_solver\_backends 与 test\_acopf\_dcopf\_crossval 驱动。
        79 个标准算例按规模粗分三档：\textbf{小规模}（$\le$500 节点，50 个，
        case4gs--case89pegase 及 case300、ACTIVSg200/500，含 IEEE 标准系与
        配网等值 case33bw/case141/case\_ieee123、RTS-GMLC）；\textbf{中规模}
        （501--4000 节点，18 个，case533mt\_hi/lo 至 case3375wp，含 RTE
        1888/1951/2848/2868、Polish 2383wp--3375wp、pegase 1354/2869、
        ACTIVSg2000）；\textbf{大规模}（$>$4000 节点，11 个，case6468rte--%
        case13659pegase、ACTIVSg10k/25k/70k、case\_SyntheticUSA）。
  \item \textbf{大系统鲁棒性}：case2869pegase/case6515rte/case2000\_acdc，
        稀疏 KKT + filter 回溯（\srcpath{docs/large\_hybrid\_ipm\_diagnostics.md}）。
\end{enumerate}

\subsection{测试矩阵}
\begin{footnotesize}
\begin{longtable}{@{}L{4.2cm}L{3.6cm}L{6.6cm}@{}}
\toprule
测试目标 & 算例 & 关键断言（容差）\\
\midrule
\endhead
test\_opf\_solver\_backends & case300\_acdc, case2000\_acdc, case2383wp, case30,
  case9, case2869pegase, case6515rte, 混合微网 &
  parity 收敛；暖启动目标差 $\le10^{-5}$；DC 参考 $|v_{dc}-v_{set}|<2\times10^{-6}$；
  KCL Jacobian FD $<10^{-6}$；case2000 $<160$ 迭代；case2869 $<100$ 迭代稀疏后端；
  case6515rte 不发散；Parity 成本 $\le$ ED$\times1.001+1$\\
test\_acopf\_dcopf\_crossval & MATPOWER case5/9/14/30/39/57 &
  ACOPF/DCOPF 调度重跑 PF 收敛（case57 WARN）；孤岛切负荷逐母线 margin $10^{-8}$；
  LMP $\ge-10^{-4}$ 且 $\le10\times$maxMC$+1$；无阻塞 LMP spread $<10^{-4}$；
  阻塞 $|flow|\le$限$\times1.01+0.1$；\fld{branch\_mu\_valid}=false 显式断言，
  且 \fld{branch\_mu\_validity\_reason} 非空、\fld{lmp\_valid}=true
  （:870--872）\\
test\_reactive\_power\_opt & case9, case300 &
  \fld{globally\_certified}=false、\fld{gap}=1.0；\fld{model\_limitations} 非空
  （:33）；目标不劣于基线 $+10^{-10}$；
  排除原因枚举；固定变比不伪造 OLTC；tap 联动 $1.03\times1.025$\\
test\_three\_phase\_hybrid\_opf & 3 母线$\times$3 相 + 2 DC 手工算例 &
  Full vs GraphReduced 目标差 $\le10^{-6}$；动态平衡残差 $<10^{-7}$；
  导数误差 $<10^{-5}$；独立 PF 复现 $<2\times10^{-6}$\\
test\_nighttime\_opf & Canvas 系统、南沙全网 &
  ng=0 时 external grid 提升；储能钉调度 WithinAbs $2\times10^{-4}$；
  DC/DC $I^2R$ 复核；24 步 UC$\to$OPF$\to$PF 全收敛\\
test\_converter\_coordination & 构造混合算例 &
  AML Vdc 松弛/AC 参考与控制角色解析器一致；DC/DC 三模式
  margin $10^{-3}$--$2\times10^{-2}$；DC/DC 数据捕获 \fld{r\_eq\_pu}=0.01
  （:1541,1555）、Power 模式 \fld{max\_dc\_p\_viol\_pu}$<10^{-4}$（:1575）\\
test\_market\_simulation & case9 24h &
  SCUC$\to$SCED/LMP$\to$N-1$\to$结算闭环；AC 认证 residual $<10^{-5}$；
  结算现金流残差 $\approx0$（margin $10^{-6}$）\\
test\_hacdcdss & 自建 2/3 母线、case9 &
  DC OPF 5 用例（:315--443，6 处 \srcpath{solve\_dc\_opf} 调用）：
  廉价机组优先调度；Native LP 目标与 LMP 为 MW 标度、\fld{lmp\_valid}=true、
  \fld{branch\_mu\_valid}=false 且 reason 非空（:350--352）；
  \fld{pwl\_segments} 分段数控制 LP 二次成本逼近；单机带载在限值内；
  case9 经济调度可行；\texttt{OPFProblem}$\to$\texttt{DefaultOPFSolver}
  接线用例（:408--424）\\
test\_crossmodule\_integration & IEEE14、case33bw &
  Pipeline 1：AC/DC OPF$\to$Newton PF$\to$碳流（:142 起），DC OPF 成本
  $\le$ AC OPF$\times1.05$（松弛性）、$|V_m^{PF}-V_m^{OPF}|<0.05$~pu；
  Pipeline 2：ONR$\to$DC OPF$\to$PF$\to$碳流（:283 起），重构后 DC OPF 成本
  $\le$基准$\times1.05$、调度在 $[P_{min},P_{max}]$\\
test\_graph\_roundtrip & 自建链式/环网、case33bw &
  AC OPF 结果经重索引反投影尺寸守恒（:223--258）；DC OPF 目标在 series
  降阶下保持 WithinAbs $10^{-3}$（:671--704）\\
test\_distribution\_pipeline & 自建 CB/开关/VSC 配网 &
  rich$\to$graph$\to$PF$\to$OPF 端到端管线（:410 起）；Stage 6 DC OPF
  DER 优先调度（:637--667），$P_g$ 在限值内、$P_g[1]\approx3.0$
  （margin 0.1）\\
test\_multiscale\_comprehensive & 多尺度 AC/DC 综合算例 &
  ParityIPM 受限 OPF 收敛（:170 起），AC/DC 电压范围有限、诊断带
  \texttt{opf\_} 前缀；毫秒至年度多尺度链（:231 起）\\
test\_solver\_capabilities & case9 等 &
  \fld{supports\_ac\_opf}/\fld{supports\_dc\_opf} 能力标志与实际求解一致
  （:135--158）；\srcpath{verify\_opf\_result} 独立审计通过/不收敛提示
  （:194 起）；混合算例抑制 AC-only 回退（:241 起）\\
test\_matpower\_cases & MATPOWER 三档 + case9/14 &
  Tier1 九算例 PF 收敛与电压合理性（另驱动潮流手册三档回归）；
  DC OPF 一致性：case9 目标有限（:618 起）、case14 调度在机组限值内
  （:626 起）\\
rpo\_gui\_e2e（E2E） & case24、case300（浏览器） &
  ctest 目标 \texttt{rpo\_gui\_e2e}（\srcpath{tests/CMakeLists.txt:502--507}，
  标签 \texttt{gui;e2e;browser;rpo;cross\_validation}，超时 180~s，Node +
  Playwright chromium 门控）：RPO 模块在稳态工作流暴露、
  \texttt{/api/session/run\_rpo} 响应契约、结果面板渲染（局部可行解/
  电压越限母线/覆盖范围/独立 OPF 重算）、case300 RPO 输入清单与 OLTC
  档位定位\\
\bottomrule
\end{longtable}
\end{footnotesize}

\subsection{交叉验证机制}
\begin{itemize}
  \item \textbf{OPF$\to$PF 回放}：AC OPF/DC OPF/三相 OPF 的调度与运行点均由独立潮流
        复现（三相电压差 $<2\times10^{-6}$）；
  \item \textbf{RPO 交叉验证}（\srcpath{docs/reactive\_power\_optimization\_validation.md}）：
        独立 OPF 复算 + PF 回放，判据目标相对差 $\le\max(10^{-3},10\epsilon_d)$、
        电压差 $\le\max(5\times10^{-4},10\epsilon_p)$~pu；富混合 AC/DC 标注“部分覆盖”；
  \item \textbf{解析导数校核}：\srcpath{check\_ac\_opf\_jacobian\_fd}（死代码，见
        第~\ref{sec:acopf}~章注）、
        DC slack KCL FD（$<10^{-6}$）、三相 Jacobian/Hessian（$<10^{-5}$，
        另用 dynamics 设备模型独立复核 $<10^{-7}$）；
  \item \textbf{独立审计}：\srcpath{verify\_opf\_result} 重算限值与目标；
  \item \textbf{能力打点}：\srcpath{tools/validate\_case\_capabilities.py}
        （OPF 收敛+成本$>$0+出力在限值内；RPO 动作非空）与
        \srcpath{tools/gui\_api\_e2e.py}（\texttt{/api/session/opf} 回归、OPF$\to$PF 一致性）。
\end{itemize}

\subsection{性能实测锚点（代码注释与文档）}
case2869pegase 经济目标精确参考值 133999（同伦 154$\to$73 迭代）；
case9241 参考目标 315837（$-0.013\%$）；case13659 经 \fld{ac\_pf\_warm\_start}
189 迭代/$\sim$40~s；case2000\_acdc $\sim$2~s；
IEEE24-3area-expanded 内层 IPM 640$\to$36 迭代（冻结 Ruiz 缩放）；
RPO 大规模（case2869+20 OLTC+10 并联）基线 74 迭代/25~s，暖启动链均值 0.46~s/次。

\subsection{已知边界与覆盖缺口}
\begin{enumerate}
  \item DC OPF \fld{branch\_mu\_valid} 恒 false（支路对偶未认证）；\fld{pwl\_segments}
        未接入构建；
  \item RPO 无全局最优证书；3000+ 母线完整搜索受时限约束；LossEconomic 内层
        不直接为电压最优；
  \item Native AC 旧路径用外部二次罚处理限值、不钉 DC\_V 电压、支路 $\pi$ 模型不含 shift；
  \item 储能无跨时段 SOC 模型（钉 UC 调度值）；能量路由器无损守恒；
  \item AML 层非生产链路；AC OPF/DC OPF/LinDistFlow/SCUC 的 AML 版本无直接测试；
        AML 混合 OPF 省略 DC/DC $I^2R$、\fld{max\_dc\_p\_viol\_pu} 未计 DC/DC 项；
        AML SCUC 无网络约束、MILP 下价格对偶通常不可得；
        LinDistFlow 头注释含 DG 而实现无 DG；
  \item Ipopt 路径无约束乘子 $\Rightarrow$ 无 LMP；
  \item \texttt{compute\_jacobian\_diagnostics} 声明无实现（实际 FD 诊断
        \srcpath{check\_ac\_opf\_jacobian\_fd} 为死代码）；
        \srcpath{external\_data/opf\_hand\_cases/} 未接入自动回归；
        \texttt{IOPFSolver} 抽象已有默认实现 \texttt{DefaultOPFSolver}
        （\srcpath{include/hacdcpf/optimal\_power\_flow/opf\_solver\_interface.hpp:30--42}，
        按 \fld{is\_ac}/\fld{is\_dc} 域转发生产 \srcpath{solve\_ac\_opf}/%
        \srcpath{solve\_dc\_opf}），并被测试消费
        （\srcpath{tests/test\_hacdcdss.cpp:408--424}），但生产 \texttt{src/}
        仍无调用方，属测试侧接线；
  \item 三相混合 OPF 活跃研发：AC 支路热限与 VSC 调制比未进入相域模型；
        Delta/ZIP 负荷、DCDC、能量路由器显式拒绝；
  \item 环境变量开关（KKT 形式、线性后端、$\mu$ 策略、inertia gate）效果因算例而异，
        属调试通道而非承诺接口。
\end{enumerate}

\subsection{如何运行验证}
configure/build/test 三段 preset（\srcpath{CMakePresets.json}，与根
\srcpath{AGENTS.md} 一致）：configure 侧 \texttt{macos-release}/
\texttt{linux-release}/\texttt{windows-msvc-release}/
\texttt{windows-vcpkg-release}/\texttt{full-dev}；test 侧同名 preset
仅前四个（\texttt{full-dev} 无 testPreset）。
\begin{quote}\footnotesize
\texttt{cmake --preset windows-vcpkg-release}\\
\texttt{cmake --build --preset windows-vcpkg-release}\\
\texttt{ctest --preset windows-vcpkg-release}
\end{quote}
（Linux/macOS 将 preset 名换成 \texttt{linux-release}/\texttt{macos-release} 即可。）
常用定向运行（ctest 正则，目标名以 \srcpath{tests/CMakeLists.txt} 实际注册为准）：
\begin{itemize}
  \item OPF 后端与大系统鲁棒性：\texttt{ctest -R opf\_solver\_backends}；
  \item AC/DC OPF 交叉验证：\texttt{ctest -R acopf\_dcopf\_crossval}；
  \item 无功优化：\texttt{ctest -R reactive\_power\_opt}；
  \item 三相混合 OPF：\texttt{ctest -R three\_phase\_hybrid\_opf}；
  \item DC OPF 基元与 \texttt{OPFProblem} 接口接线：\texttt{ctest -R hacdcdss}；
  \item 跨模块/图降阶回环/配网管线/多尺度：\texttt{ctest -R}
        \srcpath{'crossmodule\_integration|graph\_roundtrip|distribution\_pipeline|multiscale\_comprehensive'}；
  \item 能力打点与独立审计：\texttt{ctest -R solver\_capabilities}；
  \item MATPOWER 回归（含 DC OPF 一致性）：\texttt{ctest -R matpower\_cases}；
  \item 夜间 OPF 与换流器协调：\texttt{ctest -R}
        \srcpath{'nighttime\_opf|converter\_coordination'}；
  \item GUI 全链路：\texttt{ctest -R gui\_api\_e2e}（需
        \texttt{HACDCPF\_ENABLE\_ETAP=ON}（默认 ON）与 Python3 双重门控）；
  \item RPO 浏览器 E2E：\texttt{ctest -R rpo\_gui\_e2e}（需 Node 与
        \texttt{npx playwright install chromium}，带 \texttt{browser} 标签）。
\end{itemize}
外部依赖（MATLAB、gridlabd、Julia、OpenDSS、chromium）缺失时对应测试自动
skip；判读报告时须区分“通过”与“跳过”，与主文档诚实口径一致。


\end{document}
